\chapter{Biu-Mandara (42/80)}

The 80 Biu-Mandara languages make up the largest branch of Chadic languages. There are two major subbranches, North and South, and a minor subbranch consisting of two Hurza languages. In addition, there is one unclassified (and reportedly severely endangered) Biu-Mandara language, Jilbe [jilb1238], for which no relevant data are available. Within the Biu-Mandara North subbranch, there is a relatively large group of 25 languages called Margi-Mandara-Mofu. The tripartite name refers to its three subgroups: Bura-Marghi, Mandaraic and Mofuic. 

\section{Hurza (2/2)}

Grammatical descriptions are available for both Hurza languages. The two languages are fairly similar in the aspects of their grammars examined here. Both have a ventive extension of the form \textit{-ay(a)}, and in both languages the ventive extension is combined with a verb root \textit{n-} `be' to express ventive motion. %Both languages can express ventive CAM using a verb glossed `take' with the ventive suffix, and each has one additional strategy for expressing directed CAM. For distinguish BUY and SELL, both using a valency-increasing strategy with the root meaning BUY/TRADE: \textit{sukam} (Mbuko), \textit{səm-} (Vamé). 

%\begin{table}[H]	
%\caption{Hurza directional verb s}
%\label{tab:HurzaV}
%\begin{tabular}{lllllll} % add l for every additional column or remove as necessary
%	\lsptoprule
%	Language & VENT & ITIVE  & UP & DOWN & INTO & OUT  \\ %table header
%	\midrule
%	Mbuko (mbuk1243)&\textit{n-ay, nga} & \textit{zla}  & & \textit{dazay} & & \\	
%	Mbuko (mbuk1243)&\textit{nga} &  & &  & & \\	
%	Vamé (vame1236) & \textit{n-ay} & \textit{le} & & \textit{fat-} & \textit{li-de} & \textit{n-ay-ade} \\	
%	\lspbottomrule
%\end{tabular}
%\end{table}

\paragraph*{Mbuko (mbuk1243)} \label{mbuk}

Information on Mbuko is from a manuscript published online \citep{gravinaVerbPhraseMbuko2001} and from an online Mbuko lexicon \citep{gravinaMbukoLexicon2003,webonaryDictionnaireMbuko2016}. 

\subparagraph{Directional verbs}~

\begin{table}[H]
	\caption{Mbuko directional verbs} 
	\label{tab:mbukV}
	\begin{tabular}{lll} % add l for every additional column or remove as necessary
		\lsptoprule
		DIR & verb & gloss   \\ %table header
		\midrule
		ITIVE & \textit{zla} & `go, walk' \\
		DOWN & \textit{dazay} & `descend' \\
		\lspbottomrule
	\end{tabular}
\end{table}

As discussed below, the verb \textit{n-} `be' in Mbuko combines with a ventive suffix to express ventive directional motion. 

\subparagraph{Directional extensions}~

\begin{table}[H]
	\caption{Mbuko directional morphology}
	\label{tab:mbuk}
	\begin{tabular}{llll} % add l for every additional column or remove as necessary
		\lsptoprule
		forms & source gloss & direction  & other functions  \\ %table header
		\midrule
		\textit{-ay}  &  \textsc{ing}  &   \textsc{vent}  &   \textsc{subs.vent}  \\
		\lspbottomrule
	\end{tabular}
\end{table} 

\citet[7--8]{gravinaVerbPhraseMbuko2001} describes two ``direction'' suffixes in Mbuko: \textit{-ay} `ingressive' and \textit{-ák} `egressive'. These suffixes are transparently related (etymologically) to the adverbial ``destination markers'': \textit{àhày} `to here -- toward deictic centre' and \textit{áyāk} `to there -- away from deictic centre' \citep[18]{gravinaVerbPhraseMbuko2001}. The direction suffixes are said to increase the valency of verbs in that they allow verbs to take a locative complement \citep[11]{gravinaVerbPhraseMbuko2001}.

The suffix \textit{-ák} `eggressive' is ``found on a smaller number of verbs'' and is only shown with the verbs \textit{njaɗ-} `find' and \textit{təmah-} `receive'. It appears to have a translocative function (`there'), although one of the two examples given suggests a possible prior motion interpretation. More data are needed to determine if this suffix has a directional function`   . 

The more common suffix is \textit{-ay} `ingressive' which functions as a ventive directional with motion verbs, as in (\ref{ex2:mbuk1}).

\ea\label{ex2:mbuk1}
\langinfo{Mbuko verb \textit{ngəm} `leave' with ventive suffix}{}{\cite[9]{gravinaVerbPhraseMbuko2001}} \\
\gll a-ngəm-ay a-kutov wa	\\
\gl{3}\gl{sg}.\gl{pfv}-leave-\gl{vent} to-stomach \textsc{source}	\\
\glt `He came out of the stomach.'
\z

Exceptionally, the ventive suffix is said to have an itive directional interpretation with the verb \textit{haw} `run': \textit{haw-ay} `run away', in which case \citet[7]{gravinaVerbPhraseMbuko2001} says it ``functions as a derivational affix.''  With the verb `climb' the interpretation is that the motion can be taking place in several different directions: \textit{jən-ay} `climb up/down/out' \citep[8]{gravinaVerbPhraseMbuko2001}. So while the suffix \textit{-ay} is described as primarily a ventive extension, it is described as having a range of possible interpretations.
 
There are a few cases of the ventive suffix with non-motion verbs which can be interpreted as subsequent ventive motion, normally with a CAM component as well. The clearest example is the verb `buy': \textit{sukum-ay} `buy and bring back'. 

In a pattern not seen in any other Chadic language besides the closely-related Vamé, Gravina analyzes the verb \textit{nay} `come' as composed of the verb \textit{na} `be (somewhere)' and the ventive suffix, as in (\ref{ex2:mbuk3}). In this analysis, the ventive suffix can be seen as contributing not only the ventive direction, but also the fact of motion itself.

\ea\label{ex2:mbuk3}
\langinfo{Mbuko verb \textit{n-} `be' with ventive suffix}{}{\cite[8]{gravinaVerbPhraseMbuko2001}} \\
\gll a-n-ay ahay a-juvok \\	
\gl{3}\gl{sg}.\gl{pfv}-be-\gl{vent} to.here to-guest.hut	\\
\glt `He came to the guest hut.'
\z 

\subparagraph{Caused accompanied motion (CAM)}

Verbs expressing (non-directed) CAM include \textit{ray} `carry (several things)' [French: `\textit{apporter} (\textit{plusieurs choses})'] and \textit{tavak} `carry on the head' [French: `\textit{porter sur la tête}'] \citep{gravinaMbukoLexicon2003,webonaryDictionnaireMbuko2016}. Ventive CAM may be expressed by the ventive form of the verb \textit{gəɓ-} `take', as in (\ref{ex2:mbuk4}).

\ea\label{ex2:mbuk4}
\langinfo{Mbuko verb \textit{gəɓ-} `take' with ventive suffix}{}{\cite[8]{gravinaVerbPhraseMbuko2001}} \\
\gll  a-gəɓ-ay anan ahay dədew uho	\\
\gl{3}\gl{sg}.\gl{pfv}-take-\gl{vent} it to.here snake outside	\\
\glt `He brought the snake outside.'
\z 

It is also possible for the motion predicates \textit{zla} `go' and \textit{n-ay} `come (< be-\textsc{vent})' to be used transitively (without any morphological marking of valency change) in order to express itive and ventive CAM, as in (\ref{ex2:mbuk5}) and (\ref{ex2:mbuk6}).\footnote{In example (\ref{ex2:mbuk6}) Gravina analyzes the low tone on the 2\textsc{pl} suffix to be evidence that there is an underlying ventive suffix that has been deleted due to the following suffix.}

\ea\label{ex2:mbuk5}
\langinfo{Mbuko verb \textit{zla} `go' as a transitive predicate}{}{\cite[22]{gravinaVerbPhraseMbuko2001}} \\
\gll Ta-zla anan zana a-tinen ahay pə-tila aga njavan ahay\\
\gl{3}\gl{pl}.\gl{pfv}-go \gl{obj}.\gl{foc} cloth of-they \gl{pl} on-tailor house.of guinea-fowl \gl{pl}\\
\glt `They took their clothes from the tailor to the guinea-fowls’ house.'

\ex\label{ex2:mbuk6}
\langinfo{Mbuko verb \textit{nay} `come' as a transitive predicate}{}{\cite[18]{gravinaVerbPhraseMbuko2001}} \\
\gll ki n-èn anan ahay sla ahay a-guvo uno sabay \\
2\gl{pl}.\gl{imp} be-\gl{vent}.2\gl{pl} \gl{obj}.\gl{foc} to.here cow \gl{pl} to-field my no.longer	\\
\glt `Don’t bring your cows here to my field again.'
\z 

\subparagraph{Buy-sell verbs}

The verb \textit{sukom} can mean `buy' and `sell'. If followed by a causative marker \textit{sukom anan} then it means `sell' \citep{gravinaMbukoLexicon2003,webonaryDictionnaireMbuko2016}. Note that the verb can occur with the ventive suffix and still mean `buy', as in \textit{sukum-ay} `buy and bring back' \citep[8]{gravinaVerbPhraseMbuko2001}.

\paragraph*{Vamé (vame1236)} \label{vame}

Vamé grammar is described in a manuscript published online \citep{kinnairdVameVerbalSystem2006}.

\subparagraph{Directional verbs}~

\begin{table}[H]
	\caption{Vamé directional verbs} 
	\label{tab:vameV}
	\begin{tabular}{lll} % add l for every additional column or remove as necessary
		\lsptoprule
		DIR & verb & gloss   \\ %table header
		\midrule
		ITIVE & \textit{le} & `go' \\
		DOWN & \textit{fat} & `climb down, get down' \\
		\lspbottomrule
	\end{tabular}
\end{table}

As discussed below, just as in Mbuko, the verb \textit{n-} `be' in Vamé combines with a ventive suffix to express ventive directional motion. This complex form can be further combined with a morpheme \textit{-ade} to express outward motion, but this morpheme is not shown to be a productive suffix in the available data. Likewise, the verb \textit{le} `go' can combine with a morpheme \textit{-de} to express inward motion, but this suffix is not attested elsewhere. As these verb stems are potentially composed of multiple morphemes they are not counted as directional verb roots. 

\subparagraph{Directional extensions}~

\begin{table}[H]
	\caption{Vamé directional morphology}
	\label{tab:vame}
	\begin{tabular}{llll} % add l for every additional column or remove as necessary
		\lsptoprule
		forms & source gloss & direction  & other functions  \\ %table header
		\midrule
		\textit{-aya, -ki}  &  \textsc{cp}  &   \textsc{vent} &   \textsc{subs.vent}  \\
		\lspbottomrule
	\end{tabular}
\end{table} 

\citet{kinnairdVameVerbalSystem2006} describes a ``directional'' suffix with a ``centripetal'' function glossed \textsc{cp}, which is a ventive suffix. The suffix \textit{-aya} becomes \textit{-ki} before the present tense suffix \textit{-ke}.\footnote{\citet[56]{kinnairdVameVerbalSystem2006} proposes that the ventive suffix is \textit{-ayà} with intransitive verbs and \textit{-iyá} with transitive verbs. This pattern does not appear to hold in all the examples given, and often no suffix-initial vowel appears.} With a manner-of-motion verb, this suffix can have a ventive directional function, as in (\ref{ex2:vame1}).

\ea\label{ex2:vame1}
\langinfo{Vamé verb \textit{hav} `fly' with ventive suffix}{}{\cite[56]{kinnairdVameVerbalSystem2006}} \\
\gll  fan hav-ki-ke	\\
bird fly-\gl{vent}-\gl{prs}	\\
\glt `The bird is flying (towards us).'
\z 

As in Mbuko, the same suffix is used with the verb stem \textit{n-} `be' to express ventive directional motion, as in (\ref{ex2:vame2}).

\ea\label{ex2:vame2}
\langinfo{Vamé verb \textit{n-} `be, come' with ventive suffix}{}{\cite[56]{kinnairdVameVerbalSystem2006}} \\
\gll  na eveŋ ni-ki-ke	\\
behold rain be-\gl{vent}-\gl{prs}	\\
\glt `Look, the rain is coming.'
\z 

The same suffix occurs in many contexts where it is not clear that it has ventive meaning. For example, the verb \textit{mə} `remain' occurs with a ventive suffix in (\ref{ex2:vame3}) and the translation suggests a distal rather than directional interpretation.

\ea\label{ex2:vame3}
\langinfo{Vamé verb \textit{mə} `remain' with ventive suffix}{}{\cite[21]{kinnairdVameVerbalSystem2006}} \\
\gll m-aya Malataŋgwe saha	\\
remain-\gl{vent} Malatangwe \gl{dem}	\\
\glt `There was only left (there), Malatangwe.'
\z

In comparison with Wuzlam (Ouldemé), \citet[48]{kinnairdVameVerbalSystem2006} describes several other ``spatial'' suffixes in Vamé. Some of these are used in directional contexts, for example, compare the verb \textit{le} `go' with \textit{lide} `go into'. However, \citet[50]{kinnairdVameVerbalSystem2006} shows that the suffix \textit{-de} does not consistently have the same directional meaning across verb stems. There are likely etymological explanations for some of the patterns Kinnaird identifies, but none of the suffixes are considered to be productive motion-related verbal morphology in a synchronic analysis of the grammar. 

\subparagraph{Caused accompanied motion (CAM)}

In addition to a verb root \textit{mə-} `bring', CAM can be expressed using the ventive suffix with one of several possible verb roots including \textit{gi-} `do' (as in \ref{ex2:vame6}),  \textit{təv-} `take' (as in \ref{ex2:vame7}) and \textit{buri} `catch'.

\ea\label{ex2:vame6}
\langinfo{Vamé verb \textit{gi-} `do' with ventive suffix}{}{\cite[11]{kinnairdVameVerbalSystem2006}} \\
\gll ɗa gi-yə-nə-ka ahwam bunə-ɗi-l-pə-na səlay	\\
3\gl{pl} do-\gl{vent}-3\gl{sg}.\gl{dat}-\gl{pst} water wash-\gl{inf}-\gl{ins}-\gl{loc}-3\gl{sg}.\gl{dat} foot \\
\glt `They brought him water with which to wash his feet.'

\ex\label{ex2:vame7}
\langinfo{Vamé verb \textit{təv-} `take' with ventive suffix}{}{\cite[55]{kinnairdVameVerbalSystem2006}} \\
\gll  na bəskur ga əŋ təv-ya-la-ka ɗu kuza	\\
here.is bike which I take-\gl{vent}-\gl{ins}-\gl{pst} person illness \\
\glt `Here is the bike with which I brought the sick person.'
\z 

Itive CAM can be expressed by the use of an ``instrumental{\slash}comitative'' suffix \textit{-la} (also called ``causative/instrumental'') with the verb \textit{le} `go': \textit{lila} `take away' \citep[43]{kinnairdVameVerbalSystem2006}.

\subparagraph{Buy-sell verbs}

The verb \textit{səm-} is glossed `buy'. Adding the ``instrumental{\slash}comitative'' suffix \textit{-la} is said to result in a ``causative'' form, \textit{səməla} `sell' \citep[47]{kinnairdVameVerbalSystem2006}.

\section{North Biu-Mandara (27/50)}

North Biu-Mandara is the largest subbranch, consisting of 50 languages of which 27 are included here. It is further divided into seven groups, the largest one being the 25 languages of the Margi-Mandara-Mofu group. 

%\begin{table}[H]
%	\caption{North Biu-Mandara group directional extensions}
%	\label{tab:North}
%	\begin{tabular}{lllllll} % add l for every additional column or remove as necessary
%		\lsptoprule
%		language & VENT & ITIVE  & UP & DOWN & INTO & OUT  \\ %table header
%		\midrule
%		Gidar	(gida1247) & \textit{-i} & & & & & \\
%		Bana	(bana1305) & \textit{-ke}  & \textit{-r̂e}/\textit{-ghe} & & & & \\
%		Psikye	(psik1239) & 	\textit{-ke} & \textit{-gu} & \textit{-(a)me} & \textit{-awa} & \textit{-mbe}? & \textit{-mte} \\
%		Psikye	(psik1239) & &  & \textit{-ate} & \textit{-(i)yi} &  &  \\
%		Buduma	(budu1265) & & & & & \textit{-(a)l} & \\
%		Lagwan	(lagw1237) & & & & & \textit{-li} & \\
%		Makary Kotoko	(mpad1242) & & & & & & \\
%		Lamang	(lama1288) & \textit{-gha} & & \textit{-f(i)} & \textit{-gaa} & & \textit{-b(i)}  \\
%		Lamang	(lama1288) & & & & \textit{-ɗi} & &  \\
%		Bura	(bura1292) &  & \textit{-mta} & & \textit{-hi} & \textit{-(k)wa}  & \textit{-bila} \\
%		Bura	(bura1292) &  &  & &  & \textit{-u}  & \textit{-ɓəla} \\
%		Cibak	(ciba1236) & & \textit{-nta} ? & & & &  	\textit{-ba} \\
%		Margi	(marg1265) & & \textit{-nà} & & \textit{-ía} & \textit{-wá} & \textit{-bá}  \\
%		Huba	(huba1236) & & \textit{-nà} & & \textit{-yà}  & & \textit{-bìyà} \\
%		Dghwede	(dghw1239) & \textit{-àwá} & \textit{-dá} & \textit{-gwá}  & \textit{-tígè} & \textit{-mè} & \textit{-àwílè} \\
%		Dghwede	(dghw1239) & &  & \textit{-fè}  & \textit{-àwáyà} &  &  \\
%		Dghwede	(dghw1239) & &  & \textit{-fè}  & \textit{-àwí} &  &  \\
%		Parkwa	(park1239) & \textit{-ətsa} & & \textit{-u} & \textit{-aha} &  \textit{-əkwa}  & \textit{-edá} \\
%		Parkwa	(park1239) & \textit{-atsa} & & \textit{-əlu} & \textit{-əla} &  \textit{-akwa}  & \textit{-əlá} \\
%		Parkwa	(park1239) & \textit{-atsa} & & \textit{-alu} & \textit{-ada} &    & \textit{adá} \\
%		Glavda	(glav1244) & & \textit{-d(a)} & \textit{-(i)t}  & 	\textit{-(x)i} & \textit{-m} & \\
%		Malgwa	(malg1251) & & &\textit{-tə́} & &\textit{-ə́m} & \textit{-sə́} \\
%		Wandala	(wand1278) & & & & &\textit{-m} &\textit{-s}  \\
%		Merey	(mere1246) & \textit{-aw} & & & & & \\
%		Zulgo-Gemzek	(zulg1242) & \textit{-ara} & \textit{aha} & & & & \\
%		Zulgo-Gemzek	(zulg1242) &  \textit{-aw} &  & & & & \\
%		Zulgo-Gemzek	(zulg1242) & \textit{-ya} &	 & & & & \\
%		Mofu-Gudur	(mofu1248) & \textit{-wa} & & & & & \\
%		Mada	(mada1293) & \textit{-re} &\textit{-ro} & & \textit{-ra} & \textit{-(v)a} & 	\textit{-ra} \\
%		Moloko	(molo1266) &\textit{alaj} & \textit{ala} & & & & \\
%		Muyang	(muya1243) & 	\textit{-biyu} & & & & 	\textit{-iyu}  & \textit{-aya} \\
%		Wuzlam	(wuzl1236) & \textit{-ārá} & \textit{-ārāy}  & & & & \\
%		Wuzlam	(wuzl1236) & \textit{-ārá} & \textit{ə̄rə̄}  & & & & \\
%		North Giziga	(nort3047) & \textit{-áwà} & & & & & \\
%		North Giziga	(nort3047) & \textit{-óò} & & & & & \\
%		North Giziga	(nort3047) & \textit{-o} & & & & & \\
%		Mbara	(mbar1260) & \textit{-sì} & & & & & \\
%		Musgu	(musg1254) & 	\textit{sí}  &\textit{di}  & & & & \\	
%		\lspbottomrule
%	\end{tabular}
%\end{table}
	
\subsection{Gidar (1/1)}

\paragraph*{Gidar (gida1247)} \label{gida}

Gidar is described in a reference grammar \citep{frajzyngierGrammarGidar2008}. 

\subparagraph{Directional verbs}~

\begin{table}[H]
	\caption{Gidar directional verbs} 
	\label{tab:gidaV}
	\begin{tabular}{lll} % add l for every additional column or remove as necessary
		\lsptoprule
		DIR & verb & gloss   \\ %table header
		\midrule
		ITIVE & \textit{(n)d, mbat} & `go' \\
		VENT & \textit{zá} & `come' \\
		INTO& \textit{tóŋ} & `enter' \\
		OUT & \textit{gil} & `get out' \\
		\lspbottomrule
	\end{tabular}
\end{table}

\subparagraph{Directional extensions}~

%\begin{table}[H]
%	\caption{Gidar directional morphology}
%	\label{tab:gida}
%	\begin{tabular}{llll} % add l for every additional column or remove as necessary
%		\lsptoprule
%		forms & source gloss & direction  & other functions  \\ %table header
%		\midrule
%		\textit{-i} &  \textsc{vent}  & \textsc{dir.vent} (lexically restricted)  &  no data  \\
%		\lspbottomrule
%	\end{tabular}
%\end{table}

\citet[196--198]{frajzyngierGrammarGidar2008} describes a ``ventive extension'' in the form of a suffix \textit{-i}. The ventive is only found in interlinear glosses of three verbs, \textit{(n)d} `go', as in (\ref{ex2:gida1}), \textit{z} `come' (also glossed `go'), as in (\ref{ex2:gida2}) and \textit{k} `take, bring' (see caused accompanied motion below). Note that for two of these examples the form described as ventive does not actually have the phonetic form [i]. The putative \textit{-i} suffix only appears with one verb stem and so there is not sufficient evidence available to consider it part of the verbal morphology.  

\ea\label{ex2:gida1}
\langinfo{Gidar (putative) ventive form of \textit{d} ‘go’}{}{\cite[197]{frajzyngierGrammarGidar2008}} \\
\gll tə̀-ddé-kè \\
3\gl{f}-go.\gl{vent}-\gl{prf} \\
\glt `She arrived.'

\ex\label{ex2:gida2}
\langinfo{Gidar (putative) ventive form of \textit{z} ‘go, come’}{}{\cite[197]{frajzyngierGrammarGidar2008}} \\
\gll tə̀-zzé \\
3\gl{f}-go.\gl{vent}-\gl{prf} \\
\glt `She returned home.'
\z

\citet[50]{frajzyngierGrammarGidar2008} proposes a phonological alternation that could hypothetically account for this, however, \citet[313]{schuhChadicCornucopia2017} mentions in a footnote that he published an (inaccessible) review in which he shows that ``this cannot be part of Gidar phonology'' as it contradicts the data. The earlier field notes of \citet{schuhDonneesLangueGidar1984} also mention a ventive form, but only provide a few examples and a suggestion that palatal prosody may be associated with ventive direction.

\subparagraph{Caused accompanied motion (CAM)}

The ``ventive'' form of the verb \textit{k} `take' (also glossed `take away' \citep[196]{frajzyngierGrammarGidar2008}) can be used to express ventive caused accompanied motion, as in (\ref{ex2:gida6}). The verb \textit{k} `take' appears to be a CAM verb, in contrast with \textit{rmà} `take' and \textit{gə̀mə̀} `take' which express taking possession of an object. 

\ea\label{ex2:gida6}
\langinfo{Gidar ventive form of \textit{k} ‘take’}{}{\cite[196]{frajzyngierGrammarGidar2008}} \\
\gll	nə̀-k-í-ì-ká \\
1\gl{sg}-take-\gl{vent}-3\gl{pl}-\gl{prf} \\
\glt `I brought them.'
\z 

\subparagraph{Buy-sell verbs}

There are two different verb roots in Gidar: \textit{ə́lbàhà} `buy' and \textit{ə̀və́l} `sell'. 

\subsection{Higic (2/5)}

Of the five Higic languages, an analysis of the grammar is only available for two. The three languages not discussed here are: Hya (hyaa1239), Kamwe (kamw1239; formerly known as Higi) and Kirya-Konzel (kiry1234). %The languages included, Bana and Psikye, both have ventive and itive extensions. Psikye has a considerable number of additional directional extensions, and, strikingly, the forms of the directional extensions match the forms of directional verbs in Bana. The difference is that in Bana these forms appear to be restricted to one verbal stem (at most two) and so they are not productive morphology. Both languages use directional extensions with a `take' verb to express Caused Accompanied Motion. Psikye also uses directional extensions to distinguish BUY and SELL, whereas Bana uses two distinct verb stems. 

\paragraph*{Bana (bana1305)} \label{bana}

Information on Bana is from a pedagogical grammar \citep{zuchGrammaireExercicesLangue1977}, a short journal article \citep{wente-lukasZurSprachlichenStellung1973}, an MA thesis with an appended lexicon \citep{hofmannPreliminaryPhonologyBana1990}, and an unpublished lexicon \citep{gigerBanaProvisionalLexicon2003}. 

\subparagraph{Directional verbs}~

\begin{table}[H]
	\caption{Bana directional verbs} 
	\label{tab:banaV}
    \fittable{
	\begin{tabular}{lll} % add l for every additional column or remove as necessary
		\lsptoprule
		DIR & verb & gloss   \\ %table header
		\midrule
		ITIVE & \textit{dzà} & `\textit{aller} [go]' \\
		VENT & \textit{sə́kə́} & `\textit{venir} [come]' \\
		UP & \textit{dzə́mə̀} & `\textit{aller} (\textit{montant beaucoup}) [go (sharply upwards)]' \\
		UP & \textit{dzàtì} & `\textit{aller} (\textit{en montant un peu}) [go (slightly upwards)]' \\
		ITIVE.DOWN 	& \textit{dzghwà} &	`\textit{aller} (\textit{descendant beaucoup}) [go (sharply downwards)]' \\ 
		ITIVE.DOWN 	& \textit{dzì} &	`\textit{aller} (\textit{descendant un peu}) [go (slightly downwards)]' \\ 
		VENT.DOWN & \textit{skwá} & `\textit{venir en bas} [come down]' \\
		INTO & \textit{dzə̀mbə̀} & `\textit{entrer} [enter]' \\
		ITIVE.OUT & \textit{dzə̀vrì} & `\textit{sortir}, \textit{aller} (\textit{assez loin}) [go out, go (fairly far)]' \\	
		VENT.OUT & \textit{sə́vrí} & `\textit{venir dehors}, \textit{venir de loin} [come out, come from far]' \\
		\lspbottomrule
	\end{tabular}
	}
\end{table}

The verbs in Table \ref{tab:banaV} are taken from \citet{gigerBanaProvisionalLexicon2003} and are generally described in the same way in other sources.\footnote{Slightly different definitions are offered by \citet[72]{zuchGrammaireExercicesLangue1977} who proposes that a crucial element of the meaning is whether the path of motion involves a river. For example, one downward verb \textit{dzeghwà} is defined as `\textit{on descend vers un fleuve} [we descend towards a river]' and the other downward verb \textit{dzì} is defined as `\textit{on descend, mais on ne traverse pas de fleuve} [we descend, but we do not cross a river]'.} While the obvious pattern in the roots of Bana directional verbs suggests that these may be decomposed morphologically, in most cases particular verb endings are only attested once and not found with other verb stems. The one exception is the outward motion verbs which end in  \textit{-vri}. This ending is similar to the outward directional extension in the closely-related Psikye language, but in the available Bana data, no other verbs are found ending in \textit{-vri}.

\subparagraph{Directional extensions}~

\begin{table}[H]
	\caption{Bana directional morphology}
	\label{tab:bana}
	\begin{tabular}{llll} % add l for every additional column or remove as necessary
		\lsptoprule
		forms & source gloss & direction  & other functions  \\ %table header
		\midrule
		\textit{-ke}  & \textit{rapprochement}  &  \textsc{dir.vent} &   \textsc{subs.vent.cam}  \\
		\textit{-r̂e}/\textit{-ghe}  &  \textit{éloignement}  & \textsc{dir.itive}  & \textsc{subs.itive.cam}  \\
		\lspbottomrule
	\end{tabular}
\end{table} 

Bana has a ventive suffix \textit{-ke} which \citet[44]{zuchGrammaireExercicesLangue1977} says expresses ``\textit{rapprochement}'' and an itive suffix \textit{-gha} which \citet[44]{zuchGrammaireExercicesLangue1977} transcribes as \textit{-r̂a} and says expresses ``\textit{eloignement}''. Examples are included of these suffixes with six different verb stems, all transitive. With a transitive motion verb, as in (\ref{ex2:bana3}) and (\ref{ex2:bana4}), the suffix indicates the direction of the path of motion of the patient. 

\ea\label{ex2:bana3}
\langinfo{Bana ventive form of \textit{r̂uni} ‘send’}{}{\cite[44]{zuchGrammaireExercicesLangue1977}} \\
\gll a r̂uni-ke zer̂we	\\
3\gl{sg}.\gl{f} send-\gl{vent} child \\
\glt `She sent the child (towards me).' \newline [French: `\textit{Elle a envoyé l'enfant} (\textit{vers moi}).']

\ex\label{ex2:bana4}
\langinfo{Bana itive form of \textit{r̂uni} ‘send’}{}{\cite[44]{zuchGrammaireExercicesLangue1977}} \\
\gll a r̂uni-r̂e zer̂we	\\
3\gl{sg}.\gl{f} sent-\gl{itive} child	\\
\glt `She took the child away from me.' [French: `\textit{Elle m'a enlevé l'enfant}.']
\z 

Translations of some non-motion verb stems with a directional extension, as in (\ref{ex2:bana5}), (\ref{ex2:bana6}), (\ref{ex2:bana7}) and (\ref{ex2:bana8}), show an interpretation of (ventive or itive) subsequent motion. In these examples, both the agent and the patient appear to change locations after the activity predicated by the verb has ended, giving the predicate a type of caused accompanied motion sense as well. 

\ea\label{ex2:bana5}
\langinfo{Bana ventive form of \textit{pla} ‘find’}{}{\cite[44]{zuchGrammaireExercicesLangue1977}} \\
\gll Pla-ke zer̂we!	\\
seek-\gl{vent} child \\
\glt `Go look for the child (and bring him to me)!' \newline [French: `\textit{Va chercher l'enfant} (\textit{et amène-le vers moi}) !']

\ex\label{ex2:bana6}
\langinfo{Bana itive form of \textit{pla} ‘find’}{}{\cite[44]{zuchGrammaireExercicesLangue1977}} \\
\gll Pla-r̂e zer̂we! \\
seek-\gl{itive} child \\
\glt `Go look for the child and take him (to his mother, for example)!' \newline [French: `\textit{Va chercher l'enfant et emmène-le} (\textit{à sa mère par ex.}) !']

\ex\label{ex2:bana7}
\langinfo{Bana ventive form of \textit{r̂eli} ‘steal’}{}{\cite[44]{zuchGrammaireExercicesLangue1977}} \\
\gll e r̂eli-ke gena \\
1\gl{sg} steal-\gl{vent} money \\
\glt `I stole some money. (I stole it elsewhere and I kept it at my place.)' \newline [French: `\textit{J'ai volé de l'argent} (\textit{je l'ai volé ailleurs et je l'ai gardé chez moi}).']

\ex\label{ex2:bana8}
\langinfo{Bana itive form of \textit{r̂eli} ‘steal’}{}{\cite[44]{zuchGrammaireExercicesLangue1977}} \\
\gll e r̂eli-r̂e gena	\\
1\gl{sg} steal-\gl{itive} argent	\\
\glt `I stole some money. (I stole it here and I took it elsewhere.)' \newline [French: `\textit{J'ai volé de l’argent} (\textit{je l'ai volé ici et je l'ai transporté ailleurs}).']
\z 

\subparagraph{Caused accompanied motion (CAM)}

Directed CAM can be expressed by using the ventive or itive suffix with the verb \textit{ɓe} `take' [French: `\textit{prendre}'] as in (\ref{ex2:bana1}) and (\ref{ex2:bana2}). Another verb whose ventive and itive forms can express directed CAM is \textit{ime} `gather' [French: `\textit{ramasser}'].

\ea\label{ex2:bana1}
\langinfo{Bana ventive form of \textit{ɓe} ‘take’}{}{\cite[44]{zuchGrammaireExercicesLangue1977}} \\
\gll ɓe-ke shi va	\\
take-\gl{vent} thing \gl{dem}	\\
\glt `Bring (me) those things!' [French: `\textit{Apporte} (\textit{moi}) \textit{ces choses} !']

\ex\label{ex2:bana2}
\langinfo{Bana itive form of \textit{ɓe} ‘take’}{}{\cite[44]{zuchGrammaireExercicesLangue1977}} \\
\gll ɓe-r̂e shi va	\\
take-\gl{itive} things \gl{dem} \\	
\glt `Take those things away (from me)!' [French: `\textit{Enlève} (\textit{moi}) \textit{ces choses} !']
\z 

\subparagraph{Buy-sell verbs}

Buying and selling are expressed with different verbs stems: \textit{ɗə́l} `sell' and \textit{pá} `pay' or \textit{wá} `pay, buy' \citep{gigerBanaProvisionalLexicon2003}.

 \paragraph*{Psikye (psik1239)} \label{psik}

Information on Psikye grammar is from a doctoral dissertation \citep{smithKapsikiLanguage1969}.

\subparagraph{Directional verbs}~

\begin{table}[H]
	\caption{Psikye directional verbs} 
	\label{tab:psikV}
	\begin{tabular}{lll} % add l for every additional column or remove as necessary
		\lsptoprule
		DIR & verb & gloss   \\ %table header
		\midrule
		ITIVE & \textit{dze} & `go' \\
		VENT & \textit{sé} & `come' \\
		\lspbottomrule
	\end{tabular}
\end{table}

In contrast with the closely related language Bana, Psikye appears to have only two basic directional verbs. Inward/outward and upward/downward motion are expressed by adding directional suffixes to the basic directional motion verbs, as discussed below. 

\subparagraph{Directional extensions}~

\begin{table}[H]
	\caption{Psikye directional morphology}
	\label{tab:psik}
	\begin{tabular}{llll} % add l for every additional column or remove as necessary
		\lsptoprule
		forms & gloss & direction  & other functions  \\ %table header
		\midrule
		\textit{-ke}  &  \textsc{vent}  &  no data &  \textsc{subs.vent}  \\
		\textit{-gu}  &  \textsc{itive}  & \textsc{itive}  &   \textsc{subs.itive}  \\
		\textit{-(aka)me}  &  \textsc{up1}  &    \textsc{up} &  lexicalized  \\
		\textit{-(aka)ate}  &  \textsc{up2}  &   \textsc{up}  &  \textsc{subs.up}, lexicalized \\
		\textit{-(ak)awa}, \textit{-(ak)aa}   &  \textsc{down1}  &  \textsc{down}  &    \\
		\textit{-(i)yi} &  \textsc{down2}  &   \textsc{down}  &    \\
		\textit{-(aka)ve}  &  \textsc{across}  &    \textsc{out}  & lexicalized \\
		\textit{-(aka)mbe}  &  \textsc{into}  &   \textsc{into} &       \\
		\textit{-mte}  &  \textsc{bush}  &   \textsc{bush}   &   \textsc{compl}  \\
		\lspbottomrule
	\end{tabular}
\end{table} 

\citet{smithKapsikiLanguage1969} describes nine motion-related verbal suffixes in Psikye, as shown in Table \ref{tab:psik}. Smith recognizes that these suffixes have different functions with motion and non-motion verbs, but confusingly treats identical forms as unrelated suffixes based on whether they are ``used normally only with verbs of motion'' \citep[113]{smithKapsikiLanguage1969}. At the end of the dissertation, \citet[163--164]{smithKapsikiLanguage1969} briefly considers how some of these identical forms might be related. 

The suffix \textit{-ke} is described as meaning ``to do something and bring here'' and the suffix \textit{-gu} as meaning ``to do something and go away with it, or do something at one place when the effect will be felt elsewhere'' \citep[113--114]{smithKapsikiLanguage1969}. Only a few examples of these forms are found in Smith's data, and none with intransitive motion verbs. The ventive suffix \textit{-ke} is only shown in a subsequent motion/CAM interpretation, as in (\ref{ex2:psik3}), and does not appear in any examples where it has a directional function.\footnote{Note that the ventive directional suffix in the closely related language Bana has the same form \textit{-ke}.} \citet[165]{smithKapsikiLanguage1969} also mentions that a suffix \textit{-ke} is ``often used with no discernible meaning other than to mark an action as completed.'' 

\ea\label{ex2:psik3}
\langinfo{Psikye ventive form of \textit{gú} `dip'}{}{\cite[114]{smithKapsikiLanguage1969}} \\
\gll pá bde gú-ke ŋ'yɛ́ yɛmú	\\
then 3\gl{sg}.\gl{f} dip-\gl{vent} 1\gl{pl} water	\\
\glt `Then she dipped and brought us water.'
\z 	

The itive suffix \textit{-gu} is shown in a directional function in one example with a transitive motion verb, in (\ref{ex2:psik1}), and in a subsequent motion (or subsequent CAM) interpretation in another, shown in (\ref{ex2:psik2}).

\ea\label{ex2:psik1}
\langinfo{Psikye itive form of \textit{gene} `send'}{}{\cite[113]{smithKapsikiLanguage1969}} \\
\gll ’a ’yá ké-gene-gu ndá 'wú'we \\
\textsc{active} 1\gl{sg} \gl{ptcp}-send-\gl{itive} ?? 3\gl{sg}.\gl{m}	\\
\glt `I sent money to him.'

\ex\label{ex2:psik2}
\langinfo{Psikye itive form of \textit{tle} `cut'}{}{\cite[113]{smithKapsikiLanguage1969}} \\
\gll  tle-gu le xwá	\\
cut-\gl{itive} \gl{ins} knife	\\
\glt `Cut it with the knife and take it away!'
\z 

There are two suffixes described as expressing upward motion, and two described as expressing downward motion. One upward suffix \textit{-(a)me} (here glossed \textsc{up1}), as in (\ref{ex2:psik4}), is said to mean ``to go or come up precipitously'', and the other upward suffix \textit{-ate} means ``to come or go up but not too steeply'' (glossed \textsc{up2}), as in (\ref{ex2:psik5}).\footnote{Compare the directional verbs in Bana that are also described as differentiating vertical paths of motion for steepness.} Note that the steepness of the upward motion is generally not made explicit in free translations, but can be inferred from the context in a few cases.

\ea\label{ex2:psik4}
\langinfo{Psikye (steeply) upward form of \textit{dze} `go'}{}{\cite[117]{smithKapsikiLanguage1969}} \\
\gll  a ké-dze-mé va cɛ	\\
\textsc{active} \gl{ptcp}-go-\gl{up}1 tree	\\
\glt `(He) climbed the tree.'

\ex\label{ex2:psik5}
\langinfo{Psikye (slightly) upward form of \textit{sé} `come'}{}{\cite[116]{smithKapsikiLanguage1969}} \\
\gll ka-s-até xɛ́ \\
\gl{ipfv}-come-\gl{up}2 up \\
\glt `They are coming up.'
\z

Smith also notes that the \textit{-ate} `(slightly) upward' forms have additional idiomatic interpretations: `stand up' with the verb \textit{sé} `come', as in (\ref{ex2:psik6}), and  `be standing' with the verb \textit{dze} `go', as in (\ref{ex2:psik7}).

\ea\label{ex2:psik6}
\langinfo{Psikye (slightly) upward form of \textit{sé} `come'}{}{\cite[117]{smithKapsikiLanguage1969}} \\
\gll pá 'yá s-áte \\
then 1\gl{sg} come-\gl{up}2 \\
\glt `Then I got up.'

\ex\label{ex2:psik7}
\langinfo{Psikye (slightly) upward form of \textit{dze} `go'}{}{\cite[117]{smithKapsikiLanguage1969}} \\
\gll nya malé dze-mé va cɛ \\
\gl{dem} woman go-\gl{up}1 near house \\
\glt `This woman is standing near her house.'
\z

The (steep) upward suffix \textit{-amé} appears at least once with a non-motion verb, \textit{men} `do', as shown in (\ref{ex2:psik8}). There is no motion-related meaning in the translation.

\ea\label{ex2:psik8}
\langinfo{Psikye (steep) upward form of \textit{men} `do'}{}{\cite[79]{smithKapsikiLanguage1969}} \\
\gll a ké-men-amé ǹši mbe zewé	\\
\textsc{active} \gl{ptcp}-do-\gl{up}1 knot in rope	\\
\glt `He made a knot in the rope.'
\z 

The (steep) downward suffix is \textit{-awa} or \textit{-aa}. The suffix \textit{-iyi} is described as a (slightly) downward suffix (glossed \textsc{down2}) but it is unclear that there is actually any vertical dimension to its meaning. The majority of the examples found of this suffix (12/13) are translated with adverbs like `by', `over' and `across', as in (\ref{ex2:psik9}). The only case where a `down' translation is given appears to be an elicited form of the verb in isolation. \citet[139]{smithKapsikiLanguage1969} comments that ``\textit{-yi} indicates action downward in some semiological contexts and over in others'' though it is not made clear what contexts are relevant to each meaning. \citet[114--115]{smithKapsikiLanguage1969} also describes a suffix \textit{-yi}, as in (\ref{ex2:psik19}), used with non-motion verbs meaning ``to do something and leave it, not to take it up again---at least not soon; to do something and not expect to receive direct benefit from the action.'' This is presumably the (slightly) downward suffix being used with non-motion verbs, in which case it does not have any motion-related meaning. 

\ea\label{ex2:psik9}
\langinfo{Psikye (slightly) downward form of \textit{dze} `go'}{}{\cite[118]{smithKapsikiLanguage1969}} \\
\gll dz-iyí 'yá vaké \\
go-\gl{down}2 1\gl{sg} by.here \\
\glt `I am going by here...'

\ex\label{ex2:psik19} 
\langinfo{Psikye (slightly) downward form of \textit{tli} `cut'}{}{\cite[115]{smithKapsikiLanguage1969}} \\
\gll  tli-yí le xwá \\
cut-\gl{down}2 with knife \\
\glt `Cut it with a knife (grass, for example)!'
\z 

The suffix \textit{-ve} means ``to go or to come across, to come up out of'' \citep[118]{smithKapsikiLanguage1969}. The `across' meaning is found with the root \textit{dze} `go', as in (\ref{ex2:psik10}), and the outward meaning is found with the root \textit{sé} `come', as in (\ref{ex2:psik11}). Note that this suffix is identical in form to a preposition \textit{ve} `to'. 

\ea\label{ex2:psik10}
\langinfo{Psikye outward form of \textit{dze} `go'}{}{\cite[118]{smithKapsikiLanguage1969}} \\
\gll kwembɛwalé ašɛ́ ji-vé te gwa nya\\
boat \gl{dem} go-\gl{out} on river \gl{dem}	\\
\glt `That boat is going across this river.'

\ex\label{ex2:psik11}
\langinfo{Psikye outward form of \textit{sé} `come'}{}{\cite[66]{smithKapsikiLanguage1969}} \\
\gll a ké-š-avé mbe cɛ	\\
\textsc{active} \gl{ptcp}-come-\gl{out} in house	\\
\glt `(He) came out of the house.'
\z 

The \textit{-ve} suffix is also found with non-motion verbs, without a motion meaning, as in (\ref{ex2:psik24}). See also the use of \textit{-ve} to express caused accompanied motion and to express buying, as discussed below. 

\ea\label{ex2:psik24}
\langinfo{Psikye outward form of \textit{mene} `do'}{}{\cite[114]{smithKapsikiLanguage1969}} \\
\gll  'a ké-mene-vé tlené	\\
\textsc{active} \gl{ptcp}-do-\gl{out} work	\\
\glt `He did the work for wages.'
\z

The inward suffix \textit{-mbe} (identical to the preposition \textit{mbe} `in') is only attested with two verbs. It occurs with the verb \textit{dze} `go', as in (\ref{ex2:psik12}) \citep[118]{smithKapsikiLanguage1969}, and it occurs with the `take' verb to express caused accompanied motion, as discussed below. 

\ea\label{ex2:psik12}
\langinfo{Psikye inward form of \textit{dze} `go'}{}{\cite[118]{smithKapsikiLanguage1969}} \\
\gll pa dze-mbé cɛ	\\
then go-\gl{into} house	\\
\glt `Then (he) went into the house...'
\z 

The bushward suffix \textit{-mté} is only found with a directional meaning when used with the verb root \textit{dze} `go', as in (\ref{ex2:psik13}). It is sometimes translated either as `towards the bush' or, more generally, `away'. It is also found with the `take' verb to express itive caused accompanied motion, as in (\ref{ex2:psik16}) below. 

\ea\label{ex2:psik13}
\langinfo{Psikye bushward form of \textit{dze} `go'}{}{\cite[118]{smithKapsikiLanguage1969}} \\
\gll ntíŋú 'yá ka-dze-mté \\
suddenly 1\gl{sg} \gl{ipfv}-go-\gl{bush}	\\
\glt `I abruptly went to the bush.'
\z 

In addition, a suffix \textit{-mte} with non-motion verbs, as in (\ref{ex2:psik20}), is said to mean that ``the action is finished (emphatically)'' \citep[114, 163]{smithKapsikiLanguage1969}. \citet[163]{smithKapsikiLanguage1969} notes that this suffix is identical in form to the word \textit{mté} `die' which can also be used in an expression to refer to the bush.  

\ea\label{ex2:psik20}
\langinfo{Psikye bushward form of \textit{séxwé} `wipe'}{}{\cite[163]{smithKapsikiLanguage1969}} \\
\gll 'a ke-séxwé-mte xwá	\\
\textsc{active} \gl{ptcp}-wipe-\gl{bush} knife	\\
\glt `(He) wiped the knife dry.'
\z 

Finally, there is some evidence that the word \textit{gɛ} `compound, home' may be grammaticalizing since when it occurs next to a basic motion verb it can trigger palatalization of the verb (<š> = [ʃ]), as seen in (\ref{ex2:psik14}).

\ea\label{ex2:psik14}
\langinfo{Psikye verb \textit{sé} `come' phonologically bound to \textit{gɛ} `home'}{}{\cite[118]{smithKapsikiLanguage1969}} \\
\gll ša=gɛ ŋkɛ́ le yɛmú \\
come=home \gl{poss}.3\gl{sg}.\gl{f} with water	\\
\glt `She is coming home with the water.'
\z 

As indicated in Table \ref{tab:psik} some of the directional extensions in Psikye are occasionally preceded by the form \textit{ak} which has no known function elsewhere. The \textit{ak} form never appears where the suffix has a directional interpretation, and it occurs most frequently with the verb `take' to express caused accompanied motion (as discussed below) or when the interpretation is subsequent CAM, as with the verb \textit{l} `dig' in (\ref{ex2:psik17}).\footnote{Smith translates the upward form in (\ref{ex2:psik17}) as downward subsequent motion (or CAM): ``...and brought it down.'' According to Psikye speaker Zra Kodji (personal communication), Smith's translation is an error and the correct translation should involve upward motion.} However, in one example the \textit{ak} form is found with an apparently intransitive verb \textit{nef} `boil', shown in (\ref{ex2:psik18}).

\ea\label{ex2:psik17}
\langinfo{Psikye (slighty) upward form of \textit{l} `dig'}{}{\cite[120]{smithKapsikiLanguage1969}} \\
\gll a 'yá ké-l-ák-ate \\
\textsc{active} 1\gl{sg} \gl{ptcp}-dig-\textsc{ak}-\gl{up}2	\\
\glt `I dug (it) and brought it up.'

\ex\label{ex2:psik18}
\langinfo{Psikye downward form of \textit{nef} `boil'}{}{\cite[120]{smithKapsikiLanguage1969}} \\
\gll  pa ka-nef-ak-áá	\\
then ??-boil-\textsc{ak}-\gl{down}1	\\
\glt `Then (it) boiled over.'
\z 

\subparagraph{Caused accompanied motion (CAM)}

There are several ways to express directed CAM in Psikye. Most prevalent in Smith's data is the use of directional suffixes with the verbs \textit{kélé} `take (\textsc{sg})' and \textit{fu} `take (\textsc{pl})'. In two examples, the verb occurs directly with a directional suffix (no use of \textit{-ak}), as in the `across' suffix in (\ref{ex2:psik15}) and the bushward suffix in (\ref{ex2:psik16}).

\ea\label{ex2:psik15}
\langinfo{Psikye outward form of \textit{fu} `take (plural)'}{}{\cite[119]{smithKapsikiLanguage1969}} \\
\gll  fu-vé wuši	\\
take.\gl{pl}-\gl{out} things \\
\glt `Bring the things.'

\ex\label{ex2:psik16}
\langinfo{Psikye bushward form of \textit{kélé} `take (singular)'}{}{\cite[119]{smithKapsikiLanguage1969}} \\
\gll pa kélé-mté wusu mɛ	\\
then take.\gl{sg}-\gl{bush} thing mouth	\\
\glt `Then she took off the lid.'
\z 

In other examples, the directional suffix is preceded by \textit{-ak} when occurring with one of the verbs glossed `take'. This is attested with the `across/out' suffix (as in \ref{ex2:psik21}), both of the upward suffixes, the (steep) downward suffix (\textsc{down1}) and the inward suffix (as in \ref{ex2:psik22}). Note that the use of \textit{-ak} in (\ref{ex2:psik21}) contrasts with the same verb root and suffix without \textit{-ak} in (\ref{ex2:psik15}). Also, the inward suffix in (\ref{ex2:psik22}) is given an unexplained outward motion translation. 

\ea\label{ex2:psik21}
\langinfo{Psikye outward form of \textit{fu} `take (plural)'}{}{\cite[120]{smithKapsikiLanguage1969}} \\
\gll  a 'yá kē-fu-ak-avɛ xa pelɛ́ melemé	\\
\textsc{active} 1\gl{sg} \gl{ptcp}-take.\gl{pl}-\textsc{ak}-\gl{out} millet over village		\\
\glt `I brought the millet over from the village.'
\z 

\ea\label{ex2:psik22}
\langinfo{Psikye inward form of \textit{kélé} `take (singular)'}{}{\cite[120]{smithKapsikiLanguage1969}} \\
\gll a 'yá ké-ḱl-aka-mbe šagá mbe cɛ	\\
\textsc{active} 1\gl{sg} \gl{ptcp}-take.\gl{sg}-\textsc{ak}-\gl{into} pot in house	\\
\glt `I took the pot out of the house.'
\z 

One example, shown in (\ref{ex2:psik23}), indicates that directed CAM can also be expressed by combining a dative argument (recipient) with a directional motion verb (used in a transitive context).

\ea\label{ex2:psik23}
\langinfo{Psikye directional motion verb with dative recipient}{}{\cite[117]{smithKapsikiLanguage1969}} \\
\gll  s-ame ɗá yɛmú kwálinɛ́	\\
come-\gl{up}1 1\gl{sg}.\gl{dat} water cold		\\
\glt `Bring me some cold water.'
\z 

\subparagraph{Buy-sell verbs}

The verb \textit{pá} (without a directional suffix) is found translated as both `buy' and `sell' \citep[60--61, 65, 80, 123]{smithKapsikiLanguage1969}. According to \citet[114]{smithKapsikiLanguage1969} ``The verb root \textit{pá} means `sell' when used with the suffix \textit{-mte} and 'buy' when used with the suffix \textit{-ve},'' as seen in (\ref{ex2:psik25}) and (\ref{ex2:psik26}). Note, however, that in one example it appears that the ventive suffix \textit{-ke} can also be used in the context of a buying event, as shown in (\ref{ex2:psik27}).

\ea\label{ex2:psik25}
\langinfo{Psikye bushward form of \textit{pá} `buy/sell'}{}{\cite[114]{smithKapsikiLanguage1969}} \\
\gll 'a 'ya ké-pá-mte marure kwélékwelé	\\
\textsc{active} 1\gl{sg} \gl{ptcp}-buy/sell-\gl{bush} rice much \\
\glt `I sold a lot of rice.'

\ex\label{ex2:psik26}
\langinfo{Psikye `across' form of \textit{pá} `buy/sell'}{}{\cite[114]{smithKapsikiLanguage1969}} \\
\gll  pá-ve ɗá wusú \\
buy/sell-\gl{out} 1\gl{sg} thing	\\
\glt `Buy me the thing... !'

\ex\label{ex2:psik27}
\langinfo{Psikye ventive form of \textit{pá} `buy/sell'}{}{\cite[80]{smithKapsikiLanguage1969}} \\
\gll  ka-pa-ké yɛ kelepé le gwené \\
\gl{ipfv}-buy/sell-\gl{vent} \gl{pl} fish with salt \\
\glt `...(she) will buy fish and salt.'
\z 

\subsection{Kotoko-Buduma (3/10)} 

Sufficient data is only available for three of the ten Kotoko-Buduma languages. No sufficient description was found of Mser (mser1242), Jina (jina1244),\footnote{In regards to Jina, the collection of articles in \citet{schmidtAspectsGrammarZina2002} do not explicitly address the issues relevant here, and the piecemeal approach to describing the language leaves a greater degree of uncertainty about possible gaps in the description of the grammar.} Majera (maje1243), Afade (afad1236)\footnote{For Afade, \citet[91]{solkenSeetzensAffadehBeitrag1967} mentions a verbal suffix \textit{-li} which is compared to the inward suffix in Lagwan, however, no clear examples were found.} or Maslam (masl1241). I also do not include the dormant and minimally-documented language, Ngala of Lake Chad (ngal1301).

\paragraph*{Buduma (budu1265)} \label{budu}

Information on Buduma grammar is from a journal article \citep{gaudicheLangueBoudouma1938} and two grammar sketches \citep{lukasSpracheBudumaIm1939,awaganaGrammatikBudumaPhonologie2001}.

\newpage
\subparagraph{Directional verbs}~

\begin{table}[H]
	\caption{Buduma directional verbs} 
	\label{tab:buduV}
	\begin{tabular}{lll} % add l for every additional column or remove as necessary
		\lsptoprule
		DIR & verb & gloss   \\ %table header
		\midrule
		ITIVE & \textit{wa lu} & `\textit{je vais} [I go]' \\
		VENT & \textit{ua aw} & `\textit{je viens} [I come]' \\
		UP & \textit{u lugu} & `\textit{je monte} [I go up]' \\
		DOWN & \textit{wa hay} & `\textit{je descends} [I descend]' \\
		INTO & \textit{i} & `\textit{hineingehen} [go in]' \\
		\lspbottomrule
	\end{tabular}
\end{table}

The directional verbs in Table \ref{tab:buduV} are from \citet{gaudicheLangueBoudouma1938} and \citet{lukasSpracheBudumaIm1939}.

\subparagraph{Directional extensions}~

\begin{table}[H]
	\caption{Buduma directional morphology}
	\label{tab:budu}
	\begin{tabular}{llll} % add l for every additional column or remove as necessary
		\lsptoprule
		forms & gloss & direction  & other functions  \\ %table header
		\midrule
		\textit{-(a)l}  &  \textsc{into}  &  \textsc{into} &    \\
		\textit{-gə̀rə́}  &  \textsc{onto}  & \textsc{onto}   &    \\
		\lspbottomrule
	\end{tabular}
\end{table} 

There are several verbal extensions in Buduma, a few of which have motion-related translations with at least some verbs. Based on the limited data available from \citet{lukasSpracheBudumaIm1939} and \citet{awaganaGrammatikBudumaPhonologie2001}, two suffixes seem to consistently have directional interpretations: \textit{-(a)l} `inward' and \textit{-gə̀rə́} `onto'. Most of the relevant examples are of transitive motion verbs, as in (\ref{ex2:budu1}) and (\ref{ex2:budu2}).

\ea\label{ex2:budu1}
\langinfo{Buduma inward form of \textit{ta} `put' [German: `\textit{stecken}']}{}{\cite[60]{lukasSpracheBudumaIm1939}} \\
\gll ta-l \\
put-\gl{into} \\
\glt `put in' [German: `\textit{hineintun}']

\ex\label{ex2:budu2}
\langinfo{Buduma `onto' form of \textit{tà} `set' [German: `\textit{setzen}']}{}{\cite[96]{awaganaGrammatikBudumaPhonologie2001}} \\
\gll tà-gə̀rə́ àtáw lwú	\\
\gl{imp}:set-\gl{onto} there ground	\\
\glt `Set (it) there on the ground.' [German: `\textit{Setz} (\textit{es}) \textit{dort auf den Boden}.']
\z

The inward suffix is also shown with a direction verb \textit{i} glossed `\textit{hineingehen} [go in]', \textit{y-al}, with no change in meaning \citep[60]{lukasSpracheBudumaIm1939}. Both \citet[60]{lukasSpracheBudumaIm1939} and \citet{awaganaGrammatikBudumaPhonologie2001} analyze the verb \textit{kol} or \textit{kul} `be' as ending with the inward suffix, but it is unclear what the root verb would be in this case.

\subparagraph{Caused accompanied motion (CAM)}

Caused accompanied motion can be expressed by the verb \textit{ru} `bring' \citep[49]{lukasSpracheBudumaIm1939}, and directed CAM can be expressed by the causative form of a motion verb, as in (\ref{ex2:budu3}).

\ea\label{ex2:budu3}
\langinfo{Buduma causative form of \textit{wú} `come'}{}{\cite[61]{lukasSpracheBudumaIm1939}} \\
\gll u-wú-re̤ \\
1\gl{sg}-come-\gl{caus} \\
\glt `I bring.' [German: `\textit{Ich bringe}.']
\z 

\subparagraph{Buy-sell verbs}

From various passing references in \citet{lukasSpracheBudumaIm1939}, there appear to be two separate verbs in Buduma: \textit{tu} `buy' and \textit{uhi} `sell'.

\paragraph*{Lagwan (lagw1237)} \label{lagw}

Information on Lagwan grammar is from a grammar sketch \citep{lukasLogoneSpracheImZentralen1936} and a lexicon \citep{shryockLagwanFrenchEnglish2020}.

\subparagraph{Directional verbs}~

\begin{table}[H]
	\caption{Lagwan directional verbs} 
	\label{tab:lagwV}
	\begin{tabular}{lll} % add l for every additional column or remove as necessary
		\lsptoprule
		DIR & verb & gloss   \\ %table header
		\midrule
		ITIVE & \textit{gɨr, lu} & `\textit{aller}, \textit{marcher} [go, walk]' \\
		VENT & \textit{lo} & `\textit{venir} [come]' \\	
		\lspbottomrule
	\end{tabular}
\end{table}

Upward, downward, inward and outward motion are expressed by a combination of the verb \textit{sa} and a directional adverb or extension: \textit{sa ghwaa, sa kal} `climb, ascend', \textit{sa watɨn} `descend, go down', \textit{sa li} `enter', \textit{sa fɨne} `exit' \citep{shryockLagwanFrenchEnglish2020}. \citet{shryockLagwanFrenchEnglish2020} do not give a gloss for the verb \textit{sa} on its own, but \citet[116]{lukasLogoneSpracheImZentralen1936} gives the German glosses `\textit{steigen}, \textit{treten} [climb, step]'.  

\newpage
\subparagraph{Directional extensions}~

\begin{table}[H]
	\caption{Lagwan directional morphology}
	\label{tab:lagw}
	\begin{tabular}{llll} % add l for every additional column or remove as necessary
		\lsptoprule
		forms & gloss & direction  & other functions  \\ %table header
		\midrule
		\textit{-li}  &  \textsc{into}  &  \textsc{into} &    \\
		\lspbottomrule
	\end{tabular}
\end{table} 

\citet[49--50]{lukasLogoneSpracheImZentralen1936} describes a verbal suffix \textit{-li} with inward directional meaning. All motion examples given are of transitive motion verbs, as in (\ref{ex2:lagw1}). The suffix is also used with a verb of perception, as in (\ref{ex2:lagw2}). 

\ea\label{ex2:lagw1}
\langinfo{Lagwan inward form of \textit{ks̰ì} `pour out' [German: `\textit{ausgießen}']}{}{\cite[50]{lukasLogoneSpracheImZentralen1936}} \\
\gll wá-ks̰ì-lí-ya \\
1\gl{sg}-pour-\gl{into}-??	\\
\glt `I poured in.' [German: `\textit{Ich habe eingegossen}.']

\ex\label{ex2:lagw2}
\langinfo{Lagwan inward form of \textit{ŋgù} `see'}{}{\cite[49]{lukasLogoneSpracheImZentralen1936}} \\
\gll wā́-ŋgù-lí \\
??-see-\gl{into}	\\
\glt `looked in' [German: `\textit{sah hinein}']
\z 

\subparagraph{Caused accompanied motion (CAM)}

The verb \textit{do} `bring' expresses caused accompanied motion \citep{shryockLagwanFrenchEnglish2020}. No further discussion was found of ways of expressing directed CAM in the verbal morphology. 

\subparagraph{Buy-sell verbs}

There are two separate verb roots in Lagwan: \textit{tuwa} `buy' and \textit{ghwi} `sell' \citep{shryockLagwanFrenchEnglish2020}.

\paragraph*{Makary Kotoko (Mpade, mpad1242)} \label{maka}

Makary Kotoko (or Mpade) is described by \citet{allisonAspectsGrammarMakary2012,allisonGrammarMakaryKotoko2020}. Allison describes ``locative particles'' which may have historically had a directional function, but no longer systematically modify verbs to add directional meaning. Directional meaning is instead expressed through adverbials. %Motion verbs in a multiverb construction appear to express purposive motion, rather than prior motion. Ventive and itive caused accompanied motion are lexicalized as separate verbs. There is one buy-sell verb, and the `sell' meaning of the verb can be reinforced by a causative form. 

\newpage
\subparagraph{Directional verbs}~

\begin{table}[H]
	\caption{Makary Kotoko directional verbs} 
	\label{tab:makaV}
	\begin{tabular}{lll} % add l for every additional column or remove as necessary
		\lsptoprule
		DIR & verb & gloss   \\ %table header
		\midrule
		ITIVE & \textit{də́} & `go' \\
		VENT & \textit{lū} & `come' \\
		INTO& \textit{só} & `enter' \\
		\lspbottomrule
	\end{tabular}
\end{table}

No downward/upward motion verbs and no outward motion verb (e.g., `exit') were seen in the data available. 

\subparagraph{Directional extensions}~

%\begin{table}[H]
%	\caption{Makary Kotoko directional morphology}
%	\label{tab:maka}
%	\begin{tabular}{llll} % add l for every additional column or remove as necessary
%		\lsptoprule
%		forms & source gloss & direction  & other functions  \\ %table header
%		\midrule
%		\textit{ní}  &  \textsc{l.p.}  &   ?? & N/A  \\
%		\textit{=ho}  &  \textsc{l.p.}  & lexicalized (< \textsc{vent/up} ?) &   ?? \\
%		\textit{yo}  &  \textsc{l.p.}  &   lexicalized (< \textsc{itive} ?) &  ?? \\
%		\textit{=he}  &  \textsc{l.p.}  &   lexicalized (< \textsc{down} ?) &   ??  \\
%		\lspbottomrule
%	\end{tabular}
%\end{table} 

\citet[239]{allisonGrammarMakaryKotoko2020} describes four ``locative particles'', all glossed \textsc{l.p.}. The particle \textit{ní} is said to only occur with four verbs \textit{də̄} ‘go’, \textit{do} ‘take (sth/s.o.) swh’, \textit{fō} ‘run’, and \textit{kə} ‘accompany', as in (\ref{ex2:maka3}) and (\ref{ex2:maka2}).

\ea\label{ex2:maka3}
\langinfo{Makary Kotoko verb \textit{də̄} ‘go’ with ``locative particle''}{}{\cite[247]{allisonGrammarMakaryKotoko2020}} \\
\gll ē də̄ ní =wo lábā nde wo dó \\
3\gl{pl}:\gl{compl} go \textsc{l.p.} =\gl{q} or be.at.\gl{pl} village \gl{det}.\gl{f} \\
\glt `Have they left or are they (still) in the village?'

\ex\label{ex2:maka2}
\langinfo{Makary Kotoko verb \textit{fō} ‘run’ with ``locative particle''}{}{\cite[247]{allisonGrammarMakaryKotoko2020}} \\
\gll kˈani= n-ō fō ní n-ō to ro-gə́-də\\
and= 3\gl{sg}.\gl{f}-\gl{compl} run \textsc{l.p.} 3\gl{sg}.\gl{f}-\gl{compl} return.home \gl{rel}.\gl{f}-\gl{poss}-3\gl{sg}.\gl{f}		\\
\glt `Then she ran away and returned (to her) home.'
\z 

The particle \textit{ní} is analyzed as ``a syntactic requirement for the four verbs in question when no other expression of location is given to indicate in which direction the action of the motion verb is carried out'' \citep[239]{allisonGrammarMakaryKotoko2020}.\footnote{\citet[267]{allisonGrammarMakaryKotoko2020} notes an exception in which ``the reciprocal marker fulfills the requirement of a locative complement for these verbs.''} The particle \textit{ní} is often used in the context of itive directional motion, but it does not necessarily add this meaning to the predicate. 

The other three ``locative particles'' combine with a wider set of verbs, and not necessarily in the context of a motion event. When one of these three particles combines with a verb, the meaning of the predicate is not predictable based on the meaning of the verb root. In other words, there are many cases in which the ``locative particles'' have co-lexicalized with a verb root. This is illustrated by the examples Allison provides for the verb \textit{gá} ‘put’ shown in (\ref{ex2:maka1}). 

\ea\label{ex2:maka1}
\langinfo{Makary Kotoko verb \textit{gá} ‘put’ with ``locative particles''}{}{\cite[240]{allisonGrammarMakaryKotoko2020}} \\
\ea
\gll  gá =he \\
put =\textsc{l.p.}	\\ 	
\glt `build'

\ex
\gll gá =ho	\\ 
put =\textsc{l.p.}	\\
\glt `wrap up, roll up'

\ex 
\gll gá yo	\\
put \textsc{l.p.}	\\
\glt `diminish, reduce (in quantity)'
\z
\z 

\citet[239]{allisonGrammarMakaryKotoko2020} proposes that these three particles' ``primary function... is to provide spatial orientation for the situation described'' suggesting that these are downward, ventive and itive directionals: ``\textit{=he} conveys a downward direction of the action of the verb, \textit{=ho} indicates that the action of the verb occurs toward some point of reference, and \textit{yo} that the action of the verb occurs moving away from the point of reference.''\footnote{The locative particles also function as adpositions in combination with a preposition, resulting in meanings translated as ‘by’ or ‘at’ \citep[171]{allisonGrammarMakaryKotoko2020}.} There is no evidence in the data provided that these are semantically productive, creating predictably composed meanings in combination with motion predicates, so Allison’s analysis is better understood as metaphorical or etymological. 

Directional meaning in Makary Kotoko is expressed by adverbial phrases such as \textit{ʦˈe} `outside' for outward motion, \textit{tə́n} `ground' for downward motion and \textit{wō} `summit' for upward motion \citep[208, 235]{allisonGrammarMakaryKotoko2020}.

\subparagraph{Caused accompanied motion (CAM)}

The itive-ventive directional distinction is lexicalized in the CAM verbs \textit{ɗō} `bring' and \textit{do} `take away'. 

\subparagraph{Buy-sell verbs}

The Makary Kotoko verb \textit{də́wo} `buy' expresses SELL in the causative form of the verb: \textit{də́wo-l} \citep[222--223]{allisonGrammarMakaryKotoko2020}. However, the \textit{-l} does not appear when the verb is followed by an object pronoun, and in this case the verb can mean `sell' without being in a causative form. 

\subsection{Lamang-Hdi (1/3)} 

Of the three Lamang-Hdi languages, only Lamang is included here.  No description is available for Vemgo-Mabas (vemg1240). Hdi has an elaborate system of motion related verbal morphology, however the major source on Hdi grammar \citep{frajzyngierGrammarHdi2002} cannot be taken at face value due to its ``ad hoc categories and highly creative terminology'' and ``many inconsistencies with regard to the interlinear translations''  \citep[215]{wolffLanguageVariationTheoretical2011}. 

\paragraph*{Lamang (lama1288)} \label{lama}

Lamang is described by \citet{wolffGrammarLamangLanguage1983,wolffEncodingTopographyDirection2006,wolffLamangLanguageDictionary2015}. Lamang has extensive motion morphology including upward{\slash}downward and outward{\slash}inward suffixes.

\subparagraph*{\textbf{Directional verbs}}
~
\begin{table}[H]
	\caption{Lamang directional verbs} 
	\label{tab:lamaV}
	\begin{tabular}{lll} % add l for every additional column or remove as necessary
		\lsptoprule
		DIR & verb & gloss   \\ %table header
		\midrule
		ITIVE & \textit{la-, dza-} & `go' \\
		VENT & \textit{sa-, skwa-} & `come' \\
		DOWN & \textit{tsgha} & `go down' \\
		\lspbottomrule
	\end{tabular}
\end{table}

\subparagraph{Directional extensions}~

\begin{table}[H]
	\caption{Lamang directional morphology}
	\label{tab:lamang}
	\begin{tabular}{llll} % add l for every additional column or remove as necessary
		\lsptoprule
		forms & source gloss & direction  & other functions  \\ %table header
		\midrule
		\textit{-VV-}, \textit{-gha}  &  \textsc{vent}  & \textsc{home} &  lexicalized \\
		\textit{-f(i)}  &  \textsc{up}  &     \textsc{up} &   \\
		\textit{-gaa}, \textit{-ɗi}  &  \textsc{down}    & \textsc{down} &   \\
		\textit{-ŋ}  &  \textsc{ill}  &  \textsc{into}   &  \\
		\textit{-b(i)}  &  \textsc{eff}  &   \textsc{out} &   \textsc{subs.out} \\
		\lspbottomrule
	\end{tabular}
\end{table}

There are four primary directional suffixes which \citet[140--145]{wolffLamangLanguageDictionary2015} calls ``locative-directional (‘topographic') extensions.''\footnote{Elsewhere, Wolff refers to a larger category of suffixes with categorically distinct semantic properties \citep[150--153, 158, 161]{wolffLamangLanguageDictionary2015}. Here, only the shorter list of four morphemes are considered directional extensions.} These suffixes can vary in their form depending on the class of verb they occur with.\footnote{These morphemes will appear as suffixes in all examples shown here, but in the ``Subjunctive II'' form of Lamang verbs, they appear as prefixes \citep[131, 174]{wolffLamangLanguageDictionary2015}.} The most relevant verb class is formed of just two verbs, \textit{dza-} ‘go’ and \textit{skwa-} ‘come’ which cannot occur without a directional suffix \citep[139]{wolffLamangLanguageDictionary2015}. These two verbs have the exceptional form of the homeward suffix \textit{-gha}, as in (\ref{ex2:lama1}) and (\ref{ex2:lama2}), instead of a long vowel as with other verbs \citep[140]{wolffLamangLanguageDictionary2015}.

\ea\label{ex2:lama1}
\langinfo{Lamang verb \textit{dza-} ‘go’ with homeward suffix}{}{\cite[141]{wolffLamangLanguageDictionary2015}} \\
\label{lama1}
\glll dzághà \\
dza-́ghà \\
go-\gl{home}  \\
\glt `going home, going to a populated place'
\z

\ea\label{ex2:lama2}
\langinfo{Lamang verb \textit{skwa-} ‘come’ with homeward suffix}{}{\cite[141]{wolffLamangLanguageDictionary2015}} \\
\glll skwághà \\
skwa-́ghà \\
come-\gl{home} \\
\glt `coming home; coming here, hither'
\z

Wolff calls the homeward suffixes ``ventive'' but the suffixes are semantically distinct from ventive directionals in other Chadic languages. The verb in (\ref{ex2:lama1}) has an itive translation, so the fact that it can unproblematically combine with the \textit{-gha} suffix indicates that the suffix does not express the semantically contradictory notion of ventive motion (see also \ref{ex2:lama333}). The meaning of the suffix is described as ``movement into and/or towards an inhabited or sacred place'' \citep[140]{wolffLamangLanguageDictionary2015} and is elsewhere explicitly said to be neither ventive nor itive \citep[234]{wolffEncodingTopographyDirection2006}. Rather, this suffix appears more semantically similar to the homeward directional extension in Bura. With verbs other than \textit{dza-} ‘go’ and \textit{skwa-} ‘come’, the homeward form of the verb is marked by vowel lengthening and a high-low tone pattern \citep[151]{wolffLamangLanguageDictionary2015}. This form is used with the  motion verbs \textit{la-} `go' (as in \ref{ex2:lama3}) and \textit{sa-} `come'.

\ea
\langinfo{Lamang verb \textit{la-} ‘go’ with homeward suffix}{}{\cite[211]{wolffLamangLanguageDictionary2015}} \\
\ea\label{ex2:lama333}
\gll 
gú plís dzá-ghà dáa Wándàlà \\
\gl{narr} horse go-\gl{home} \gl{prep} Wandala \\
\glt `And then the horse(man) went home to Wandala.' 

\ex\label{ex2:lama3}
\gll \textbf{l-áa-kéghə̀-ɗ} dáa Wándàl ná ámghàm ká-kà \\
go-\gl{home}-and-\gl{3}\gl{sg} \gl{prep} Wandala \gl{top} chief \gl{quot}-\gl{3}\gl{sg}	 \\
\glt `When he had gone home to Wandala, ‘Sultan!’ he said.'
\z
\z 

Three other suffixes have a directional function when combined with motion verbs: the outward suffix \textit{-b(i)} (Wolff calls this ``efferential'' glossed \textsc{eff}), as in (\ref{ex2:lama6}), the upward suffix \textit{-f(i)} (in \ref{ex2:lama7}) and the downward suffix \textit{-gaa} (in \ref{ex2:lama8}).\footnote{The downward suffix is \textit{-ɗi} with the two irregular motion verbs \textit{dza-} ‘go’ and \textit{skwa-} ‘come’.}

\ea\label{ex2:lama6}
\langinfo{Lamang verb \textit{s-} ‘come’ with outward suffix}{}{\cite[223]{wolffLamangLanguageDictionary2015}} \\
\gll 
hà-ɗ únd-áa sìɗ kə̀ sə́-b-tà gùléŋ-úwó sé Kádù ŋ̀gáná-yá	\\
\gl{neg}-3\gl{sg} person-\gl{asoc} certain and come-\gl{out}-\gl{compl} also-\gl{neg} only Kadu \gl{dem}-\gl{dist}		\\
\glt `Again there was nobody to come out, only that Kadu.'
\z

\ea\label{ex2:lama7}
\langinfo{Lamang verb \textit{s-} ‘come’ with upward suffix}{}{\cite[273]{wolffLamangLanguageDictionary2015}} \\
\gll 
skwé-f-á-mb-ìn ná sə́-f-gə̀ Mə́lgwà-hàŋ	\\
come-\gl{up}-\gl{asoc}-\gl{coll}-3\gl{sg}.\gl{poss} \gl{top} come-\gl{up}-\gl{prep} Mulgwi-3\gl{pl}.\gl{excl}	\\
\glt `Their (collective) arrival to up here, they came up from Mulgwi.'
\z

\ea\label{ex2:lama8}
\langinfo{Lamang verb \textit{s-} ‘come’ with downward suffix}{}{\cite[144]{wolffLamangLanguageDictionary2015}} \\
\gll 
sa-gáa-tá	\\
come-\gl{down}-\gl{compl} \\
\glt `coming down'
\z

There are a few examples of the outward suffix used with transitive verbs in which the object can be interpreted as a figure on a path of motion, as in (\ref{ex2:lama10}) and (\ref{ex2:lama11}).

\ea\label{ex2:lama10}
\langinfo{Lamang verb \textit{pás-} ‘scrape’ with outward suffix}{}{\cite[277]{wolffLamangLanguageDictionary2015}} \\
\gll 
vítá pásá-p-pása-l-iya	\\
when \gl{pfv}.scrape-\gl{out}-scrape-3\gl{pl}.\gl{incl}-\gl{dist}	\\
\glt `When one has scraped (it) out...'
\z 

\ea\label{ex2:lama11}
\langinfo{Lamang verb \textit{ɓìts-} ‘press’ with outward suffix}{}{\cite[332]{wolffLamangLanguageDictionary2015}} \\
\gll
ɓìtsə̀-b-kégə̀-ɗè t’ síf-ín-ìyà	\\
press-\gl{out}-\gl{seq}-3\gl{sg}.\gl{excl} \gl{obj} liquid-3\gl{sg}.\gl{poss}-\gl{dist}		\\
\glt `When she has pressed out its \textit{sifi} (the sieved-through liquid)...'
\z

With the majority of verbs, directional forms do not appear to have any motion meaning, as in (\ref{ex2:lama4}), (\ref{ex2:lama13}) and (\ref{ex2:lama5}). There were no examples found of a non-motion verb with the downward suffix.

\ea\label{ex2:lama4}
\langinfo{Lamang verb \textit{d-} ‘cook’ with homeward morphology}{}{\cite[200]{wolffLamangLanguageDictionary2015}} \\
\gll 
víté-yá ŋgáné tá-d-áa-tə̀-l tə̀ {mə́gh dzásl}	 \\
day-\gl{dist} \gl{dem} \gl{ipfv}-cook-\gl{vent}-\gl{compl}-3\gl{pl}.\gl{incl} ?? {magh dzasla}	\\
\glt `On that day, \textit{magh dzasla} (i.e., beer mash) is prepared.'
\z

\ea\label{ex2:lama13}
\langinfo{Lamang verb \textit{hàn-} ‘slaughter’ with outward morphology}{}{\cite[246]{wolffLamangLanguageDictionary2015}} \\
\gll 
ùnd-úndà hànə́-b tə́ Ighə̀ŋ-úw ná	\\
person-\gl{rel} slaughter.\gl{pl}-\gl{out} \gl{obj} bull-\gl{neg} \gl{top}		\\ 
\glt `A person who had not killed a bull...'
\z

\ea\label{ex2:lama5}
\langinfo{Lamang verb \textit{yə-} ‘give  birth’ with upward morphology}{}{\cite[411]{wolffLamangLanguageDictionary2015}} \\
\gll 
gú lə̀-ŋ-Də̀gháh yə́-f-tá tə́ Kámbà	\\
\gl{narr} people-\gl{prep}-Dghaha give.birth-\gl{up}-\gl{compl} \gl{prep} Kamba		\\
\glt `The Dghaha people multiplied at Kamba.'
\z 

In addition to the above directional suffixes, there are also a few other suffixes with possible directional interpretations. The inward suffix \textit{-ŋ} (with a high-low tonal pattern) is called the ``illative extension'', glossed \textsc{ill} \citep[153]{wolffLamangLanguageDictionary2015}. The suffix can have an inward interpretation, but only in reference to an object, as in (\ref{ex2:lama18}). This suffix is not found with intransitive motion verbs, and with many transitive verbs there is no indication of any motion meaning related to the suffix, as in (\ref{ex2:lama19}). %In all of these characteristics, the inward suffix is more similar to the ``additive extensions'' \citep[158]{wolffLamangLanguageDictionary2015} than it is to the directional suffixes.

\ea\label{ex2:lama18}
\langinfo{Lamang verb \textit{dz-} ‘beat’ with inward suffix}{}{\cite[265]{wolffLamangLanguageDictionary2015}} \\
\gll 
ɘ̀=ŋ-sə́rá-ta=haŋ=uwo	\\
beat=\gl{into}-foot-\gl{compl}=3\gl{pl}=\gl{neg} \\
\glt `They did not put a foot into it.'
\z 

\ea\label{ex2:lama19}
\langinfo{Lamang verb \textit{kùzl-} ‘feel pain’ with inward suffix}{}{\cite[270]{wolffLamangLanguageDictionary2015}} \\
\gll 
ɓátsá-ɓátsá kùzlə́-ŋ-t-âa ín-áná kéy ná	\\
\gl{pfv}.diminish-diminish feel.pain-\gl{into}-\gl{compl}-\gl{asoc} thing-\gl{prox} \gl{dem} \gl{top}		\\
\glt `(As) the pain from that thing (i.e., scorpion’s sting) had diminished...'
\z 

Finally, \citet[159]{wolffLamangLanguageDictionary2015} describes a vowel lengthening process with high tone called the ``reductive-causative-downward extension''. This form is said to be used to express downward directionality with the caused accompanied motion (CAM) forms of motion verbs, but other than the isolated illustrative examples given, elsewhere the CAM motion verbs use the downward suffix \textit{-ɗ}, as in (\ref{ex2:lama20}) \citep[see also][271]{wolffLamangLanguageDictionary2015}. In other examples, the function of this vowel lengthening is opaque with no indication of a motion-related function in the translation. 

Another particularity of the Lamang directional suffixes is that the data available does not show them freely being used with manner-of-motion verbs. Rather, the directionality of a manner-of-motion verbs is primarily expressed through a multiverb construction (MVC) which \citet[146]{wolffLamangLanguageDictionary2015} calls the ``adverbial use'' of the motion verbs \textit{dza-} ‘go’ and \textit{skwa-} ‘come’. Like serial verb constructions, these constructions do not use any morphological marking of dependency relations between the verbs. In example (\ref{ex2:lama14}), the verb `run' is followed by the itive motion verb with an outward directional suffix indicating that the running occurred along a path both away from the deictic center and outward from an enclosed space.

\ea\label{ex2:lama14}
\langinfo{Lamang verb \textit{rg-} ‘run’ in a directional MVC}{}{\cite[143]{wolffLamangLanguageDictionary2015}} \\
\gll 
ɗíŋ gà-háŋ ná tə́-rg yâgh dzè-b	\\
\gl{ideo} \gl{quot}-3\gl{pl} \gl{top} \gl{ipfv}-run squirrel go-\gl{out}	\\
\glt `(While) they looked closely, squirrel kept running away (out of a hole in the ground).'
\z

A similar construction is seen in example (\ref{ex2:lama15}). In this case, the manner-of-motion verb `fly' has the same directional suffix (\textsc{down}) as the following directional verb. This is the only context where a manner-of-motion verb is seen with a directional suffix in the data examined.

\ea\label{ex2:lama15}
\langinfo{Lamang verb \textit{ndr-} ‘fly’ in a directional MVC}{}{\cite[283]{wolffLamangLanguageDictionary2015}} \\
\gll ɗíyák ndə̀rà-gáa-tá skwè-ɗè	\\
bird fly-\gl{down}-\gl{compl} come-\gl{down}	\\
\glt `The bird flew down.'
\z 


\subparagraph{Caused accompanied motion (CAM)}
\largerpage[2]

There are at least three different possible strategies for expressing caused accompanied motion (CAM) in Lamang: a causative form of motion verbs, the use of directional morphology with the verb \textit{ksa-} `catch', and the use of the verb \textit{kl-} ‘take’ in a directional multiverb construction. 

The motion verbs  \textit{la-} `go' and \textit{sa-} `come' have a CAM form in which a suffix \textit{-ŋa} is added to the stem. Wolff calls this a ``causative'' suffix, but this form is limited to these two verbs and apparently only used to express CAM. The CAM forms behave in the same way as the motion verbs when occurring with any of the four directional suffixes, as in (\ref{ex2:lama9}). 

\ea\label{ex2:lama9}
\langinfo{Lamang verb \textit{s-} ‘come’ with causative and directional suffixes}{}{\cite[142]{wolffLamangLanguageDictionary2015}} \\

\ea
\gll 
sá-ŋá-ghà \\
come-\gl{caus}-\gl{home} \\
\glt `bringing home (coming)'

\ex
\gll
sá-ŋé-bè \\
come-\gl{caus}-\gl{out} \\
\glt `bringing out (coming)'

\ex
\gll
sá-ŋé-fé \\
come-\gl{caus}-\gl{up} \\
\glt `bringing up (coming)'

\ex \label{ex2:lama20}
\gll
sá-ŋé-ɗé \\
come-\gl{caus}-\gl{down} \\
\glt `bringing down (coming)'
\z 
\z 

It is also possible to derive a CAM meaning by the use of directional morphology with the verb \textit{ksa-} `catch', as in (\ref{ex2:lama12}).

\ea\label{ex2:lama12}
\langinfo{Lamang verb \textit{ksa-} ‘catch’ with directional suffix and CAM meaning}{}{\cite[78, 153]{wolffLamangLanguageDictionary2015}} \\

\ea
\gll
káa-ksà	\\
\gl{pfv}.catch.\gl{home}-catch	\\
\glt `have caught (and brought hither)'

\ex
\gll
ksa-gáa-tá	\\
catch-\gl{down}-\gl{compl}		\\
\glt `catching and taking downhill'
\z 
\z

\largerpage
The directional MVC described above can be used to express CAM when the first verb is \textit{kl-} ‘take’. The second verb specifies both a deictic orientation (ventive or itive) and another directional feature, such as the outward motion in (\ref{ex2:lama16}).

\ea\label{ex2:lama16}
\langinfo{Lamang verb \textit{kl-} ‘take’ in a directional MVC}{}{\cite[260]{wolffLamangLanguageDictionary2015}} \\
\gll hə́má bò-l t-kə́l skwé-b m-d-e wo	\\
repeat \gl{neg}.\gl{foc}-2\gl{pl}.\gl{incl} \gl{ipfv}-take come-\gl{out} \gl{prep}-3\gl{sg}-\gl{dist} \gl{neg} \\
\glt `Never again will one take (it) out from within it.'
\z 

\subparagraph{Buy-sell verbs}

Two separate verbs are used in Lamang: \textit{skwa} ‘buy’ and \textit{dzawa} `sell'. Note that these forms are very similar to the verb roots \textit{skwa-} `come' and \textit{dza-} `go'. 

\subsection{Margi-Mandara-Mofu (17/25)}

The Margi-Mandara-Mofu languages form the largest group in this classification of Chadic languages. They are further divided into three subgroups: Bura-Marghi, Mandaraic and Mofuic. 

\subsubsection{Bura-Marghi (4/7)}

The Bura-Margi languages can be further split into two clusters: four Buraic languages and three Marghic languages. Of the four Buraic languages, Bura and Cibak are included here. No information was found for Nggwahyi (nggw1242) or Putai (puta1243). Of the three Marghic languages, Huba (Kilba) and Margi Central are included below, but no information was found for Margi South (marg1266).

\paragraph*{Bura (bura1292)} \label{bura}

Bura grammar was described by \citet{hoffmannUntersuchungenZurStruktur1955} then again by \citet{blenchBuraVerbalExtensions2010} without reference to Hoffmann's data.

\subparagraph{Directional verbs}~

\begin{table}[H]
	\caption{Bura directional verbs} 
	\label{tab:buraV}
	\begin{tabular}{lll} % add l for every additional column or remove as necessary
		\lsptoprule
		DIR & verb & gloss   \\ %table header
		\midrule
		ITIVE & \textit{li} & `go' \\
		VENT & \textit{si} & `come' \\
		\lspbottomrule
	\end{tabular}
\end{table}

\largerpage[2]
\subparagraph{Directional extensions}~

\begin{table}[H]
	\caption{Bura directional morphology}
	\label{tab:bura}
	\begin{tabular}{llll} % add l for every additional column or remove as necessary
		\lsptoprule
		forms & gloss & direction  & other functions  \\ %table header
		\midrule
		\textit{-mta}  &  \textsc{itive}  &  \textsc{itive} &   lexicalized  \\
		\textit{-hi}  &  \textsc{down}  &   \textsc{down}  &    \\
		\textit{-wa, -kwa, -u}  &  \textsc{into}  & \textsc{into}  &    \\
		\textit{-bila, -ɓəla}  &  \textsc{out}  &  \textsc{out} & \textsc{compl}  \\
		\textit{-nkir, -nkər}  &  \textsc{onto}  &   \textsc{onto}  &  \textsc{loc, add} \\		
		\textit{-vi}  &  \textsc{home}  &   \textsc{home}  &   \\	
		\textit{-dza}  &  \textsc{side}  &   \textsc{side}  & \textsc{loc}, lexicalized  \\			
		\lspbottomrule
	\end{tabular}
\end{table} 

The inventory of motion-related verbal suffixes in Bura is fairly large, and stands out in not having a ventive suffix.\footnote{\citet[3]{blenchBuraVerbalExtensions2010} lists extensions with the verb `go' of which several only appear once in the data: \textit{-kəra} `under', \textit{-ma} `somewhere', \textit{-nda} `yonder'. The suffix \textit{-kəma} `forward' is apparently the same as what \citet[284]{hoffmannUntersuchungenZurStruktur1955} transcribes as \textit{-kuma}. Hoffmann gives a few examples of this form, proposes that it is related to the noun \textit{kuma} `face' and concludes ``\textit{Es bleibt aber fraglich, ob es sich bei} -kuma \textit{um ein echtes Suffix handelt}... [However, it remains questionable whether -\textit{kuma} is a real suffix...]'' Blench's list also includes the suffix \textit{-ha} which is translated as `to an inhabited place, home' but is elsewhere described as meaning `together, collectively'.} The itive suffix \textit{-mta} is attested with an itive interpretation when used with transitive motion verbs, as in (\ref{ex2:bura1}) and (\ref{ex2:bura2}). 

\ea\label{ex2:bura1}
\langinfo{Bura verb \textit{ndzi-} ‘throw’ with itive suffix}{}{\cite[287]{hoffmannUntersuchungenZurStruktur1955}} \\
\gll ndzi-mta \\
throw-\gl{itive}	\\
\glt `to throw away' [German: `\textit{wegwerfen}']

\ex\label{ex2:bura2}
\langinfo{Bura verb \textit{carə-} ‘knock down’ with suffix}{}{\cite[6]{blenchBuraVerbalExtensions2010}} \\
\gll carə-mta \\
knock.down-\gl{itive} \\
\glt `to be knocked away from you'
\z 

\citet[289]{hoffmannUntersuchungenZurStruktur1955} also records the form \textit{limta} `go up, rise' from the verb \textit{li} `go'. According to Hoffmann, this is not the itive suffix, but a phonetic contraction of the verb with the word \textit{amta} `above' and etymologically unrelated to the itive suffix \textit{-mta}. More research would be needed to verify that analysis. 

The inward suffix has several forms, \textit{-kwa}, \textit{-wa}, and \textit{-u}. The first form in particular is similar to the preposition \textit{akwa} `in', and as such \citet[285]{hoffmannUntersuchungenZurStruktur1955} raises doubts about whether \textit{-kwa} should be treated as a verbal suffix at all. Notably, the \textit{-kwa} form is the only one found with intransitive motion verbs, as in (\ref{ex2:bura3}). 

\ea\label{ex2:bura3}
\langinfo{Bura verb \textit{li-} ‘go’ with inward suffix}{}{\cite[285]{hoffmannUntersuchungenZurStruktur1955}} \\
\gll lu-kwa \\
go-\gl{into}	\\
\glt `go into' [German: `\textit{hineingehen in}']
\z 

The \textit{-wa} form appears to be the most broadly used, but both \textit{-wa} and \textit{-u} are found with directional inward meaning when occurring with transitive motion verbs, as in (\ref{ex2:bura4}) and (\ref{ex2:bura15}). \citet{blenchBuraVerbalExtensions2010} characterizes the use of \textit{-wa} with non-motion verbs as meaning `to fix, repair, repay' although this pattern is only clear for a handful of verbs \citep[see also][320]{hoffmannUntersuchungenZurStruktur1955}.  

\ea\label{ex2:bura4}
\langinfo{Bura verb \textit{psu-} ‘throw’ with inward suffix}{}{\cite[301]{hoffmannUntersuchungenZurStruktur1955}} \\
\gll psu-wa \\
throw-\gl{into}	\\
\glt `throw in' [German: `\textit{hineinwerfen}']

\ex\label{ex2:bura15}
\langinfo{Bura verb \textit{fi-} ‘blow’ with inward suffix}{}{\cite[300]{hoffmannUntersuchungenZurStruktur1955}} \\
\gll fi-u \\
blow-\gl{into} \\
\glt `to blow in, to breathe in' [German: `\textit{einblasen, einhauchen}']
\z 

A similar pattern is seen with the outward suffix, \textit{-bila} (in Hoffmann) or \textit{-ɓəla} (in Blench), as in (\ref{ex2:bura5}) and (\ref{ex2:bura16}). Both \citet[274]{hoffmannUntersuchungenZurStruktur1955} and \citet[3]{blenchBuraVerbalExtensions2010} generalize the use of the outward suffix with non-motion verbs as meaning `thoroughly, completely' and note the similar extended use of \textit{out} in English. 

\ea\label{ex2:bura5}
\langinfo{Bura verb \textit{si-} ‘come’ with outward suffix}{}{\cite[274]{hoffmannUntersuchungenZurStruktur1955}} \\
\gll  si-bila \\
come-\gl{out} \\
\glt `come out' [German: `\textit{herauskommen}']

\ex\label{ex2:bura16}
\langinfo{Bura verb \textit{buka-} ‘push’ with outward suffix}{}{\cite[4]{blenchBuraVerbalExtensions2010}} \\
\gll  buka-ɓəla \\
push-\gl{out} \\
\glt `to push outside'
\z

The downward suffix \textit{-hi} can have a directional function with both intransitive and transitive motion verbs, as in (\ref{ex2:bura6}) and (\ref{ex2:bura17}). \citet[282]{hoffmannUntersuchungenZurStruktur1955} notes that the downward suffix is identical in form to the noun \textit{hi} `earth, soil'. 

\ea\label{ex2:bura6}
\langinfo{Bura verb \textit{si-} ‘come’ with downward suffix}{}{\cite[281]{hoffmannUntersuchungenZurStruktur1955}} \\
\gll si-hi \\
come-\gl{down} \\
\glt `come down' [German: `\textit{herunterkommen}']

\ex\label{ex2:bura17}
\langinfo{Bura verb \textit{ndzi-} ‘throw’ with downward suffix}{}{\cite[281]{hoffmannUntersuchungenZurStruktur1955}} \\
\gll ndzi-hi \\
throw-\gl{down} \\
\glt `throw down' [German: `\textit{niederwerfen}']
\z 

The `onto' suffix \textit{-nkir, -nkər} is mostly found with transitive motion verbs (like the itive suffix), as in (\ref{ex2:bura7}), but one example is given of an apparent use of this suffix with an intransitive motion verb, (\ref{ex2:bura8}). It occurs with several non-motion verbs without a clear pattern in its semantic contribution. \citet[292]{hoffmannUntersuchungenZurStruktur1955} notes that the suffix is similar in form to the noun \textit{kir} `head'. 

\ea\label{ex2:bura7}
\langinfo{Bura verb \textit{ndzihi-} ‘throw’ with `onto' suffix}{}{\cite[291]{hoffmannUntersuchungenZurStruktur1955}} \\
\gll ndzihi-nkir \\
throw-\gl{onto} \\
\glt `throw on' [German: `\textit{werfen auf}']

\ex\label{ex2:bura8}
\langinfo{Bura verb \textit{si-} ‘come’ with `onto' suffix}{}{\cite[7]{blenchBuraVerbalExtensions2010}} \\
\gll ndla-nkər \\
fall-\gl{onto} \\
\glt `to fall on top of something'
\z 

Both Hoffmann and Blench describe a verbal suffix \textit{-vi} which has a homeward meaning when used with intransitive or transitive motion verbs, as in (\ref{ex2:bura9}) and (\ref{ex2:bura10}).\footnote{Hoffmann takes his examples from a Bura translation of the New Testament, so the context of this example is a story in which friends of a sick man carry his stretcher onto the roof of a crowded house and remove part of the roof in order to lower the friend into the house on a stretcher, thereby gaining access to Jesus who miraculously heals the man.} It is also shown with non-motion verbs in which case its semantic contribution is not transparent, as in (\ref{ex2:bura11}). \citet[301]{hoffmannUntersuchungenZurStruktur1955} also describes a word \textit{avi} `in the house, at home' and \citet[7]{blenchBuraVerbalExtensions2010} cites a noun \textit{vi} `place'. 

\ea\label{ex2:bura9}
\langinfo{Bura verb \textit{li-} ‘go’ with homeward suffix}{}{\cite[301]{hoffmannUntersuchungenZurStruktur1955}} \\
\gll li-vi \\
go-\gl{home} \\
\glt `go in, go home' [German: `\textit{hineingehen, heimgehen}']

\ex\label{ex2:bura10}
\langinfo{Bura verb \textit{ha-} ‘put’ with homeward suffix}{}{\cite[301]{hoffmannUntersuchungenZurStruktur1955}} \\
\gll ha-vi \\ 
put-\gl{home} \\
\glt `let down (into the house)' [German: `\textit{hinunterlassen} (\textit{scil. ins Haus})']

\ex\label{ex2:bura11}
\langinfo{Bura verb \textit{ntsi-} ‘squeeze’ with homeward suffix}{}{\cite[7]{blenchBuraVerbalExtensions2010}} \\
\gll ntsi-vi \\
squeeze-\gl{home}\\
\glt `to crowd out'
\z

Finally, Blench discusses a suffix \textit{-dza} `beside'. This suffix also occurs with at least one intransitive motion verb, as in (\ref{ex2:bura12}), as well as transitive motion verbs, as in (\ref{ex2:bura13}). It also occurs with non-motion verbs in which case it does not have a straightforward directional meaning, as in (\ref{ex2:bura14}). \citet[275]{hoffmannUntersuchungenZurStruktur1955} also discusses the \textit{-dza} suffix but describes its meaning as more generally locative, and etymologically related to a preposition \textit{adza} `to, at'. Hoffmann raises doubts about the analysis of \textit{-dza} as a suffix and explicitly says that the form \textit{lidza}, as in (\ref{ex2:bura12}), is not a suffix but merely a phonetic amalgamation of the verb and preposition. 

\ea\label{ex2:bura12}
\langinfo{Bura verb \textit{li-} ‘go’ with `beside' suffix}{}{\cite[4]{blenchBuraVerbalExtensions2010}} \\
\gll li-dza \\
go-\textsc{side} \\
\glt `to go close to'

\ex\label{ex2:bura13}
\langinfo{Bura verb \textit{bwal-} ‘move’ with `beside' suffix}{}{\cite[4]{blenchBuraVerbalExtensions2010}} \\
\gll bwal-dza \\
move-\textsc{side} \\
\glt `to place an object beside another'

\ex\label{ex2:bura14}
\langinfo{Bura verb \textit{dzib-} ‘to fill to repletion with food’ with `beside' suffix}{}{\cite[4]{blenchBuraVerbalExtensions2010}} \\
\gll dzib-dza \\
fill.w.food-\textsc{side} \\
\glt `to plaster a mat house'
\z 

\subparagraph{Caused accompanied motion (CAM)}

There are several ways to express CAM meaning in Bura. In addition to verbs glossed `bring', there is also a causative form of the motion verb \textit{si} `come', \textit{sinta} `bring (here)' \citep[294]{hoffmannUntersuchungenZurStruktur1955}. Finally, there are also several examples of the verb \textit{fa} `take' translated with directed CAM meaning when combined with a directional suffix, as in (\ref{ex2:bura18}), (\ref{ex2:bura19}) and (\ref{ex2:bura20}).

\ea\label{ex2:bura18}
\langinfo{Bura verb \textit{fa-} ‘take’ with itive suffix}{}{\cite[287]{hoffmannUntersuchungenZurStruktur1955}} \\
\gll fa-mta \\
take-\gl{itive} \\
\glt `take away' [German: `\textit{wegnehmen}']

\ex\label{ex2:bura19}
\langinfo{Bura verb \textit{fa-} ‘take’ with inward suffix}{}{\cite[301]{hoffmannUntersuchungenZurStruktur1955}} \\
\gll fa-wa \\
take-\gl{into} \\
\glt `take (and put in another place)' [German: `\textit{nehmen} (\textit{und an einen andern Platz legen})']

\ex\label{ex2:bura20}
\langinfo{Bura verb \textit{fa-} ‘take’ with outward suffix}{}{\cite[273]{hoffmannUntersuchungenZurStruktur1955}} \\
\gll  fa-bila \\
take-\gl{out} \\
\glt `take out, get out' [German: `\textit{herausnehmen, -holen}']
\z

\subparagraph{Buy-sell verbs}

There are two separate verbs in Bura: \textit{masa} `buy' and \textit{dila} `sell'. The verbs can take various extensions without changing their root meaning, as \textit{dilə-mta} `to sell out, to sell away' \citep[6]{blenchBuraVerbalExtensions2010}.

	\paragraph*{Cibak (ciba1236)} \label{ciba}

Information on Cibak grammar is from a 29-page sketch \citep{hoffmannZurSpracheCibak1955}. 

\subparagraph{Directional verbs}~

\begin{table}[H]
	\caption{Cibak directional verbs} 
	\label{tab:cibaB}
	\begin{tabular}{lll} % add l for every additional column or remove as necessary
		\lsptoprule
		DIR & verb & gloss   \\ %table header
		\midrule
		ITIVE & \textit{li} & `go' \\
		VENT & \textit{si} & `come' \\
		\lspbottomrule
	\end{tabular}
\end{table}

\subparagraph{Directional extensions}~

\begin{table}[H]
	\caption{Cibak directional morphology}
	\label{tab:ciba}
	\begin{tabular}{llll} % add l for every additional column or remove as necessary
		\lsptoprule
		forms & source gloss & direction  & other functions  \\ %table header
		\midrule
%		\textit{-nta}  &  \textsc{itive}  &  ?? &   ??  \\
		\textit{-ba}  &  \textsc{out}  &  \textsc{out}  &   \\
		\lspbottomrule
	\end{tabular}
\end{table} 

Just two examples are given of the outward suffix \textit{-ba}. One is with a buy/sell verb (see below) and the other with the ventive motion verb \textit{si} `come', shown in (\ref{ex2:ciba2}). 

\ea\label{ex2:ciba2}
\langinfo{Cibak verb \textit{si} ‘come’ with outward suffix}{}{\cite[133]{hoffmannZurSpracheCibak1955}} \\
\gll su-ba  \\
come-\gl{out} \\
\glt `come out'
\z

In addition, the form \textit{-nta} is compared to the Bura itive suffix \textit{-mta} \citep[134]{hoffmannZurSpracheCibak1955}. Six examples are given, including three transitive motion verbs, but none of the translations give an explicit directional meaning. In (\ref{ex2:ciba1}), Hoffman includes in the free translation a comparison to a similar form in Margi which is explicitly translated with directional meaning.

\ea\label{ex2:ciba1}
\langinfo{Cibak verb \textit{dàl} ‘throw’ with itive suffix}{}{\cite[137]{hoffmannZurSpracheCibak1955}} \\
\gll dàl-ntà \\
throw-\gl{itive} \\
\glt `throw' [German: `\textit{werfen}'], cf. Margi \textit{ndalna} `throw (away)' [German: `(\textit{weg})\textit{werfen}']
\z 

\citet[136]{hoffmannZurSpracheCibak1955} discusses several other forms with directional meaning but is hesitant to analyze them as suffixes. The most relevant example is the verb \textit{lə̀kwá} `enter' which is treated as a contraction of \textit{li} `go' and the preposition \textit{akwa} `into' (cf. Bura example (\ref{ex2:bura3}) above). It appears that Hoffmann did not have enough data to further describe this aspect of the verbal morphology. 

\citet[140]{hoffmannZurSpracheCibak1955} briefly mentions a particular combination of verbs involving a directional verb and a motion verb, shown in (\ref{ex2:ciba5}). The example resembles a directional multiverb construction, but no further information is given.

\ea\label{ex2:ciba5}
\langinfo{Cibak directional multiverb construction?}{}{\cite[140]{hoffmannZurSpracheCibak1955}} \\
\gll in yi kə̀ hùy ā́ tə́rà \\
?? 1\gl{sg} ?? run \gl{prep} go.away \\
\glt `I ran away.'
\z 

\subparagraph{Caused accompanied motion (CAM)}

Expressions of (directed) CAM are not sufficiently covered by \citet{hoffmannZurSpracheCibak1955} to make any generalizations about how CAM is expressed in Cibak.

\subparagraph{Buy-sell verbs}

Hoffmann's few examples of the Cibak outward extension show that a verb root \textit{dəl} which means `sell' on its own, as in (\ref{ex2:ciba3}), is interpreted as `buy' with an outward suffix, as in (\ref{ex2:ciba4}). This contrasts with the closely related Bura language where the verb root \textit{dila} `sell' can take directional extensions without this semantic change, and instead a different verb root is used to express `buy'.

\ea\label{ex2:ciba3}
\langinfo{Cibak verb \textit{də̀l} ‘sell’ with TAM suffix}{}{\cite[133]{hoffmannZurSpracheCibak1955}} \\
\gll də̀l-tì \\
sell-?? \\
\glt `sell'

\ex\label{ex2:ciba4}
\langinfo{Cibak verb \textit{də̀l} ‘sell’ with outward suffix}{}{\cite[133]{hoffmannZurSpracheCibak1955}} \\
\gll də́l-bà \\
sell-\gl{out} \\
\glt `buy'
\z 

\paragraph*{Margi Central (marg1265)} \label{marg}

Information on Margi grammar is from a reference grammar \citep{hoffmannGrammarMargiLanguage1963}. The description of Margi does not provide morpheme breaks or interlinearized glossing, so it is not always clear whether a particular verb ending is a directional suffix or not. For this reason, the discussion of Margi data is primarily limited to sections of the grammar where Hoffmann directly discusses relevant parts of the grammar. I add segmentation and glosses (as possible) to the examples below, using the directional concepts Hoffmann refers to as labels for the suffixes \citep[121--149]{hoffmannGrammarMargiLanguage1963}. 

\subparagraph{Directional verbs}~

\begin{table}[H]
	\caption{Margi directional verbs} 
	\label{tab:margV}
	\begin{tabular}{lll} % add l for every additional column or remove as necessary
		\lsptoprule
		DIR & verb & gloss   \\ %table header
		\midrule
		ITIVE & \textit{mái} & `go (away)' \\
		VENT & \textit{shìlí} & `come' \\
		ITIVE.UP & \textit{dù} & `go up' \\
		VENT.UP & \textit{sə̀dù} & `come up' \\
		DOWN & \textit{dzìa} & `go down, come down' \\
		ITIVE.INTO & \textit{lì} & `go in, go (home)' \\
		VENT.INTO & \textit{sì} & `come in, come (home)' \\
		ITIVE.OUT & \textit{bà} & `go out, go across' \\
		VENT.OUT & \textit{zə̀bà} & `come out' \\ 
		\lspbottomrule
	\end{tabular}
\end{table}

\citet[124, 149]{hoffmannGrammarMargiLanguage1963} notes the shared forms in \textit{dù} `go up' and \textit{sə̀dù} `come up', in \textit{bà} `go out, go across'  and \textit{zə̀bà} `come out', as well as the shared vowel of \textit{lì} `go in, go (home)' and \textit{sì} `come in, come (home)', but treats this similarity as a question of etymology, not as part of the synchronic morphological system of Margi. 

\newpage
\subparagraph{Directional extensions}~

\begin{table}[H]
	\caption{Margi directional morphology}
	\label{tab:marg}
	\begin{tabular}{llll} % add l for every additional column or remove as necessary
		\lsptoprule
		forms &  source gloss & direction  & other functions  \\ %table header
		\midrule
		\textit{-nà}  &  & \textsc{itive}  &  lexicalized  \\
		\textit{-ía}  &   &   \textsc{down}  &    \\
		\textit{-wá}  &    & \textsc{into}   &  lexicalized  \\		
		\textit{-bá}  &   &  \textsc{out}  &  \\
		\textit{-ŋgéri}  &    &  \textsc{onto}  &  \textsc{add}, lexicalized  \\
		%\textit{-ari}  &   & lexicalized?   &   lexicalized? \\
		\lspbottomrule
	\end{tabular}
\end{table} 

\citet{hoffmannGrammarMargiLanguage1963} describes six motion-related suffixes in Margi, but the directional functions of these suffixes are lexically restricted to a relatively small class of verbs. With most verbs, the suffixes have a lexicalized, non-motion meaning. It is also worth noting that no forms are described as ventive. 

The suffixes that show evidence of having directional meaning with a manner-of-motion verb are the itive (with an obligatory de-transitivizing morpheme), outward and downward suffixes. All the directional suffixes are shown to be able to have a directional function when the patient of a transitive verb is interpreted as a figure on a path of motion.

The itive suffix \textit{-nà} can have a directional interpretation with a manner-of-motion verb, but only when the de-transitivizing morpheme \textit{kə́r} (elsewhere meaning `head') is present, as in (\ref{ex2:marg1}). Under the same restriction, the itive suffix can also appear with the intrinsically itive directional verb \textit{mái} `go (away)', as in (\ref{ex2:marg2}). This is the only example found of a directional suffix with an intransitive directional motion verb.

\ea\label{ex2:marg1}
\langinfo{Margi verb \textit{vəl} `fly' with itive suffix}{}{\cite[153]{hoffmannGrammarMargiLanguage1963}} \\
\gll  vəl-ná *(kə́r) \\
fly-\gl{itive} head	\\
\glt `to fly away'

\ex\label{ex2:marg2}
\langinfo{Margi verb \textit{mái} `go' with itive suffix}{}{\cite[153]{hoffmannGrammarMargiLanguage1963}} \\
\gll mái-ná *(kə́r)	\\
go-\gl{itive} head		\\
\glt `to go away'
\z

The itive suffix can also express itive direction when referring to the path of motion of the patient of a transitive verb, as in (\ref{ex2:marg3}).

\ea\label{ex2:marg3}
\langinfo{Margi verb \textit{ndàl} `throw' with itive suffix}{}{\cite[131]{hoffmannGrammarMargiLanguage1963}} \\
\gll ndàl-nà	\\
throw-\gl{itive}		\\
\glt `to throw away'
\z 

The other place where the suffix \textit{-ná} has a directional meaning is in a caused accompanied motion context, discussed below. Otherwise, the suffix frequently appears with no discernible motion-related meaning, as in (\ref{ex2:marg4}) and (\ref{ex2:marg5}).

\ea\label{ex2:marg4}
\langinfo{Margi verb \textit{dlà} `fall' with itive suffix}{}{\cite[130]{hoffmannGrammarMargiLanguage1963}} \\
\gll dlà-nà	\\
fall-\gl{itive}		\\
\glt `to overthrow'

\ex\label{ex2:marg5}
\langinfo{Margi verb \textit{ghə̀ɗá} `get tired' with itive suffix}{}{\cite[130]{hoffmannGrammarMargiLanguage1963}} \\
\gll ghə̀ɗá-ná	\\
get.tired-\gl{itive}		\\
\glt `to make tired; to tire, to fatigue'
\z

Like the itive suffix, the outward suffix \textit{-bá} can have a directional function with manner-of-motion verbs (but without needing the de-transitivizer), as in (\ref{ex2:marg6}), with a transitive verb whose patient is the figure on a path of motion, as in (\ref{ex2:marg7}), or in a CAM context (discussed below).\footnote{\citet[125]{hoffmannGrammarMargiLanguage1963} notes that the outward suffix is the same form as the outward directional verb \textit{bá} and likely etymologically related.} 

\ea\label{ex2:marg6}
\langinfo{Margi verb \textit{mə̀̀} `run' with outward suffix}{}{\cite[122]{hoffmannGrammarMargiLanguage1963}} \\
\gll mə̀-bá	\\
run-\gl{out}\\
\glt `to start up and run out'

\ex\label{ex2:marg7}
\langinfo{Margi verb \textit{ndàl} `throw' with outward suffix}{}{\cite[123]{hoffmannGrammarMargiLanguage1963}} \\
\gll ndàl-bá	\\
throw-\gl{out}	\\
\glt `to throw out'
\z

However, with most verbs, the outward suffix does not add any directional meaning, as in (\ref{ex2:marg8}).

\ea\label{ex2:marg8}
\langinfo{Margi verb \textit{ŋá} `hear' with outward suffix}{}{\cite[123]{hoffmannGrammarMargiLanguage1963}} \\
\gll ŋá-bá	\\
hear-\gl{out} \\
\glt `to understand; to hear'
\z

Like the outward suffix, the downward suffix \textit{-ía} can have a directional interpretation with a manner-of-motion verb, as in (\ref{ex2:marg9}), or when the patient of a transitive verb is the figure on a path of motion, as in (\ref{ex2:marg10}), or in a CAM context (discussed below). 

\ea\label{ex2:marg9}
\langinfo{Margi verb \textit{və̀l} `jump' with downward suffix}{}{\cite[128]{hoffmannGrammarMargiLanguage1963}} \\
\gll  və̀l-ía	\\
jump-\gl{down}		\\
\glt `to jump down'

\ex\label{ex2:marg10}
\langinfo{Margi verb \textit{p} `put' with downward suffix}{}{\cite[127]{hoffmannGrammarMargiLanguage1963}} \\
\gll p-ía	\\
put-\gl{down}	\\
\glt `to put down (many)'
\z

Again, with most verbs, the downward suffix does not have a motion-related meaning, as in (\ref{ex2:marg11}). 

\ea\label{ex2:marg11}
\langinfo{Margi verb \textit{ts} `kill' with downward suffix}{}{\cite[127]{hoffmannGrammarMargiLanguage1963}} \\
\gll ts-ía	\\
kill-\gl{down} \\
\glt `to kill, slaughter'
\z 

The suffixes \textit{-wá} `into' and \textit{-ŋgéri} `onto' are only found with directional meaning in examples where the patient of a transitive verb is the figure on the path of motion, as in (\ref{ex2:marg12}) and (\ref{ex2:marg13}).  In other contexts they do not appear to have any motion related meaning.

\ea\label{ex2:marg12}
\langinfo{Margi verb \textit{f} `put' with inward suffix}{}{\cite[148]{hoffmannGrammarMargiLanguage1963}} \\
\gll  f-wá	\\
put-\gl{into}		\\
\glt `to put (one) into'

\ex\label{ex2:marg13}
\langinfo{Margi verb \textit{ndàl} `throw' with `onto' suffix}{}{\cite[135]{hoffmannGrammarMargiLanguage1963}} \\
\gll  ndàl-ŋgə́ri	\\
throw/twist-\gl{onto} \\
\glt `to throw on top, over; to twist (p.'s neck)'
\z

The inward suffix can also indicate that the object is split into parts as a result of the action of the predicate, as in (\ref{ex2:marg02}). 

\ea\label{ex2:marg02}
\langinfo{Margi verb \textit{ɓàts} `break' with inward suffix}{}{\cite[148]{hoffmannGrammarMargiLanguage1963}} \\
\gll ŋàhyú-wá \\
trample-\gl{into}	\\
\glt `to trample on something and divide into parts'
\z 

%The upward suffix \textit{-ari} only appears with motion meaning in the context of CAM, as discussed below. Even with a manner-of-motion verb, as in (\ref{ex2:marg14}), or a transitive motion verb, as in (\ref{ex2:marg15}), the suffix \textit{-ari} does not appear to convey upward directional meaning.

%\ea\label{ex2:marg14}
%\langinfo{Margi verb \textit{ɓ} `walk' with upward suffix \textit{-arí}}{}{\cite[120]{hoffmannGrammarMargiLanguage1963}} \\
%\gll   ɓ-àr ɓù	\\
%walk-\gl{up} walk	\\
%\glt `to walk a bit'
%
%\ex\label{ex2:marg15}
%\langinfo{Margi verb \textit{ŋə̀rz} `push' with upward suffix \textit{-arí}}{}{\cite[120]{hoffmannGrammarMargiLanguage1963}} \\
%\gll   ŋə̀rz-àrì	\\
%push-\gl{up}	\\
%\glt `to push along'
%\z

\subparagraph{Caused accompanied motion (CAM)}

Caused accompanied motion (CAM) can be expressed by the combination of a motion suffix with the verb \textit{hù} take' (\textit{fá} in the plural form). This is attested with the itive suffix (\ref{ex2:marg16}), outward suffix (\ref{ex2:marg17}) and downward suffix (\ref{ex2:marg18}). 

\ea\label{ex2:marg16}
\langinfo{Margi verb \textit{hù} `take' with itive suffix}{}{\cite[114]{hoffmannGrammarMargiLanguage1963}} \\
\gll  hɘ̀-nà \\
take-\gl{itive} \\
\glt `to take away'

\ex\label{ex2:marg17}
\langinfo{Margi verb \textit{hù} `take' with outward suffix}{}{\cite[114]{hoffmannGrammarMargiLanguage1963}} \\
\gll hə́-bá \\
take-\gl{out} \\
\glt `to take out'

\ex\label{ex2:marg18}
\langinfo{Margi verb \textit{hù} `take' with downward suffix}{}{\cite[114]{hoffmannGrammarMargiLanguage1963}} \\
\gll hy-ía	\\
take-\gl{down}	\\
\glt `to take down'
%
%\ex\label{ex2:marg19}
%\langinfo{Margi verb \textit{hù} `take' with upward suffix}{}{\cite[114]%{hoffmannGrammarMargiLanguage1963}} \\
%\gll h-àrì	\\
%take-\gl{up}	\\
%\glt `to lift up (a bit)'
\z

There is one example of a manner-of-motion verb `swim' with an outward suffix, shown in (\ref{ex2:marg20}), in which the translation suggests a caused accompanied motion function. 

\ea\label{ex2:marg20}
\langinfo{Margi verb \textit{ŋkə̀} `swim' with outward suffix}{}{\cite[123]{hoffmannGrammarMargiLanguage1963}} \\
\gll ŋkə̀-bá \\
swim-\gl{out} \\
\glt `to take p. across the river by swimming (also: to swim across ?)'
\z

Caused accompanied motion can also be expressed by combining the same verb \textit{hù} `take' with a motion verb in a form that Hoffmann calls ``adverbial present'' indicated by a prefix \textit{a-}, as in (\ref{ex2:marg21}).

\ea\label{ex2:marg21}
\langinfo{Margi verb \textit{hù} `take' multiverb construction}{}{\cite[267]{hoffmannGrammarMargiLanguage1963}} \\
\gll ỳo! mà mə̀mù̥ kam, hə̀r-ɗà á-shìlí  \\
?? if honey indeed(Hausa) take-\gl{obj}.1\gl{sg} \textsc{adv}.\gl{prs}-come \\
\glt `All right! if (it is) honey, bring me (it)!'
\z 

\subparagraph{Buy-sell verbs}

The verb \textit{ɗə̀l} can mean `buy' or `sell', but with an itive directional suffix, \textit{ɗə̀lnà}, it means `sell' \citep[122]{hoffmannGrammarMargiLanguage1963}.

\paragraph*{Huba (Kilba) (huba1236)} \label{huba}

Information on Huba grammar is from a grammar sketch \citep{muazuGrammarKilbaLanguage2009} and a lexicon \citep{muazuModernKilbaDictionary2015}.

\subparagraph{Directional verbs}~

\begin{table}[H]
	\caption{Huba directional verbs} 
	\label{tab:hubaV}
	\begin{tabular}{lll} % add l for every additional column or remove as necessary
		\lsptoprule
		DIR & verb & gloss   \\ %table header
		\midrule
		ITIVE & \textit{mai} & `to go' \\
		VENT & \textit{àndà} & `come' \\
		INTO& \textit{gwà(kwà)} & `enter' \\
		UP & \textit{igí} & `to go up' \\
		\lspbottomrule
	\end{tabular}
\end{table}

\citet[40--41]{muazuModernKilbaDictionary2015} give three separate entries for a verb glossed `to go': \textit{ma'i}, \textit{máí} and \textit{mai}. The English-Kilba section of the dictionary lists \textit{àndà} as `come', but this word is not found in the Kilba-English section. The closest word is \textit{àndá} `there, yonder'. \citet[24]{muazuGrammarKilbaLanguage2009} lists another word, \textit{siya} `come down' which would appear to be a downward suffix attached to a root verb \textit{si} `come' (cf. Bura and Cibak), but this analysis cannot be confirmed. A very similar word, \textit{síyí} `to go down' appears in \citet[61]{muazuModernKilbaDictionary2015}. The same gloss is given to \textit{màìse} `go down' \citep[81]{muazuModernKilbaDictionary2015}. 

In addition to \textit{gwà} `enter', the dictionary also lists \textit{gwákwà} `to enter, to go in or among' and \textit{gwàkwà} `to go into, to enter' \citep[25, 80]{muazuModernKilbaDictionary2015}. As in other Bura-Marghi languages, there is some doubt as to whether this should be analyzed as a verbal suffix or as a verb plus the preposition \textit{akwa} `into'. 

\newpage
\subparagraph{Directional extensions}~

\begin{table}[H]
	\caption{Huba directional morphology}
	\label{tab:huba}
	\begin{tabular}{llll} % add l for every additional column or remove as necessary
		\lsptoprule
		forms & source gloss & direction  & other functions  \\ %table header
		\midrule
		\textit{-nà}  &  \textsc{itive}  &  \textsc{itive} &     \\
		%\textit{-wa, -kwa, -u}  &  \textsc{into}  & \textsc{into}  &    \\
		\textit{-yà}  &  \textsc{down}  &   \textsc{down}  &  \textsc{compl}   \\
		\textit{-bìyà}  &  \textsc{out}  &  \textsc{out}   & \textsc{compl} \\
		\textit{-ngərì}  &  \textsc{onto}  &   \textsc{onto}   &  \textsc{loc, add}, lexicalized \\		
		%\textit{-vi}  &  \textsc{home}  &   \textsc{home}  &   \\	
		\lspbottomrule
	\end{tabular}
\end{table} 

The itive suffix \textit{-nà}, the outward suffix \textit{-bìyà} and the \textsc{onto} suffix \textit{-ngərì} are shown with directional meanings with transitive motion verbs, as in (\ref{ex2:huba1}) and (\ref{ex2:huba2}). No examples are given of these suffixes with intransitive motion verbs. With other verbs, the meanings of the suffixes are not transparent, although in several cases the outward suffix has a meaning similar to `completely' (as in \ref{ex2:huba2}) and the \textsc{onto} suffix can mean `additional'. 

\ea\label{ex2:huba1}
\langinfo{Huba verb \textit{pù} `pour' with itive suffix}{}{\cite[109]{muazuGrammarKilbaLanguage2009}} \\
\gll pù-nà \\
pour-\gl{itive} \\
\glt `to pour away'

\ex\label{ex2:huba2}
\langinfo{Huba verb \textit{vàk} `throw' with outward suffix}{}{\cite[113]{muazuGrammarKilbaLanguage2009}} \\
\gll vàk-bìyà	 \\
throw-\gl{out} \\
\glt `you throw out/all'
\z 

The one directional suffix shown with an intransitive motion verb is the downward suffix in (\ref{ex2:huba3}). The downward suffix can also have a directional meaning with transitive motion verbs, as in (\ref{ex2:huba4}). 

\ea\label{ex2:huba3}
\langinfo{Huba verb \textit{fəl} `jump' with downward suffix}{}{\cite[108]{muazuGrammarKilbaLanguage2009}} \\
\gll fəl-yà \\
jump-\gl{down} \\
\glt `to jump down'

\ex\label{ex2:huba4}
\langinfo{Huba verb \textit{pù} `put' with downward suffix}{}{\cite[108]{muazuGrammarKilbaLanguage2009}} \\
\gll pù-yà \\
put-\gl{down} \\
\glt `to put down (many)'
\z 

In addition to the four directional suffixes described by \citet{muazuGrammarKilbaLanguage2009}, there is other verbal morphology with similarities to directional extensions described in Bura-Marghi languages. The example of \textit{gwákwà} `to enter, to go in or among' is mentioned above as possible evidence of an inward suffix (although the verb \textit{gwà} already means `to enter'). \citet[24]{muazuModernKilbaDictionary2015} give the form \textit{gúví} `to go into a house or room' which seems to show a parallel with the homeward suffixes found in other Bura-Marghi languages, but in Kilba it is not analyzed as a suffix.

\subparagraph{Caused accompanied motion (CAM)}

The verb \textit{shìlnì} `bring' expresses (ventive) caused accompanied motion \citep[79]{muazuModernKilbaDictionary2015}. It also appears that CAM may be expressed by combining a directional suffix with a verb glossed `take' such as \textit{hù-ya} `take down' and \textit{hù-nà} `take away'  \citep[108--109]{muazuGrammarKilbaLanguage2009}. These examples are not given in the context of a full sentence so it is somewhat uncertain that they necessarily express a translational motion event. 

\subparagraph{Buy-sell verbs}

A verb transcribed as \textit{ɗəlá} or \textit{dəlà} or \textit{dùlà} means `buy' \citep[19, 79]{muazuModernKilbaDictionary2015}. This verb root is found with an outward suffix, \textit{ɗələ-bíyà}, in which case it still means `buy' but possibly with an additional meaning of `completely'. With the itive suffix, \textit{dúl-nà}, the meaning is `sell' \citep[17]{muazuModernKilbaDictionary2015}.

\subsubsection{Mandaraic (6/9)}

Six of nine Mandaric languages are discussed here. Short grammar sketches of Cineni (cine1238) and Gvoko (gvok1239) do not address directional extensions \citep{kimChineneCineneBugdong2010,kimBugdongNaijiliaChadigeodeului2013}. There is no available description of the grammar of Guduf-Gava (gudu1252).

\paragraph*{Dghwede (dghw1239)} \label{dghw}

Most information on Dghwede grammar is from a journal article \citep{frickVerbalSystemDghwede1978} which describes many motion-related verbal suffixes in Dghwede. However, the analysis contains a few unsubstantiated claims, and there are no morpheme breaks or glosses included, nor is there a lexicon of the language available. This makes it difficult, and, in some cases, impossible, to determine the lexical semantics of the verb roots appearing in each example sentence. Where possible, glosses are added based on \citet{hartmannInformationStructureLeft2006}.

\subparagraph{Directional verbs}~

\begin{table}[H]
	\caption{Dghwede directional verbs} 
	\label{tab:dghwV}
	\begin{tabular}{lll} % add l for every additional column or remove as necessary
		\lsptoprule
		DIR & verb & gloss   \\ %table header
		\midrule
		ITIVE & \textit{jə̀wá, dá} & `go' \\
		VENT & \textit{sá} & `come' \\
		INTO& \textit{tá} & `enter' \\
		\lspbottomrule
	\end{tabular}
\end{table}

The verb roots \textit{dá} `go', \textit{sá} `come' and \textit{tá} `descend' are said to be ``defective'' in that they do not occur without a suffix \citep[12]{frickVerbalSystemDghwede1978}. The verb \textit{tá} `descend' is transcribed \textit{txá} by \citet[27]{frickPhonologyDghwede1977}.

\subparagraph{Directional extensions}~

\begin{table}[H]
	\caption{Dghwede directional morphology}
	\label{tab:dghw}
	\begin{tabular}{llll} % add l for every additional column or remove as necessary
		\lsptoprule
		forms & source gloss & direction  & other functions  \\ %table header
		\midrule
		\textit{-(d)àwá}  &  \textsc{vent.across}  &   \textsc{vent}  &   \textsc{subs.vent}  \\		
		\textit{-gré}  &  \textsc{dir}  &   \textsc{itive} &   \textsc{subs.vent}?  \\
		\textit{-dá}  &  \textsc{itive.across}  &  \textsc{itive}  &    \\
		\textit{-(de)gwá}  &  \textsc{up}  & \textsc{up}  &  \textsc{conc.up}, \textsc{subs.up} \\
		\textit{-(de)fè} & \textsc{up2} & \textsc{up}  &  \\
		\textit{-tígè}  &  \textsc{itive.down}  &   \textsc{dir.itive.down} &   \\
		\textit{-àwáyà, -dáyà,}  &  \textsc{vent.down}  &  \textsc{dir.(vent.)down}  & \textsc{subs.vent.down}   \\
		\textit{-àwí, -dí} & & & \\
		\textit{-mè}  &  \textsc{into}  &   \textsc{into}  &    \\		
		\textit{-àwílè, -dílè}  &  \textsc{out}  &  \textsc{(vent.)out}  &  \textsc{subs.vent.out}, lex. \\	
		\textit{-ségè}  &  \textsc{under}  &  \textsc{under}   &   \textsc{loc} \\
		\textit{-gá}  &  \textsc{onto}  &  \textsc{onto}   &  \\
		\lspbottomrule
	\end{tabular}
\end{table} 

\citet{frickVerbalSystemDghwede1978} describes many directional extensions in Dghwede. Of these, several have two forms, one of which has an additional initial consonant \textit{d} and some other morphophonological alternation, as shown in Table \ref{tab:dghw}. Possible analyses of these forms will be discussed at the end of this section. 

The suffix \textit{-dá} is described as an itive directional used for paths without a significant change in the vertical dimension.\footnote{There appears to be an error in one place where \textit{-dá} is labeled ``across to [the speaker]'' \citep[20]{frickVerbalSystemDghwede1978}. Elsewhere the description is ``away from the speaker'' \citep[5, 18, 19]{frickVerbalSystemDghwede1978}.} Translations for examples with \textit{-dá} are often not explicitly itive (e.g., ``away''), but are compatible with itive interpretations, as in (\ref{ex2:dghw10}). The suffix is similar in form to the preposition \textit{dà} `to(ward)'. 

\ea\label{ex2:dghw10}
\langinfo{Dghwede verb \textit{wə̀rə̀} `return' with itive suffix}{}{\cite[20]{frickVerbalSystemDghwede1978}} \\
\gll kâ' wə̀rə̀-dá dà lùwá	\\
\gl{narr} return-\gl{itive} to home \\
\glt `He went across back home.'
\z 

Frick presents the suffix \textit{-àwá} as the ventive counterpart of \textit{-dá}. The example in (\ref{ex2:dghw11}) seems to show an explicit ventive interpretation of \textit{-àwá} though the meaning of the verb root has not been verified. 

\ea\label{ex2:dghw11}
\langinfo{Dghwede verb \textit{mbə́r} `hurry?' with ventive suffix}{}{\cite[22]{frickVerbalSystemDghwede1978}} \\
\gll ʽà mbə́r-èn-áwá gè \\
\gl{narr} ??-\gl{tr}-\gl{vent} ?? \\
\glt `I hurried across here.'
\z 

Frick describes the upward suffix \textit{-gwá} as occurring in both itive and ventive contexts. The sentence in (\ref{ex2:dghw3}) is used twice, once with a ventive translation, and once with a translation meant to illustrate itive motion. The simplest analysis of this data is that this upward suffix \textit{-gwá} expresses neither itive nor ventive directional meaning. 

\ea\label{ex2:dghw3}
\langinfo{Dghwede verb \textit{xwáyə̀} `run' with upward suffix}{}{\cite[20, 23]{frickVerbalSystemDghwede1978}} \\
\gll ká' gà lə́kúve cì xwáyə̀-gwá	\\
\gl{narr} ?? baboon ?? run-\gl{up}	\\
\glt `The baboons ran upwards.'  \citep[20]{frickVerbalSystemDghwede1978} \\
`The baboons came running up.'	\citep[23]{frickVerbalSystemDghwede1978}	 	
\z 

There is one example of \textit{-gwá} with a non-motion verb where the translation indicates an upward concurrent motion meaning, shown in (\ref{ex2:dghw5}). Other examples of this suffix with non-motion verbs do not give any clear motion-related interpretation. 

\ea\label{ex2:dghw5}
\langinfo{Dghwede verb \textit{táwə̀} `weep' with upward suffix}{}{\cite[6]{frickVerbalSystemDghwede1978}} \\
\gll kâ' táwə̀-gwá táwà dà lùwá \\
\gl{narr} weep-\gl{up} \gl{cont} to home \\
\glt `He was crying all the way up home.'
\z 

Frick describes two downward motion directional suffixes: the itive downward suffix, \textit{-tígè}, and the ventive downward suffix, \textit{-àwáyà} or \textit{-àwí}. However, the translations given in the relevant examples do not always make the itive-ventive distinction explicit, as seen in (\ref{ex2:dghw6}) and (\ref{ex2:dghw7}). 

\ea\label{ex2:dghw6}
\langinfo{Dghwede verb \textit{bə̀l'} `escape' with (itive?) downward suffix}{}{\cite[20]{frickVerbalSystemDghwede1978}} \\
\gll  ski cé gè bə́l'ə́-tíge dà mbé \\
thing \gl{pst} ?? escape-\gl{itive}.\gl{down} to ?? \\
\glt `My thing has escaped down into it.'

\ex\label{ex2:dghw7}
\langinfo{Dghwede verb \textit{bə̀l'} `escape' with (ventive?) downward suffix}{}{\cite[22]{frickVerbalSystemDghwede1978}} \\
\gll  ká' bə̀l'-àwí dà xáyà \\
\gl{narr} escape-\gl{vent}.\gl{down} to ground	\\
\glt `He escaped down to the ground.'
\z

There is a pair of suffixes that express inward-outward directional meaning. \citet[19]{frickVerbalSystemDghwede1978} expresses some uncertainty about the inward suffix, \textit{-mè}.\footnote{ ``...\textit{-mè}, -\textit{gá}, \textit{-sə́gè}, \textit{-défè} and \textit{-fè} do have some directional meaning, although it is not as clearly definable...'' \citep[19]{frickVerbalSystemDghwede1978}.} This may be because it is only shown with one motion verb, seen in (\ref{ex2:dghw8}), and its function with non-motion verbs is not transparent. 

\ea\label{ex2:dghw8}
\langinfo{Dghwede verb \textit{ə̀yə̂} `fall' with inward suffix}{}{\cite[25]{frickVerbalSystemDghwede1978}} \\
\gll ə̀yə̂-me cé ne dá me kárà... \\
fall-\gl{into} \gl{pst} ?? to ?? fire \\
\glt `When he had fallen into the fire.'
\z

The suffix \textit{-àwílè} expresses outward directional meaning, but is also translated with ventive direction, as in (\ref{ex2:dghw9}).

\ea\label{ex2:dghw9}
\langinfo{Dghwede verb \textit{xwáy} `run' with (ventive?) outward suffix}{}{\cite[23]{frickVerbalSystemDghwede1978}} \\
\gll kâ' xwáy-àwilè	\\
\gl{narr} run-\gl{vent}.\gl{out} \\
\glt `He came running out.'
\z 

The suffixes \textit{-fè} `up', \textit{-sege} `under' and \textit{-ga} `upon' may function as directional markers, at least with transitive motion verbs, as shown in (\ref{ex2:dghw14}), (\ref{ex2:dghw12}) and (\ref{ex2:dghw13}). 

\ea\label{ex2:dghw14}
\langinfo{Dghwede verb \textit{ndə́ɗə́} `pull?' with upward suffix}{}{\cite[25]{frickVerbalSystemDghwede1978}} \\
\gll kâ' ndə́ɗə́-nə́-fe ɬágə́là	\\
\gl{narr} ??-\gl{tr}-\gl{up}2 grass \\
\glt `He pulled up grass.'

\ex\label{ex2:dghw12}
\langinfo{Dghwede verb \textit{kàt} `push?' with `under' suffix}{}{\cite[26]{frickVerbalSystemDghwede1978}} \\
\gll kàt-àrə̀-sə́ge kátá kárà	\\
??-\gl{dat}-\gl{under} ?? fire	 \\
\glt `Push the fire under it!'

\ex\label{ex2:dghw13}
\langinfo{Dghwede verb \textit{pàɗ} `pour' with `onto' suffix}{}{\cite[26]{frickVerbalSystemDghwede1978}} \\
\gll ká' tè pàɗ-àrə̀-gà páɗá yíwè nè kwíré yíwè \\
\gl{narr} ?? pour-\gl{dat}-\gl{onto} pour water ?? stone water	\\
\glt `They were repeatedly pouring water on the water stone.'
\z 

Finally, the suffix \textit{-gré} is described as meaning both ``away from the speaker'' and ``toward the speaker''. Note that in examples (\ref{ex2:dghw1}) and (\ref{ex2:dghw2}) the same suffix is used with the same verb, \textit{xwáyə̀} `run'. In (\ref{ex2:dghw1}) an itive interpretation is given, while in (\ref{ex2:dghw2}) a ventive interpretation is given. Since (\ref{ex2:dghw2}) is the only ventive use of \textit{-gré}, it may be mistranslated. 

\ea\label{ex2:dghw1}
\langinfo{Dghwede verb \textit{xwáyə̀} `run' with \textit{-gré}}{}{\cite[20]{frickVerbalSystemDghwede1978}} \\
\gll 	kâ' də́gwe cì xwáyə̀-gŕe dà lùwá	\\
\gl{narr} ?? ?? run-\textsc{gre} to home	\\
\glt `The girl ran away home(wards).'

\ex\label{ex2:dghw2}
\langinfo{Dghwede verb \textit{xwáyə̀} `run' with \textit{-gré}}{}{\cite[21]{frickVerbalSystemDghwede1978}} \\
\gll nîne xwáyə̀-gr̀e m̀dá dà màná \\
?? run-\textsc{gre} ?? to ?? 	\\
\glt `This is why we (\textsc{excl}) have run toward here.'
\z 

%\begin{table}[H]
%	\caption{Dghwede transitive (ventive) directional morphology}
%	\label{tab:dghwD}
%	\begin{tabular}{llll} % add l for every additional column or remove as necessary
%		\lsptoprule
%		forms & gloss & \textit{d}-form  & \textit{d-}form function  \\ %table header
%		\midrule
%		\textit{-gré}  &  \textsc{dir}  &  \textit{-də̀-gré} &  \textsc{vent.(subs.)cam}, ?? \\
%		\textit{-dá}  &  \textsc{itive.across}  &  NA  &    \\
%		\textit{-àwá}  &  \textsc{vent.across}  &  \textit{-d-àwá}  &   \textsc{vent.(subs.)cam}  \\		
%		\textit{-gwá}  &  \textsc{up}  & \textit{-də̀-gwá}  & \textsc{vent.up.(subs.)cam}, ?? \\
%		\textit{-fè}  & \textsc{up2} & \textit{-də́-fè} & \textsc{dir.up} ?  \\
%		\textit{-tígè}  &  \textsc{itive.down}  &   NA &   \\
%		\textit{-àwáyà, -àwí}  &  \textsc{vent.down}  &   \textit{d-áyà, d-í}  &    \textsc{vent.down.(subs.)cam},  \textsc{dir.vent.down} (trans.)  \\
%		\textit{-àwílè}  &  \textsc{out}  & \textit{d-ílè}  &  \textsc{vent.out.(subs.)cam},  \textsc{dir.vent.out} (trans.) \\
%		\textit{-mè}  &  \textsc{into}  &   NA  &    \\		
%		\textit{-ségè}  &  \textsc{under}  & NA  &  \\
%		\textit{-gá}  &  \textsc{onto}  & NA  &  \\
%		\lspbottomrule
%	\end{tabular}
%\end{table} 

As shown in Table \ref{tab:dghw}, some of the directional extensions have a form which begins with an initial consonant \textit{d}.\footnote{\citet[19]{frickVerbalSystemDghwede1978} also discusses a causative suffix, \textit{-(d)úwè} in this group. It is not considered here because it does not have motion-related meaning.} As \citet[19]{frickVerbalSystemDghwede1978} notes, the forms beginning with \textit{d} are only used with transitive predicates, and correspond to translations with ventive motion (except perhaps in the case of \textit{-də́-fè} where the translations of the available examples do not make an explicit itive-ventive distinction). The ventive meaning of the \textit{d}-forms is clear when we compare the ventive (subsequent motion) interpretation of the suffix \textit{-də̀gŕe} in (\ref{ex2:dghw18}) with the itive interpretations of \textit{-gré} in most examples. 

\ea\label{ex2:dghw18}
\langinfo{Dghwede verb \textit{pàɗə̀} `draw (water)' with transitive ventive suffix}{}{\cite[21]{frickVerbalSystemDghwede1978}} \\
\gll  ká' tè pàɗə̀-də̀gŕe yíwè \\
\gl{narr} ?? draw-\gl{tr}.\gl{vent} water \\
\glt `They drew water (and brought it here).'
\z 

Likewise, compare the ventive interpretation of the suffix \textit{-də̀gwá} in (\ref{ex2:dghw19}) with the (non-ventive) interpretation of \textit{-gwá} in (\ref{ex2:dghw3}).

\ea\label{ex2:dghw19}
\langinfo{Dghwede verb \textit{yáxàgə̀} `call' with transitive ventive upward suffix}{}{\cite[22]{frickVerbalSystemDghwede1978}} \\
\gll  kâ' Kálákwà yáxàgə̀-də̀gwá \\
\gl{narr} Kalakwa call-\gl{tr}.\gl{vent}.\gl{up} \\
\glt `Kalakwa called me (to come) up.'
\z 

It appears to be the case that the \textit{-d} morpheme only (but not always) appears when the object (or patient) is a figure on the path of motion. The following examples show various interpretations of the suffix \textit{-dáyà} (transitive ventive downward) according to the semantics of the verb it attaches to. In (\ref{ex2:dghw15}), the verb is a transitive (caused) motion verb, and it is only the object that is moving. In (\ref{ex2:dghw16}), the meaning of the verb dictates that both the subject and object are moving (caused accompanied motion). In (\ref{ex2:dghw17}), the translation indicates that the motion occurs after the event predicated by the verb (subsequent motion) and includes both the subject and object (caused accompanied motion). 

\ea\label{ex2:dghw15}
\langinfo{Dghwede verb \textit{wə̀jə̀} `throw' with transitive ventive downward suffix}{}{\cite[22]{frickVerbalSystemDghwede1978}} \\
\gll  kâ' plíse dà wə̀jə̀-dáyà \\ 
\gl{narr} horse ?? throw-\gl{vent}.\gl{down} \\
\glt `The horse will throw him off.'

\ex\label{ex2:dghw16}
\langinfo{Dghwede verb \textit{zə̀} `carry' with transitive ventive downward suffix}{}{\cite[5]{frickVerbalSystemDghwede1978}} \\
\gll ʽà zə̀-dáyà	\\
3\gl{sg} carry-\gl{vent}.\gl{down}	\\
\glt `He brought it down.'

\ex\label{ex2:dghw17}
\langinfo{Dghwede verb \textit{càcə̀} `chase' with transitive ventive downward suffix}{}{\cite[22]{frickVerbalSystemDghwede1978}} \\
\gll  ká' mdà càcə̀-dáyà \\
\gl{narr} ?? chase-\gl{vent}.\gl{down}	\\
\glt `We chased them (and brought them) down.'
\z 

When a \textit{d}-form is used with a non-motion verb, there is a subsequent (ventive) CAM interpretation, as in (\ref{ex2:dghw18}), (\ref{ex2:dghw17}) and (\ref{ex2:dghw21}). This contrasts with the use of the shorter directional suffixes (without \textit{d}) with non-motion verbs. In most cases the translations of these examples do not have any clear motion meaning, but in two examples the translation is of subsequent ventive motion of the subject/agent (not accompanied by the object/patient), as in (\ref{ex2:dghw23}). Note that (\ref{ex2:dghw23}) also includes a `partial' suffix \textit{-n}.\footnote{\citet[27]{frickVerbalSystemDghwede1978} refers to \textit{-n} in (\ref{ex2:dghw23}) as a transitivizing suffix resulting in a perplexing analysis of the data. Dghwede has both a transitivizing suffix \textit{-n} and a `partial' suffix of the same form \citep[14, 16]{frickVerbalSystemDghwede1978}. Whether this is a case of polysemy or homophony is not completely clear, but, in any case, the assumption that the \textit{-n} here has `partial' meaning seems to offer a much simpler explanation of the data.} Semantically, it makes sense that if the releasing event was not completed, the patient would not be free to accompany the agent's path of motion subsequent to the event. In contrast, in (\ref{ex2:dghw22}), there is no partial suffix and the transitive form of the directional suffix is used giving a subsequent ventive CAM interpretation. 

\ea\label{ex2:dghw23}
\langinfo{Dghwede verb \textit{və́də̀} `release?' with partial and ventive suffixes}{}{\cite[27]{frickVerbalSystemDghwede1978}} \\
\gll ʽà və́də̀-n-àwà gə̀rà cè nè me fŕsə́nà \\
\gl{narr} release?-\textsc{partial}-\gl{vent} friend ?? ?? ?? prison \\
\glt He released his friend from prison (and came) across.

\ex\label{ex2:dghw22}
\langinfo{Dghwede verb \textit{və́də̀} `release?' with transitive ventive suffix}{}{\cite[27]{frickVerbalSystemDghwede1978}} \\
\gll ʽà və̀də̀-dàwà gə̀rà cè nè me fŕsə́nà \\
\gl{narr} release?-\gl{tr}.\gl{vent} friend ?? ?? ?? prison \\
\glt He released his friend from prison (and brought him) across.
\z 

\subparagraph{Caused accompanied motion (CAM)}

There is a verb, \textit{náxa}, glossed `bring!' \citep[14]{frickPhonologyDghwede1977}, but not seen in any example sentences. More commonly, caused accompanied motion (CAM) is expressed with the verbs \textit{za} `carry' and \textit{xə̀l} `move'. Both verbs are found with  directional suffixes to express directed CAM, as in (\ref{ex2:dghw16}) and (\ref{ex2:dghw20}).\footnote{In most examples, a transitive \textit{d}-form of the directional suffix is used, but Frick also gives examples with the suffixes \textit{dá} `\textsc{itive.across}', \textit{-tígè} `\textsc{itive.down}', \textit{-gwá} `\textsc{up}' and \textit{-gré} `\textsc{dir}'. When attached to CAM verbs, \textit{-gré} is translated with an itive interpretation contrasting with the ventive interpretation of \textit{-də̀gré}. }

\ea\label{ex2:dghw20}
\langinfo{Dghwede verb \textit{xəl} `move' with transitive ventive outward suffix}{}{\cite[23]{frickVerbalSystemDghwede1978}} \\
\gll xə̀lə̀-mə̀-díle xə́lá kfì \\
move-2\gl{pl}.\gl{imp}-\gl{tr}.\gl{vent}.\gl{out} move food	\\
\glt `Bring (\textsc{pl}) the food out!'
\z

\subparagraph{Buy-sell verbs}

The verb \textit{skwá} is glossed `buy' and maintains this meaning with a transitive ventive suffix, interpreted as subsequent (caused accompanied) motion, as in (\ref{ex2:dghw21}). The same verb has a `sell' meaning with the (transitive) causative suffix, \textit{'à skwə̀-dúwè} `he sold it' \citep[43]{frickVerbalSystemDghwede1978}. There is also one example of another verb in a causative form meaning `sell', but it is unclear what the meaning of the verb root is in other contexts \citep[25]{frickVerbalSystemDghwede1978}. 

\ea\label{ex2:dghw21}
\langinfo{Dghwede verb \textit{skwə̀} `buy' with transitive ventive suffix}{}{\cite[21]{frickVerbalSystemDghwede1978}} \\
\gll skwə̀-də̀gr̀e mè kàskwà gè	\\
buy-\gl{vent} ?? ?? ??	\\
\glt `I bought it in the market (and brought it here).'
\z 

\paragraph*{Matal (mata1306)} \label{mata}

The grammar of Matal is described in a PhD thesis \citep{dougopheGrammarMatal2021}. There is also a BA thesis describing Matal grammar based on texts from a New Testament translation \citep{verdizadeSelectedTopicsGrammar2018}. Data below will primarily be taken from \citet{dougopheGrammarMatal2021} given that it based on first-hand observation of oral use of the language, but some differences found in the translated texts used by \citet{verdizadeSelectedTopicsGrammar2018} will be noted as well.

\subparagraph{Directional verbs}~

\begin{table}[H]
	\caption{Matal directional verbs} 
	\label{tab:mataV}
	\begin{tabular}{lll} % add l for every additional column or remove as necessary
		\lsptoprule
		DIR & verb & gloss   \\ %table header
		\midrule
		ITIVE & \textit{d-} & `go' \\
		VENT & \textit{s-} & `come' \\
		\lspbottomrule
	\end{tabular}
\end{table}

\citet{dougopheGrammarMatal2021} also lists a verb \textit{mapiyay} glossed `to continue, to move on, to enter' but confirmed via personal communication that this is not likely a translocative motion verb. 

\subparagraph{Directional extensions}~

\begin{table}[H]
	\caption{Matal directional morphology}
	\label{tab:mata}
	\begin{tabular}{llll} % add l for every additional column or remove as necessary
		\lsptoprule
		forms & source gloss & direction  & other functions  \\ %table header
		\midrule
		\textit{awâŋ}  &  \textsc{vnt}  &  \textsc{vent} &   \\
%		\textit{álâ} & \textsc{itv} & & \\
		\textit{həŋ}  &  {down}  &    \textsc{down} &  lexicalized \\
		\textit{ábà} & {inside} & \textsc{into} (uncovered) & \textsc{loc}  \\
		\textit{águ} & {inside} & \textsc{into} (covered) &   \\
		\textit{àwdá} & {out} & \textsc{out} &   \\
		\lspbottomrule
	\end{tabular}
\end{table} 

\citet{dougopheGrammarMatal2021} describes two directional extensions: the ventive \textit{awâŋ} and the itive \textit{álâ} (sometimes glossed `away'). The latter is not given in any examples that clearly demonstrate a directional function, so it is not included as a directional extension in Table \ref{tab:mata}. Rather, examples given of \textit{álâ} seem to frequently involve events which have a transforming effect on the object, similar in distribution to aspectual markers sometimes called ``totality'' markers, as in (\ref{ex2:mata4}).

\ea\label{ex2:mata4}
\langinfo{Matal verb \textit{pàɗ} `eat' with putative itive extension}{}{\cite[219]{dougopheGrammarMatal2021}} \\
\gll à-dá-pàɗ-àx-àl-àŋ álá	\\
3\gl{sg}-\gl{pfv}-eat-\gl{pl}-3\gl{sg}.\gl{dat}-3\gl{sg}.\gl{obj} away \\
\glt `She had eaten him up.'
\z 

Most examples of the ventive extension occur with the verb stem \textit{s-} `come', as in (\ref{ex2:mata1}). In this case, the semantic contribution of the extension appears to be redundant. The verb \textit{s-} nearly always occurs with a directional extension of some type, but there are a few exceptions showing that this is not a morphosyntactic requirement. 

\ea\label{ex2:mata1}
\langinfo{Matal verb \textit{s} `come' with ventive extension}{}{\cite[26]{dougopheGrammarMatal2021}} \\
\gll  kà-s-àwáŋ mà là mtə̀gà=áj \\
2\gl{sg}-come-\gl{vent} \gl{top} \gl{loc} home=\gl{q} \\
\glt `Are you coming from home?'
\z 

There are also a few examples of the ventive extension occurring with other motion verbs, as in (\ref{ex2:mata2}), in which case its directional function is more evident. 

\ea\label{ex2:mata2}
\langinfo{Matal verb \textit{nə́f} `follow' with ventive extension}{}{\cite[188]{dougopheGrammarMatal2021}} \\
\gll ánw kə̀njj à-dá-də̀ mə̀tə́gá ɮà mə̀-nə́f-ák-àwáŋ 	\\
1\gl{pl} too 3\gl{sg}-\gl{pfv}-go home then 1\gl{sg}-follow-2\gl{sg}.\gl{obj}-\gl{vent}		\\
\glt `As for us, we thought you went home and we followed you this way.'
\z 

The ventive extension is normally transcribed as a suffix on the verb, in some cases following other verbal suffixes, as in (\ref{ex2:mata2}). There are also a couple examples where an object noun phrase occurs between the verb and the ventive extension, as in (\ref{ex2:mata3}). 

\ea\label{ex2:mata3}
\langinfo{Matal verb \textit{tə̀ɗ} `fetch' with ventive extension}{}{\cite[190]{dougopheGrammarMatal2021}} \\
\gll tá-tə̀ɗ-nì jàw àwáŋ mà-dàb-ə̀ŋ à bə́zj ágàj	\\
3\gl{pl}-fetch.\gl{prs}-\textsc{partitive} water \gl{vent} \gl{inf}-taste-3\gl{sg}.\gl{obj} to child already		\\
\glt `They fetch some water and give it to the child already.'
\z 

The directional morphemes \textit{həŋ} `down' and \textit{àwdá} `out' could also be considered directional extensions, at least in some cases (Séraphine Dougophe, personal communication). The directional use of the downward extension can be seen with intransitive verbs, as in (\ref{ex2:mata5}), as well as with transitive verbs, as in (\ref{ex2:mata6}).\footnote{The morpheme \textit{həŋ} also has other functions, including being used as part of an existential predicate and in demonstrative phrases \citep[39, 95, 157]{dougopheGrammarMatal2021}.} 

\ea\label{ex2:mata5}
\langinfo{Matal verb \textit{d} `go' with downward extension}{}{\cite[208]{dougopheGrammarMatal2021}} \\
\gll áfàʦ á-d xə̀ŋ là gàj gwə́də̀ŋ rə́əzj àdágàj \\
sun 3\gl{sg}-go.\gl{prs} down \gl{loc} mouth mountain \gl{ideo} already \\ 
\glt `The sun is setting. (lit., As the sun was going down the mouth of the mountain.)'

\ex\label{ex2:mata6}
\langinfo{Matal verb \textit{bək} `put' with downward extension}{}{\cite[421]{dougopheGrammarMatal2021}} \\
\gll Ma-bək həŋ	\\
\gl{inf}-put \gl{down}	\\
\glt `to put down'
\z 

The outward extension \textit{àwdá} is also used in a directional function with intransitive verbs, as in (\ref{ex2:mata7}), and transitive verbs, as in (\ref{ex2:mata8}). 

\ea\label{ex2:mata7}
\langinfo{Matal verb \textit{vàl} `jump' with outward extension}{}{\cite[206]{dougopheGrammarMatal2021}} \\
\gll gə̀-vàl àwdá káŋà mà-d-àj mà-də́ɮ àzə́g wànáj	\\
1\gl{sg}-jump \gl{out} \gl{purp} \gl{inf}-go-?? \gl{inf}-collect okra \gl{dem}	\\
\glt `I jumped out to go and collect this okra.'

\ex\label{ex2:mata8}
\langinfo{Matal verb \textit{vàl} `jump' with outward extension}{}{\cite[69]{dougopheGrammarMatal2021}} \\
\gll təla masok yaw awudâ 		\\
month pour water \gl{out} \\
\glt `month to pour water outside'
\z 

There also appear to be two inward directional extensions. The form \textit{ába} is used for uncovered containers and \textit{águ} is used for covered containers \citep[76]{dougopheGrammarMatal2021}. Their use as directional extensions can be seen in (\ref{ex2:mata10}) and (\ref{ex2:mata11}).\footnote{Another possible directional extensions is the morpheme \textit{afik} which is variously glossed, `up', `on' or `top'. However, the available examples are somewhat ambiguous as to the function and morphosyntactic properties of this morpheme.}

\ea\label{ex2:mata10}
\langinfo{Matal verb \textit{d} `go' with inward extension}{}{\cite[180]{dougopheGrammarMatal2021}} \\
\gll áŋá ákàl à-dá-d ábà jkə̀nj á-s-àl=àw	\\
if fire 3\gl{sg}-\gl{pfv}-go \gl{into} too 3\gl{sg}-please-3\gl{sg}.\gl{dat}=\gl{neg}	\\
\glt `If heat enters (the pipe), it does not like/it is not good.'

\ex\label{ex2:mata11}
\langinfo{Matal verb \textit{sàkw} `put' with inward extension}{}{\cite[202]{dougopheGrammarMatal2021}} \\
\gll  wànáj xə̀ŋ ka mə̀-sàkw-àx-àj ka kwàndúrkú águ	\\
\gl{dem} \gl{down} \gl{top} 1\gl{sg}-put-\gl{pl}-?? \gl{top} dolo \gl{into}		\\
\glt `As for this one, we put dolo (an unfermented sorghum beverage) inside.'
\z 

In one example, an inward suffix is used with a static verb \textit{xə̀ŋ} `stay', shown in (\ref{ex2:mata17}). In this case, the interpretation is locative rather than directional. It may be possible to analyze \textit{ábà} as a post-position in this context, especially if the verb \textit{xə̀ŋ} `stay' is an intransitive verb. However, this would not fit with descriptions of Matal as a prepositional language.

\ea\label{ex2:mata17}
\langinfo{Matal verb \textit{xə̀ŋ} `stay' with inward extension}{}{\cite[291]{dougopheGrammarMatal2021}} \\
\gll hwə̀jp à-dá-xə̀ŋ jáw ábà kà	\\
\gl{ideo}(spend.one.day) 3\gl{sg}-\gl{pfv}-stay water inside \gl{top}		\\
\glt `After it spends one day in the water...'
\z 

There is some ambiguity about the morphosyntactic status of these directional extensions. The ventive and itive extensions are described as ``full phonological words that can stand alone but are not immune to phonological reductions that transform them into proclitics.'' In certain syntactic contexts, these morphemes might not be considered directional extensions, such as (\ref{ex2:mata9}) where the outward extensions is preceded by a preposition and has a locative interpretation. In the same example, the downward extension appears following a demonstrative and is also used in combination with a preposition to function as an existential predicate. 

\ea\label{ex2:mata9}
\langinfo{Matal outward extension in prepositional phrase}{}{\cite[110]{dougopheGrammarMatal2021}} \\
\gll lábátáj xə̀ŋ là àwdá ka ákàl là=xə̀ŋ	\\
\gl{dem}.\gl{dist} \gl{down} \gl{loc} \gl{out} \gl{top} fire \gl{prep}=\gl{down}		\\
\glt `There outside, there is fire.'
\z 

\citet[35]{verdizadeSelectedTopicsGrammar2018} describes these morphemes as ``locative adverbs'' which must be preceded by a preposition \textit{la} or \textit{à}. This restrictive pattern from the Matal New Testament translation data used by \citet[]{verdizadeSelectedTopicsGrammar2018} does not hold for Dougophe's Matal data. 

\subparagraph{Caused accompanied motion (CAM)}

Directed CAM can be expressed in Matal by the combination of a directional extension with a verb meaning `take', as in (\ref{ex2:mata12}) and (\ref{ex2:mata13}).

\ea\label{ex2:mata12}
\langinfo{Matal verb \textit{xə̀l} `take' with ventive extension}{}{\cite[294]{dougopheGrammarMatal2021}} \\
\gll  tá-xə̀l gwàʦàk àwáŋ áŋà  mà-ʦə̀ɗ-àj 	\\
3\gl{pl}-take hen \gl{vent} \gl{purp} \gl{inf}-cut-??		\\
\glt `They bring a hen to slaughter...'

\ex\label{ex2:mata13}
\langinfo{Matal verb \textit{kap} `take' with downward extension}{}{\cite[162]{dougopheGrammarMatal2021}} \\
\gll ma-kap həŋ 		\\
\gl{inf}-take \gl{down} \\
\glt `to take down'
\z 

There are also examples which appear to express directed CAM by the use of motion verbs with a comitative prepositional phrase, as in (\ref{ex2:mata14}). 

\ea\label{ex2:mata14}
\langinfo{Matal verb \textit{d} `go' with comitative phrase}{}{\cite[219]{dougopheGrammarMatal2021}} \\
\gll  à-dá-də̀-là-àtà mjà \\
3\gl{sg}-\gl{pfv}-go-with-3\gl{pl} where \\ 
\glt `where had she gone with them (where had she taken them)'
\z

\subparagraph{Buy-sell verbs}

There are two separate verb stems in Matal, \textit{sə̀kw} `buy' and \textit{daw} `sell, throw'. In one case, example (\ref{ex2:mata15}), the verb \textit{sə̀kw} `buy' appears with a ventive extension. There is also one case, example (\ref{ex2:mata16}), in which the verb \textit{daw} `sell' appears with the putative itive extension. This is said to be ``because when something is sold, it is taken away from the vendor'' \citep[189]{dougopheGrammarMatal2021}. However, in most examples these verbs appear without an extension. 

\ea\label{ex2:mata15}
\langinfo{Matal verb \textit{sə̀kw} `buy' with ventive extension}{}{\cite[188]{dougopheGrammarMatal2021}} \\
\gll ʣá 	dájdàj  ábàj 	 kà á-dàj	 à 	kwàsə̀kwà 	ábàj	 à-sə̀kw-àwàŋ \\
person	any	too	\gl{top} 3\gl{sg}-go	to	market 		too	3\gl{sg}-buy.\gl{prs}-\gl{vent} \\
\glt `Anyone can go and get it from the market.'

\ex\label{ex2:mata16}
\langinfo{Matal verb \textit{də̀w} `sell' with putative itive extension}{}{\cite[189]{dougopheGrammarMatal2021}} \\
\gll ŋgàxà ká 	ká-dàj 	bj 	jsə́ŋw 	kwə̀lw 	bj 	ʣə́mj	tá-də̀w-àkà 	álá \\
and 	\gl{top} 	2\gl{sg}-go 	either	 money 	ten 	or 	five 	3\gl{pl}-sell-2\gl{sg}.\gl{dat} \gl{itive} \\
\glt `Then you go with ten or five francs they will sell it (give away) to you.'
\z 

\paragraph*{Parkwa/Podoko (park1239)} \label{park}

Information on Parkwa (also known as Podoko) is from two journal articles by the same author \citep{jarvisPodokoVerbalDirectionals1983,jarvisEsquisseGrammaticalePodoko1989}. The two publications take slightly different approaches to parsing the verbal morphology which complicates the discussion of directional extensions below. 

\subparagraph{Directional verbs}~

\begin{table}[H]
	\caption{Parkwa directional verbs} 
	\label{tab:parkV}
	\begin{tabular}{lll} % add l for every additional column or remove as necessary
		\lsptoprule
		DIR & verb & gloss   \\ %table header
		\midrule
		ITIVE & \textit{d-} & `go' \\
		VENT & \textit{s-} & `come' \\
		\lspbottomrule
	\end{tabular}
\end{table}

\subparagraph{Directional extensions}~

\begin{table}[H]
	\caption{Parkwa directional morphology}
	\label{tab:park}
	\begin{tabular}{llll} % add l for every additional column or remove as necessary
		\lsptoprule
		forms & source gloss & direction  & other functions  \\ %table header
		\midrule
		\textit{-ətsa, -atsa}  &  {to here}  &  \textsc{vent} & \textsc{subs.vent}  \\
		\textit{-u, -əlu, -alu}  &  {up}  & \textsc{up}  &  \textsc{subs.up}, lexicalized \\
		\textit{-aha, -əla, -ada}  &  {down}  &    \textsc{down} &  \textsc{subs.down} \\
		\textit{-əkwa, -akwa}  &  {in}  &   \textsc{into}  &  \textsc{subs.into}, \textsc{loc}   \\
		\textit{-edá, -əlá, -adá} & {out} & \textsc{out} &  \textsc{subs.out}  \\
		\textit{-arə} & on & \textsc{onto} & \textsc{subs.onto}, lexicalized \\
		\textit{-asə} & under & \textsc{under} &  \textsc{loc}, lexicalized   \\
		\lspbottomrule
	\end{tabular}
\end{table} 

\citet{jarvisPodokoVerbalDirectionals1983} and \citet{jarvisEsquisseGrammaticalePodoko1989} both describe motion-related verbal morphology in Parkwa (or Podoko) with slightly different approaches to the morphological parsing. \citet[322--323]{jarvisPodokoVerbalDirectionals1983} simplifies the view of directional morphology segmenting the phonemes that are common to each of the verbal endings in Table \ref{tab:park} and treating only those segments as directional suffixes. This results in the suffixes \textit{-tsa} `to here', \textit{-u}, `up', \textit{-a} `down', \textit{-kwa} `in' and \textit{dá} `out'. However, the preceding morphemes \textit{ə}, \textit{a}, \textit{əl} and \textit{ad} are left unexplained. 

\citet[89, 92]{jarvisEsquisseGrammaticalePodoko1989} takes two approaches to these verbal endings. On the one hand, \textit{-əl} and \textit{-ad} are analyzed as marking the (in)definiteness of the object, which Jarvis describes (in French) as \textit{-əl} `\textit{entier} [whole]' for definite and \textit{-ad} `\textit{partiel} [partial]' for indefinite. When the (in)definite object marker analysis is applied, \citet[92]{jarvisEsquisseGrammaticalePodoko1989} essentially ignores the fact that there is always a vowel following the \textit{-əl}/\textit{-ad} which makes the form identical to a directional suffix: \textit{á} `\textsc{out}', \textit{a} `\textsc{down}' or \textit{u} `\textsc{up}'. This is written off as a case of lexically-conditioned allomorphy (``\textit{selon le verbe} [according to the verb]'').\footnote{There are actually several exceptions to this where the final vowel is glossed as a directional suffix \citep[94, 96, 117]{jarvisEsquisseGrammaticalePodoko1989}. None of the exceptions have translations that include motion meaning.} 

The second approach that \citet[92]{jarvisEsquisseGrammaticalePodoko1989} applies is reflected in Table \ref{tab:park} in which the verbal endings are not segmented so that \textit{-əl} and \textit{-ad} are treated as part of a directional suffix. The result of applying two different analyses to the same verbal endings is that the ``\textit{partiel}'' and ``\textit{entier}'' suffixes are identical to forms of the upward, downward and outward directional suffixes. \citet[89]{jarvisEsquisseGrammaticalePodoko1989} admits that this dual-meaning analysis of the verbal endings results in a certain amount of subjectivity.\footnote{{}``\textit{Il est souvent difficile de savoir si la voyelle finale fait partie des marques du l'entier et du partiel, ou si c'est un suffix directionnel qui a perdu son sens premier.} [It is often difficult to know if the final vowel is part of the markers of the whole (definite) and the partial (indefinite), or if it is a directional suffix that has lost its primary sense.]'' } \citet[92]{jarvisEsquisseGrammaticalePodoko1989} further claims that five of the directional suffixes have intranstive and transitive forms, but her approach does not account for all of the directional forms shown in Table \ref{tab:park}. 

%\begin{table}[H]
%	\caption{Parkwa directional morphology (Jarvis 1989)}
%	\label{tab:parkJarvis}
%	\begin{tabular}{lll} % add l for every additional column or remove as necessary
	%		\lsptoprule
	%	 source gloss & intransitive  & transitive \\ %table header
	%		\midrule
	%	  {to here}  &	\textit{-ətsa}  &  	\textit{-atsa}   \\
	%	  {up}  & \textit{-u}  &  \textit{-əlu} \\
	%	 {down}  &    \textit{-aha} &  \textit{-əla} \\
	%	  {in}  &   	\textit{-əkwa}  &  	\textit{-akwa}  \\
	%	{out} & \textit{-edá} &   \textit{-əlá}  \\
	%		\lspbottomrule
	%	\end{tabular}
%\end{table} 

An improved account of the different forms of directional suffixes in Parkwa is shown in Table \ref{tab:park2}. In this approach, it is assumed that directional extensions may have allomorphs depending on whether the verb is intranstive, transitive with a definite object, or transitive with an indefinite object. Note that for some directional extensions, the same form is used in two or three of these contexts. 

\begin{table}
	\caption{Parkwa directional morphology (re-analyzed)}
	\label{tab:park2}
	\begin{tabular}{llll} % add l for every additional column or remove as necessary
		\lsptoprule
		source gloss & intransitive  & definite object & indefinite object \\ %table header
		\midrule
		\textsc{vent}  &	\textit{-ətsa}  &  	\textit{-atsa}  & no data  \\
		{\textsc{up}}  & \textit{-u}  &  \textit{-əlu, -alu} & \textit{-u} \\
		{\textsc{down}}  &    \textit{-aha} &  \textit{-əla} & \textit{-ada} \\
		{\textsc{in}}  &   	\textit{-əkwa}  &  	\textit{-əkwa, -akwa} & no data \\
		{\textsc{out}} & \textit{-edá} &   \textit{-əlá} & \textit{-adá}  \\
		{\textsc{onto}} & \textit{-arə} &  \textit{-arə}  & \textit{-arə}   \\
		{\textsc{under}} & \textit{-asə} &  \textit{-asə} &  \textit{-asə}   \\
		\lspbottomrule
	\end{tabular}
\end{table} 

One exception to the pattern is shown in (\ref{ex2:park1}) where the verb \textit{p-} `put' (here glossed `take') has the outward suffix \textit{-adá} despite the translation indicating a definite object.

\ea\label{ex2:park1}
\langinfo{Parkwa verb \textit{p} `put' with transitive ventive suffix}{}{\cite[115]{jarvisEsquisseGrammaticalePodoko1989}} \\
\gll  dá p-adá pə bwərngwə-bwərngwa məná sərá	\\
\gl{fut} take-\gl{out} take fetish 3\gl{sg}.\gl{poss} two \\
\glt `[When you go into the granary,] take his two fetishes.'
\z 

It is also left unexplained why the inward suffix \textit{-əkwa} is used for both intransitive predicates and those with a definite object despite their apparently being a separate form \textit{-akwa} which is only used when a definite object is present. Without more data, it is unclear how to further improve on Jarvis's analyses. 

The first five suffixes in Table \ref{tab:park} (ventive, upward, downward, inward, outward) can have a directional meaning with intransitive or transitive motion verbs, as in (\ref{ex2:park2}) and (\ref{ex2:park3}). 

\ea\label{ex2:park2}
\langinfo{Parkwa verb \textit{mbar} `leave' with ventive suffix}{}{\cite[93]{jarvisEsquisseGrammaticalePodoko1989}} \\
\gll   ʸmbar-ətsa mətá	\\
leave-\gl{vent} 3\gl{pl}	\\
\glt `They came here.'

\ex\label{ex2:park3}
\langinfo{Parkwa verb \textit{p} `put' with ventive suffix}{}{\cite[92]{jarvisEsquisseGrammaticalePodoko1989}} \\
\gll a ʸp-atsa ʸpá	\\
\gl{foc} put.\gl{prf}-\gl{vent} put \\
\glt `He put it here.'
\z 

With the verb \textit{kəs-} `take', they have a directed caused accompanied motion (CAM) meaning, as discussed below, and with non-motion verbs they can have either no motion meaning, as in (\ref{ex2:park4}), or subsequent directed CAM meaning, as in (\ref{ex2:park5}).  

\ea\label{ex2:park4}
\langinfo{Parkwa verb \textit{rəw} `die' with downward suffix}{}{\cite[95]{jarvisEsquisseGrammaticalePodoko1989}} \\
\gll rəw-aha mətá	\\
die.\gl{pl}-\gl{down} 3\gl{pl} \\
\glt `They died.'

\ex\label{ex2:park5}
\langinfo{Parkwa verb \textit{t} `cook' with downward suffix}{}{\cite[322]{jarvisPodokoVerbalDirectionals1983}} \\
\gll t-ad-a mayə dafá	\\
cook-??-\gl{down} 1\gl{sg} fufu	 \\
\glt `I cooked some fufu and took it down.'
\z 

Jarvis gives several examples of the verb stem \textit{ngwad} `attach, bind' with a directional suffix that has a subsequent directed CAM meaning, as in (\ref{ex2:park6}). However, one example of the same verb with a directional suffix does not have a subsequent directed CAM interpretation, shown in (\ref{ex2:park7}). In this case, the locative phrase is a container, so it is not interpreted as the destination of the agent's path of motion, but rather the predicate is interpreted as a type of motion in which the patient alone is moving along a path. This indicates that the subsequent directed CAM interpretation is, at least to some degree, dependent on a particular context. 

\ea\label{ex2:park6}
\langinfo{Parkwa verb \textit{ngwad} `bind' with inward suffix}{}{\cite[315]{jarvisPodokoVerbalDirectionals1983}} \\
\gll ngwaɗ-akwa mayə́ paní daká.	\\
bind-\gl{into} 1\gl{sg} stalk into.house	\\
\glt `I tied up the stalks and took them into the house.'

\ex\label{ex2:park7}
\langinfo{Parkwa verb \textit{ngwad} `bind' with inward suffix}{}{\cite[315]{jarvisPodokoVerbalDirectionals1983}} \\
\gll ʸngwaɗ-akwa mayə́ həyá da bəhwa \\
bind-\gl{into} 1\gl{sg} corn \gl{prep} sack \\
\glt `I put the corn in the sack and tied it up.'
\z

The suffix \textit{-arə} `onto' has the same functions as other directionals, except that no CAM examples found. The suffix \textit{-asə} `under' can have a directional interpretation with motion verbs, as in (\ref{ex2:park8}). In most cases where it occurs with non-motion verbs, there is no motion-related meaning,\footnote{\citet[320]{jarvisPodokoVerbalDirectionals1983} notes that the `under' suffix may be used when referring to any part of the body below the waist.} but in just one example, shown in (\ref{ex2:park9}), the interpretation appears to be locative. 

\ea\label{ex2:park8}
\langinfo{Parkwa verb \textit{d} `go' with `under' suffix}{}{\cite[315]{jarvisPodokoVerbalDirectionals1983}} \\
\gll d-asə mayə́ akə́ sla də́ zlə́ma \\
go-\gl{under} 1\gl{sg} \gl{prep} cow \gl{prep} stable	\\
\glt `I went into the cowshed. (Literally, `I went under to the cow in the stable.'' The stable is lower than the rest of the house.)'

\ex\label{ex2:park9}
\langinfo{Parkwa verb \textit{ngwad} `bind' with `under' suffix}{}{\cite[320]{jarvisPodokoVerbalDirectionals1983}} \\
\gll vats-asə mayə́ kará. \\
light-\gl{under} 1\gl{sg} fire \\
\glt `I lit the fire underneath.'
\z 

%324 "Whether or not -əl- can be rightly termed the 3rd.sg. object marker in series with the object marker for other persons, it is certainly some sort of object marker, for it occurs only with transitive verbs and never co-occurs (at least not in any unambiguous examples) with one of the other object markers. In contrast with -d- it carries an idea of definiteness or wholeness as opposed to indefiniteness or partitiveness... Whatever the precised meaning of these markers, the directionals -a and -á cannot occur without them if there is no other suffix present in the verb." 

%325 DOWN "It occurs in particular with verbs that have to do with the passing of time."

\subparagraph{Caused accompanied motion (CAM)}

Five of the directional suffixes (ventive, upward, downward, inward, outward) are shown to have a directed caused accompanied motion (CAM) meaning when attached to the verb \textit{kəs-} `take', as in (\ref{ex2:park10}) and (\ref{ex2:park11}). However, the directional suffix is not required for this verb to take on a directed CAM meaning. In (\ref{ex2:park12}) the verb \textit{kəs-} `take' with no directional suffix still has an outward CAM interpretation when it occurs with a prepositional phrase expressing outward motion. 

\ea\label{ex2:park10}
\langinfo{Parkwa verb \textit{kəs} `take' with downward suffix}{}{\cite[93]{jarvisEsquisseGrammaticalePodoko1989}} \\
\gll kəs-əla mayə́ kwərá \\
take-\gl{down} 1\gl{sg} stone	\\
\glt `I carried the stone down.'

\ex\label{ex2:park11}
\langinfo{Parkwa verb \textit{kəs} `take' with outward suffix}{}{\cite[314]{jarvisPodokoVerbalDirectionals1983}} \\
\gll kəs-əlá mayə́ həyə́ da udá \\
take-\gl{out} I corn to outside \\
\glt `I carried the corn outside.'

\ex\label{ex2:park12}
\langinfo{Parkwa verb \textit{kəs} `take' with outward prepositional phrase}{}{\cite[314]{jarvisPodokoVerbalDirectionals1983}} \\
\gll a ʸkəsə həyə́ yá da udá \\
\gl{foc} take corn 1\gl{sg} \gl{prep} outside \\
\glt `I am carrying the corn outside.'
\z 

With non-motion verbs, directional suffixes may have a subsequent direct CAM interpretation, as discussed in regard to (\ref{ex2:park5}) and (\ref{ex2:park6}) above.

%315 ``Some of the directionals have figurative as well as literal usages. For instance, -arə "on" can have the figurative meaning of "in addition to"''

%316 ``Another, -asə `under', sometimes seems to have the figurative meaning of `under the direction or leadership of someone'.''

%CAM metaphor? Kəs-ad-á mayə́ ɗalá	take-?-\gl{out} 1\gl{sg} song		I chose a song.		325

\subparagraph{Buy-sell verbs}

There are two verb roots in Parkwa, \textit{vəl-} `sell' and \textit{səkw-} `buy'. As seen in the following examples, both verbs may be used with directional suffixes without any change to their main meaning. The suffix \textit{-asə} `under' in (\ref{ex2:park14}) is said to be used because a skirt is worn on the lower part of the body. The \textit{-arə} `onto' in (\ref{ex2:park15}) is said to contribute the meaning `in addition to'. 

\ea\label{ex2:park13}
\langinfo{Parkwa verb \textit{vəl} `sell' with upward suffix}{}{\cite[319]{jarvisPodokoVerbalDirectionals1983}} \\
\gll ʸvəl-əlu mayə́ ndərə mayá 	\\
sell-\gl{up} 1\gl{sg} peanut 1\gl{sg}.\gl{poss}	\\
\glt `I sold my peanuts.'

\ex\label{ex2:park14}
\langinfo{Parkwa verb \textit{səkw-} `buy' with `under' suffix}{}{\cite[320]{jarvisPodokoVerbalDirectionals1983}} \\
\gll səkw-asə mayə́ patari akə́ nəs-ála. \\
buy-\gl{under} 1\gl{sg} skirt \gl{prep} wife-1\gl{sg}.\gl{poss}	\\
\glt `I bought a skirt for my wife to wear.'

\ex\label{ex2:park15}
\langinfo{Parkwa verb \textit{səkw-} `buy' with `onto' suffix}{}{\cite[94]{jarvisEsquisseGrammaticalePodoko1989}} \\
\gll  səkw-arə mayə́ burə akə́ sləɓa \\
buy-\gl{onto} 1\gl{sg} salt \gl{prep} meat	\\
\glt `I bought salt in addition to meat.'
\z 

\paragraph*{Glavda (glav1244)} \label{glav}

There are three short sketches of aspects of Glavda morphosyntax \citep{rappDictionaryGlavdaLanguage1968,bubaGlavdaMorphology2007,nghagyiyaSketchGlavdaGrammar2011} and a collection of word lists and transcribed texts with partial glossing \citep{owensGlavdaRuralUrban2011}.\footnote{In particular, the document ``verb\_vocab-final''. See also https://www.arabistik.uni-bayreuth.de/en/research/previous-research-projects/glavda/index.html [accessed 5 April 2022].} Note that these sources do not agree on the phonological inventory of the language, particularly its vowels. \citet{bubaGlavdaMorphology2007} include an appendix which shows that Glavda has three to five short vowels and four or five long vowels. This sharply contrasts with \citet[8]{nghagyiyaSketchGlavdaGrammar2011} who analyzes Glavda as only having one phonemic vowel /a/ occurring in long and short form. All other vowels are said to be allophones or epenthetic vowels. The sources do not claim to present a comprehensive analysis of the verbal morphology, and not only do they differ in which morphemes they analyze, but they also occasionally contradict each other in their semantic analyses. This could potentially be (dialect) variation or there may be gaps in the available data which would explain the apparent discrepancies. 

\subparagraph{Directional verbs}~

\begin{table}[H]
	\caption{Glavda directional verbs} 
	\label{tab:glavV}
	\begin{tabular}{lll} % add l for every additional column or remove as necessary
		\lsptoprule
		DIR & verb & gloss   \\ %table header
		\midrule
		ITIVE & \textit{d-} & `go' \\
		VENT & \textit{s-} & `come' \\
		\lspbottomrule
	\end{tabular}
\end{table}

\newpage
\subparagraph{Directional extensions}~

\begin{table}[H]
	\caption{Glavda directional morphology}
	\label{tab:glav}
	\begin{tabular}{llll} % add l for every additional column or remove as necessary
		\lsptoprule
		forms & source gloss & direction  & other functions  \\ %table header
		\midrule
		%\textit{-al}  &  \textsc{v.ex}  &  ?  &   ? \\
		\textit{-da}  &  \textsc{ext, v.ex}  &  \textsc{itive}  &  \textsc{subs.itive}, lexicalized   \\
		\textit{-m, -dəm}  &  \textsc{ext, v.ex}  & \textsc{into}   &  \textsc{loc}  \\	
		\textit{-it, -dit}  &  \textsc{ext, up}  & \textsc{up}  & \textsc{subs.up}, lexicalized    \\
		\textit{-xi, -di}  &  \textsc{ext, down}  &  \textsc{down}  & \textsc{subs.down}  \\
		%\textit{-di}  &  \textsc{down, di}  &   &   \\
		\textit{-s}  &  \textsc{v.ex, s}  & \textsc{under, itive}   &  lexicalized  \\
		\lspbottomrule
	\end{tabular}
\end{table} 

\citet[125--126]{rappDictionaryGlavdaLanguage1968} list eight extensions in Glavda. \citet[662]{bubaGlavdaMorphology2007} “collected about 30 unanalyzed verbal extensions” in Glavda.\footnote{A note on glossing and segmentation: \citet{bubaGlavdaMorphology2007} use \textsc{ext} to gloss all verbal ``extensions'', \citet{nghagyiyaSketchGlavdaGrammar2011} uses \textsc{v.ex} for most with the exception of \textsc{up} and \textsc{down}. I resegment their examples according to the analysis explained in this section, separating the \textit{-d(a)} from the directional suffixes. The directional suffixes are labeled \textsc{itive}, \textsc{up}, \textsc{down} and \textsc{into}. Other suffixes are labeled by their form, as in \textsc{al} for \textit{-al}.} \citet[22--27]{nghagyiyaSketchGlavdaGrammar2011} has a list of 16 Glavda verbal extensions. In all three sources, several of the verbal extensions have directional meanings. In this section, those forms said to have directional meaning are discussed. It is noteworthy that none of the extensions described are said to express ventive direction.

The suffix \textit{-da} is said to express itive directional meaning, but it appears that \textit{-da} only has a directional interpretation when it combines with a transitive verb root whose patient can be construed as a figure on a path of motion, as in (\ref{ex2:glav1}). All examples of \textit{-da} involve transitive verbs with the exception of examples where it combines with the verb \textit{d-} `go' or \textit{s-} `come' to express caused accompanied motion, discussed below.

\ea\label{ex2:glav1}
\langinfo{Glavda verb \textit{ɬəg-} `push' with \textit{-da} suffix }{}{\cite[662]{bubaGlavdaMorphology2007}} \\
\gll ɬəg-a-dá-ɬəga\\
push-3-\textsc{da}-push \\
\glt `He has pushed away.'
\z

In two examples, the translation indicates a subsequent motion interpretation of \textit{-da}, as in (\ref{ex2:glav11}). 

\ea\label{ex2:glav11}
\langinfo{Glavda verb \textit{čag-} `choose' with \textit{-da} suffix }{}{\cite[]{owensGlavdaRuralUrban2011}} \\
\gll čag-da	\\ 
choose-\gl{itive}	\\
\glt `pick and carry away'
\z 

There are also many examples where the form \textit{-da} combines with a non-motion verb and makes no discernible change to the meaning of the predicate, as in (\ref{ex2:glav10}), or forms a seemingly unrelated predicate, as in (\ref{ex2:glav100}). 

\ea\label{ex2:glav10}
\langinfo{Glavda verb \textit{mbəď-} `converse' with \textit{-da} suffix }{}{\cite[]{owensGlavdaRuralUrban2011}} \\
\gll mbəď-da	 \\
converse-\gl{itive}	 \\
\glt `converse'

\ex\label{ex2:glav100}
\langinfo{Glavda verb \textit{badz-} `spoil' with \textit{-da} suffix }{}{\cite[]{owensGlavdaRuralUrban2011}} \\
badz-da	\\
spoil-\gl{itive} \\
\glt `anger'
\z 

Upward directional motion with an intransitive verb is expressed by the suffix \textit{-it}, as in (\ref{ex2:glav17}) and (\ref{ex2:glav18}). 

\ea\label{ex2:glav17}
\langinfo{Glavda verb \textit{d-} `go' with \textit{-ít} suffix }{}{\cite[]{owensGlavdaRuralUrban2011}} \\
\gll d-it	\\
go-\gl{up} \\
\glt `go up, go up into, climb up'

\ex\label{ex2:glav18}
\langinfo{Glavda verb \textit{ďal-} `climb' with \textit{-ít} suffix }{}{\cite[]{owensGlavdaRuralUrban2011}} \\
\gll ďal-it	\\
climb-\gl{up}		\\
\glt `climb up something (tree, mountain)'
\z

\citet[662]{bubaGlavdaMorphology2007} describe a morpheme \textit{-dit} as having a meaning of either ‘upward’ or ‘westward’, as in (\ref{ex2:glav7}), since Glavda speakers live immediately to the east of the Mandara mountains.\footnote{Nghagyiya associates the ‘upward’ morphology with ‘eastward’ which is presumably a mistranslation given the geography of the region.} As was the case with \textit{-da} above, the distribution of \textit{-dit} appears to be limited to cases where the patient of a transitive verb is a figure on the path of motion. The suffix \textit{-dit} therefore appears to be an alternate from of the upward suffix \textit{-it} with their distribution determined by the transitivity of the verb.

\ea\label{ex2:glav7}
\langinfo{Glavda verb \textit{ɬəg-} `push' with upward suffix }{}{\cite[663]{bubaGlavdaMorphology2007}} \\
\gll ɬəg-a-dít-ɬəga \\
push-3-\gl{up}-push \\
\glt `He pushed s.t. up/from east to west.'
\z

\citet[26]{nghagyiyaSketchGlavdaGrammar2011} proposes the form \textit{-t} for upward motion when preceded by \textit{-d} (glossed `go') and as expressing `on' otherwise. \citet[662]{bubaGlavdaMorphology2007} describe \textit{-it} as expressing `on, partitive, reflexive' in addition to its upward function. These descriptions are likely due to the fact that most occurrences of \textit{-it} do not show any motion meaning related to the suffix, as in (\ref{ex2:glav19}), (\ref{ex2:glav20}) and (\ref{ex2:glav21}).

\ea\label{ex2:glav19}
\langinfo{Glavda verb \textit{łəg-} `push' with upward suffix }{}{\cite[]{owensGlavdaRuralUrban2011}} \\
\gll łəg-it	\\
push-\gl{up}	\\
\glt `it is pushable, it is enough'

\ex\label{ex2:glav20}
\langinfo{Glavda verb \textit{n-} `become' with upward suffix }{}{\cite[]{owensGlavdaRuralUrban2011}} \\
\gll n-it	\\
become-\gl{up} \\
\glt `become spoiled (of a person)'

\ex\label{ex2:glav21}
\langinfo{Glavda verb \textit{gat-} `seek' with upward suffix }{}{\cite[]{owensGlavdaRuralUrban2011}} \\
\gll gat-it	\\
seek-\gl{up} \\
\glt `find s.t.'
\z 

The suffix \textit{-xi} expresses downward motion, as in (\ref{ex2:glav4}). In (\ref{ex2:glav5}) the suffix is combined with a verb of breaking. The translation suggests that the directional meaning is applied to the orientation of the event, construing the object doing the breaking as a figure on a path of motion. 

\ea\label{ex2:glav4}
\langinfo{Glavda verb \textit{s} `come' with downward suffix }{}{\cite[24]{nghagyiyaSketchGlavdaGrammar2011}} \\
\gll s-àn-xí	sᵊ-g-à \\
come-1\gl{sg}-\gl{down}	come-??-\textsc{final.vowel}	\\
\glt `I came down.'

\ex\label{ex2:glav5}
\langinfo{Glavda verb \textit{kəl-} `break' with downward suffix }{}{\cite[]{owensGlavdaRuralUrban2011}} \\
\gll kəl-xi	\\
break-\gl{down}	\\
\glt `break downwards'
\z

\citet[24]{nghagyiyaSketchGlavdaGrammar2011} presents another downward motion suffix \textit{-di}, suggesting that it might be parsed \textit{-d} `go' and \textit{-i} `downward'. Note that, like the suffix \textit{-dit} above, verbs with a \textit{-di} form seem to only occur with transitive verbs, as in (\ref{ex2:glav6}). 

\ea\label{ex2:glav6}
\langinfo{Glavda verb \textit{ɗù:l} `throw' with downward suffix }{}{\cite[24]{nghagyiyaSketchGlavdaGrammar2011}} \\
\gll ɗù:l-á-dì	ɗú:l-g-à \\
throw-1\gl{sg}-\gl{down}	throw-??-\textsc{final.vowel}	 \\
\glt `I threw it downwards.'
\z 

It is also possible for the \textit{-di} to have an idiomatic interpretation as in (\ref{ex2:glav6a}). 

\ea\label{ex2:glav6a}
\langinfo{Glavda verb \textit{ts} `clap/strike' with \textit{-di}  suffix }{}{\cite[125]{rappDictionaryGlavdaLanguage1968}} \\
\gll ts-di \\
clap/strike-\gl{down}	\\
\glt `to subdue'
\z

%In support of an analysis of \textit{-di} as two morphemes,  \citet[125]{rappDictionaryGlavdaLanguage1968} show a suffix \textit{-i} with the verb root \textit{s} `come' with downward directional meaning, `to come down'. There are no other examples of the \textit{-i} suffix occurring on its own, but this may be due to a phonological rule deleting the /x/ phoneme. 

Nghagyiya includes two examples of a verbal suffix \textit{-m} with inward directional meaning, one with an intransitive verb, in (\ref{ex2:glav14}), and one with a transitive verb, as in (\ref{ex2:glav9}).

\ea\label{ex2:glav14}
\langinfo{Glavda verb \textit{d-} `go' with inward suffix }{}{\cite[27]{nghagyiyaSketchGlavdaGrammar2011}} \\
\gll dᵊ-m dᵊ-g dámá ɬó:dʒìà	\\
go-\gl{into} go-?? go.in soldier	\\
\glt `He has joined the army.'

\ex\label{ex2:glav9}
\langinfo{Glavda verb \textit{f-} `put' with inward suffix }{}{\cite[27]{nghagyiyaSketchGlavdaGrammar2011}} \\
\gll f-àrᵊ-m	fᵊ-g-à \\
put-3\gl{pl}-\gl{into}	put-?-FV \\
\glt `They put it in.'
\z 

\citet[663]{bubaGlavdaMorphology2007} propose two forms for the inward extension describing the form \textit{-im} as `in, into, reflexive' and \textit{-dəm} as `inside, in, inner surface', but do not include examples using these suffixes. Two examples of \textit{-dəm} appear in \citet{owensGlavdaRuralUrban2011}, including (\ref{ex2:glav22}). Note that the same verb appears in (\ref{ex2:glav9}). The two examples given of \textit{-dəm} are with transitive predicates which suggests a pattern as above where a \textit{d} morpheme appears with transitive predicates, however, there are other transitive predicates shown with the suffix \textit{-m} in both sources. This leaves the distributions of \textit{-m} and \textit{-dəm} unclear. 

\ea\label{ex2:glav22}
\langinfo{Glavda verb \textit{tal-} `put' with inward suffix }{}{\cite[]{owensGlavdaRuralUrban2011}} \\
\gll f-dəm \\
put-\gl{into} \\
\glt `put into'
\z 

Nghagyiya analyzes a verbal suffix \textit{-s} as having a direction-related meaning `under', as in (\ref{ex2:glav8}). The morpheme \textit{-s} is described by \citet[126]{rappDictionaryGlavdaLanguage1968} as having itive directional meaning when combined with the root \textit{ḅil} ‘send’, and when combined with the non-motion verb \textit{ngwud} ‘tie’ it has the idiomatic meaning ‘wear’. Examples of the \textit{-s} suffix in \citet{owensGlavdaRuralUrban2011} do not indicate any motion-related meaning.

\ea\label{ex2:glav8}
\langinfo{Glavda verb \textit{f-} `put' with \textit{-s}  suffix }{}{\cite[24]{nghagyiyaSketchGlavdaGrammar2011}} \\
\gll f-àn-ár-s	fᵊ-g	k-ká:rá \\
put-1\gl{sg}-??-\gl{under}	put-??	\gl{obj}-fire \\
\glt `I put fire under it.'
\z

\citet[25]{nghagyiyaSketchGlavdaGrammar2011} also describes \textit{-al} as a suffix that ``conveys the idea of `separation or apart'{''}, giving two examples with possible itive motion translations, shown in (\ref{ex2:glav3}) and (\ref{ex2:glav2}). \citet{bubaGlavdaMorphology2007} describe the function of \textit{-al} as ``apart, -able, reflexive, north–south or south–north direction'' but do not offer any examples with a directional meaning. \citet{owensGlavdaRuralUrban2011} has at least 25 examples of a verb root combining with \textit{-al}, none of which express itive motion, including the verb \textit{d-} `go', as in (\ref{ex2:glav13}). Note that the use of \textit{-al} correlates with intransitive verbs. This suggests a possible analysis of \textit{-al} as an intransitive allomorph of the itive directional \textit{-da}, but (\ref{ex2:glav24}) shows that the two suffixes can apparently co-occur.

\ea\label{ex2:glav3}
\langinfo{Glavda verb \textit{kl-} `break' with \textit{-al} suffix }{}{\cite[25]{nghagyiyaSketchGlavdaGrammar2011}} \\
\gll kl-ál	kᵊl-g-à \\
break-\textsc{al}	break-?-\textsc{final.vowel}	 \\
\glt `It broke away.'

\ex\label{ex2:glav2}
\langinfo{Glavda verb \textit{t͡ʃì-} `stand' with \textit{-al}  suffix }{}{\cite[25]{nghagyiyaSketchGlavdaGrammar2011}} \\
\gll t͡ʃì-ál	t͡ʃí:-g-à \\
stand-\textsc{al}	stand-?-\textsc{final.vowel} \\
\glt `He stood up and went away.'

\ex\label{ex2:glav13}
\langinfo{Glavda verb \textit{d-} `go' with \textit{-al}  suffix }{}{\cite[]{owensGlavdaRuralUrban2011}} \\
\gll d-al \\
go-\textsc{al} \\
\glt `finish (of an event), be over'
\z 

\subparagraph{Caused accompanied motion (CAM)}

Data from \citet{owensGlavdaRuralUrban2011} contains several morphological variations of the verbs \textit{d-} `go' and \textit{s-} `come' translated as itive and ventive caused accompanied motion, respectively. This usually involves the suffix \textit{-da}, as in (\ref{ex2:glav23}) and (\ref{ex2:glav12}). 

\ea\label{ex2:glav23}
\langinfo{Glavda verb \textit{d-} `go' with \textit{-da}  suffix }{}{\cite[]{owensGlavdaRuralUrban2011}} \\
\gll də-da \\
go-\gl{itive} \\
\glt `carry away, take away'

\ex\label{ex2:glav12}
\langinfo{Glavda verb \textit{s-} `come' with \textit{-da} suffix }{}{\cite[]{owensGlavdaRuralUrban2011}} \\
\gll s-da \\
come-\gl{itive} \\
\glt `bring, bring out, cause, explain'
\z 

Other suffixes can also be used in combination with \textit{-da}, as in (\ref{ex2:glav24}). 

\ea\label{ex2:glav24}
\langinfo{Glavda verb \textit{s-} `come' with \textit{-da}  suffix }{}{\cite[]{owensGlavdaRuralUrban2011}} \\
\gll d-al-da \\
go-\textsc{al}-\gl{itive} \\
\glt `carry off to'
\z 

When combined with an intransitive motion verb, the (transitive) downward suffix \textit{-di} also has a CAM interpretation, as in (\ref{ex2:glav16}) \citep[cf.][125]{rappDictionaryGlavdaLanguage1968}. This morphological analysis is complicated by a verb phrase identical except for tone which has the translation `get down', shown in (\ref{ex2:glav17a}). 

\ea\label{ex2:glav16}
\langinfo{Glavda verb \textit{s-} `come' with \textit{-di}  suffix (high tone) }{}{\cite[]{owensGlavdaRuralUrban2011}} \\
\gll sə-dií	\\
come-\gl{down}		\\
\glt `bring down'

\ex\label{ex2:glav17a}
\langinfo{Glavda verb \textit{s-} `come' with \textit{-di}  suffix (low tone) }{}{\cite[]{owensGlavdaRuralUrban2011}} \\
\gll sə-dii	\\
come-\gl{down}	\\
\glt `get down'
\z 

A suffix \textit{-v} with directional motion verbs can also have a CAM interpretation as in (\ref{ex2:glav25}) and (\ref{ex2:glav26}). \citet[22]{nghagyiyaSketchGlavdaGrammar2011} describes \textit{-v} as meaning `on or around oneself' and \citet[662]{bubaGlavdaMorphology2007} describe it as meaning `on, benefactive, reflexive'. 

\ea\label{ex2:glav25}
\langinfo{Glavda verb \textit{d-} `go' with \textit{-v}  suffix }{}{\cite[]{owensGlavdaRuralUrban2011}} \\
\gll d-əv \\
go-\textsc{v} \\
\glt `carry to, take away to, be gone'

\ex\label{ex2:glav26}
\langinfo{Glavda verb \textit{s-} `come' with \textit{-v}  suffix }{}{\cite[]{owensGlavdaRuralUrban2011}} \\
\gll sә-sә-v	\\ 	
come-??-\textsc{v} \\
\glt `bring with'
\z

\subparagraph{Buy-sell verbs}

The verb root \textit{səgw-} can mean both `buy' and `sell' \citep[648, 672]{bubaGlavdaMorphology2007}. Selling can also be expressed by the verb \textit{və̀l} `give, sell' \citep[642]{bubaGlavdaMorphology2007}.

\paragraph*{Malgwa (malg1251)} \label{malg}

Information on Malgwa grammar is from a reference grammar \citep{lohrSpracheMalgwaNara2002}.

\subparagraph{Directional verbs}~

\begin{table}[H]
	\caption{Malgwa directional verbs} 
	\label{tab:malgV}
	\begin{tabular}{lll} % add l for every additional column or remove as necessary
		\lsptoprule
		DIR & verb & gloss   \\ %table header
		\midrule
		ITIVE & \textit{də́} & `\textit{gehen} [go]' \\
		VENT & \textit{sa} & `\textit{kommen} [come]' \\
		UP & \textit{ɗála} & `\textit{hinaufgehen}, \textit{aufsteigen}, \textit{klettern} [go up, rise, climb]' \\
		DOWN & \textit{tsə́kwa} & `\textit{heruntersteigen} [descend]' \\
		\lspbottomrule
	\end{tabular}
\end{table}

\subparagraph{Directional extensions}~

\begin{table}[H]
	\caption{Malgwa directional morphology}
	\label{tab:malg}
	\begin{tabular}{llll} % add l for every additional column or remove as necessary
		\lsptoprule
		forms & gloss & direction  & other functions  \\ %table header
		\midrule
		\textit{-tə́}  &  \textsc{up}  &   \textsc{up}   &    lexicalized \\
		\textit{-ə́m}  &  \textsc{into}  & \textsc{into}  &  \textsc{split}, lexicalized  \\
		\textit{-sə́}  &  \textsc{out}  &    \textsc{out} &  \textsc{ben}, lexicalized \\
		\textit{-ar}  &  \textsc{onto}  &   \textsc{onto}  &  lexicalized   \\
		\lspbottomrule
	\end{tabular}
\end{table} 

There are four verbal suffixes in Malgwa that are shown to have a directional function with motion verbs: upward, as in (\ref{ex2:malg1}), inward, as in (\ref{ex2:malg2}), outward, as in (\ref{ex2:malg3}), and `onto', as in (\ref{ex2:malg4}). Notably, none of these are a ventive suffix. 

\ea\label{ex2:malg1}
\langinfo{Malgwa verb \textit{d-} `go' with upward suffix }{}{\cite[153]{lohrSpracheMalgwaNara2002}} \\
\gll də́-tə́ \\
go-\gl{up} \\
\glt `to climb (e.g., up a tree)'

\ex\label{ex2:malg2}
\langinfo{Malgwa verb \textit{d-} `go' with inward suffix }{}{\cite[154]{lohrSpracheMalgwaNara2002}} \\
\gll  yaa d-ə́m adə́m bəré	 \\
1\gl{sg}.\gl{pfv} go-\gl{into} \gl{prep} house \\
\glt `I entered a house.' [German: `\textit{Ich betrat ein Haus}.']

\ex\label{ex2:malg3}
\langinfo{Malgwa verb \textit{d-} `go' with outward suffix }{}{\cite[161]{lohrSpracheMalgwaNara2002}} \\
\gll də́-sə́ \\
go-\gl{out} \\
\glt `to go out' [German: `\textit{hinausgehen}']

\ex\label{ex2:malg4}
\langinfo{Malgwa verb \textit{d-} `go' with \textsc{onto} suffix }{}{\cite[159]{lohrSpracheMalgwaNara2002}} \\
\gll bə́za-ár \\
jump-\gl{onto} \\
\glt `to jump over' [German: `\textit{darüberspringen}']
\z 

The above examples show the suffixes with intransitive motion verbs. Examples with transitive motion verbs were found with the inward and outward suffixes, as in (\ref{ex2:malg5}) and (\ref{ex2:malg6}), and possibly the `onto' suffix depending on the interpretation of (\ref{ex2:malg7}), but not the upward suffix. 

\ea\label{ex2:malg5}
\langinfo{Malgwa verb \textit{f-} `put' with inward suffix }{}{\cite[151]{lohrSpracheMalgwaNara2002}} \\
\gll fa-n-ə́m-fá \\
put-1\gl{sg}.\gl{prf}-\gl{into}-put \\
\glt `I put it in.'

\ex\label{ex2:malg6}
\langinfo{Malgwa verb \textit{f-} `put' with outward suffix }{}{\cite[152]{lohrSpracheMalgwaNara2002}} \\
\gll tá púwa-sə́ \\
3\gl{pl}.\gl{prs} pour-\gl{out} \\
\glt `They pour it out.'

\ex\label{ex2:malg7}
\langinfo{Malgwa verb \textit{f-} `put' with `onto' suffix }{}{\cite[151]{lohrSpracheMalgwaNara2002}} \\
\gll fa-́∅-tər-ár-fé \\
put-3\gl{sg}.\gl{prf}-\gl{obj}-\gl{onto}-put	\\
\glt `He raised it for her.'
\z 

When the directional suffixes combine with non-motion verbs, they receive non-motion interpretations which are frequently idiomatic, as in (\ref{ex2:malg8}) and (\ref{ex2:malg9}). 

\ea\label{ex2:malg8}
\langinfo{Malgwa verb \textit{sha} `drink' with upward suffix }{}{\cite[139]{lohrSpracheMalgwaNara2002}} \\
\gll sha-n-án-tə́-sha \\
drink-1\gl{sg}.\gl{prf}-\gl{tr}-\gl{up}-drink \\
\glt `I have received.' [German: `\textit{Ich habe erhalten}.']

\ex\label{ex2:malg9}
\langinfo{Malgwa verb \textit{pəsha} `sprinkle' with outward suffix }{}{\cite[152]{lohrSpracheMalgwaNara2002}} \\
\gll pəsha-r-sə́-pə́sha ə́ɬa \\
sprinkle-3\gl{pl}.\gl{prf}-\gl{out}-sprinkle cow \\
\glt `They milked the cow.' [German: `\textit{Sie haben die Kuh gemolken}.']
\z 

In example (\ref{ex2:malg18}), the inward directional suffix results in an interpretation that the object is split in two parts. 

\ea\label{ex2:malg18}
\langinfo{Malgwa verb \textit{kya} `break' with inward suffix }{}{\cite[162]{lohrSpracheMalgwaNara2002}} \\
\gll kya-an-ə́m	\\
break-\gl{tr}-\gl{into}	\\
\glt `to break in two' [German: `\textit{in zwei Teile brechen}']
\z

In at least some cases, directional suffixes can have an idiomatic non-motion interpretation even when suffixed to a motion verb, as in (\ref{ex2:malg10}), (\ref{ex2:malg11}) and (\ref{ex2:malg12}) (cf. \ref{ex2:malg6} and \ref{ex2:malg7}).

\ea\label{ex2:malg10}
\langinfo{Malgwa verb \textit{f-} `put' with outward suffix }{}{\cite[151]{lohrSpracheMalgwaNara2002}} \\
\gll yaa fá-n-sə́-fá \\
1\gl{sg}.\gl{pfv} put-\gl{dat}-\gl{out}-put \\
\glt `I built it for him' [German: `\textit{Ich baute es für ihn}.']

\ex\label{ex2:malg11}
\langinfo{Malgwa verb \textit{f-} `put' with upward and `onto' suffixes}{}{\cite[139]{lohrSpracheMalgwaNara2002}} \\
\gll fə̄-té-re \\
put-\gl{up}-\gl{onto}	\\
\glt `trust' [German: `\textit{vertrauen}']

\ex\label{ex2:malg12}
\langinfo{Malgwa verb \textit{f-} `put' with `onto' suffix }{}{\cite[160]{lohrSpracheMalgwaNara2002}} \\
\gll fa-́∅-tər-ár-fé \\
put-3\gl{sg}.\gl{prf}-\gl{dat}-\gl{onto}-put \\
\glt `He increased it for them.' [German: `\textit{Er hat es für sie} (\gl{pl}) \textit{erhöht}.']
\z 

\subparagraph{Caused accompanied motion (CAM)}

Directed caused accompanied motion (CAM) appears to be most regularly expressed by the use of two verbal suffixes that increase the valency of the intranstive motion verbs `go' and `come', as in (\ref{ex2:malg13}) and (\ref{ex2:malg14}). 

\ea\label{ex2:malg13}
\langinfo{Malgwa verb \textit{d-} `go' with transitive and outward suffixes}{}{\cite[139]{lohrSpracheMalgwaNara2002}} \\
\gll d-an-sé \\
go-\gl{tr}-\gl{out} \\
\glt `take away' [German: `\textit{wegnehmen}']

\ex\label{ex2:malg14}
\langinfo{Malgwa verb \textit{s-} `come' with applicative suffix}{}{\cite[156]{lohrSpracheMalgwaNara2002}} \\
\gll yá sa-kur-áá ɗáfa \\ 
1\gl{sg}.\gl{prs} come-\gl{dat}-\gl{appl} food \\
\glt `I bring you food.' [German: `\textit{Ich bringe euch Speisen}.']
\z 

\subparagraph{Buy-sell verbs}

There are two separate verbs in Malgwa: \textit{shə́kwa} `buy' and \textit{və́la} `sell'. The verb  \textit{və́la} `sell' is mostly attested with the totality suffix \textit{-úú}, as in (\ref{ex2:malg15}), and the verb \textit{shə́kwa} `buy' is mostly found with a suffix not found elsewhere, \textit{-fə́}, as in (\ref{ex2:malg16}), though this is not obligatory, as seen in (\ref{ex2:malg17}) \citep[158, 160]{lohrSpracheMalgwaNara2002}. 

\ea\label{ex2:malg15}
\langinfo{Malgwa verb \textit{və́la} `sell' with totality suffix}{}{\cite[159]{lohrSpracheMalgwaNara2002}} \\
\gll tá vəla-tər-úú-vəlé \\
3\gl{pl}.\gl{prs} sell-\gl{dat}-\gl{tot}-sell \\
\glt `They will sell it to you.' [German: `\textit{Sie werden es ihnen verkaufen}.']

\ex\label{ex2:malg16}
\langinfo{Malgwa verb \textit{shakwa} `buy' with fossilized \textit{-fə́} suffix}{}{\cite[190]{lohrSpracheMalgwaNara2002}} \\
\gll itäre shakwa-r-fa-ga näyaba \\
\gl{sbj} buy-3\gl{pl}.\gl{prf}-\textsc{fa}-\gl{neg} banana  \\
\glt `they didn't buy bananas.' [German: `\textit{Sie kauften keine Bananen}.']

\ex\label{ex2:malg17}
\langinfo{Malgwa verb \textit{shakwa} `buy' }{}{\cite[232]{lohrSpracheMalgwaNara2002}} \\
\gll yaa shə́kwa baúwəkəní ɓáŋa ba náyaba \\
1\gl{sg}.\gl{pfv} buy \gl{indf} ?? \gl{foc} banana \\
\glt `I bought everything but bananas.' [German: `\textit{Ich kaufte alles außer Bananen}.']
\z 

\paragraph*{Wandala (wand1278)} \label{wand}

Information on Wandala grammar is primarily from a reference grammar \citep{frajzyngierGrammarWandala2012}.

\subparagraph{Directional verbs}~

\begin{table}[H]
	\caption{Wandala directional verbs} 
	\label{tab:wandV}
	\begin{tabular}{lll} % add l for every additional column or remove as necessary
		\lsptoprule
		DIR & verb & gloss   \\ %table header
		\midrule
		ITIVE & \textit{d-} & `go' \\
		VENT & \textit{s-} & `come' \\
		\lspbottomrule
	\end{tabular}
\end{table}

\subparagraph{Directional extensions}~

\begin{table}[H]
	\caption{Wandala directional morphology}
	\label{tab:wand}
	\begin{tabular}{llll} % add l for every additional column or remove as necessary
		\lsptoprule
		forms & source gloss & direction  & other functions  \\ %table header
		\midrule
		\textit{-m}  &  \textsc{in}  & \textsc{into}  &  \\
		\textit{-s}  &  \textsc{s}  &    \textsc{out} &   \\
		\textit{-ar}  &  \textsc{on}  &  \textsc{onto} &   lexicalized  \\
		\lspbottomrule
	\end{tabular}
\end{table} 

There are three motion-related verbal suffixes in Wandala: \textit{-m} `inward', \textit{-s} outward and \textit{-ar} `onto'.\footnote{\citet{frajzyngierGrammarWandala2012} refers to the inward suffix as the ``inner-space extension'' glossed \textsc{in}, the outward suffix as the ``source-oriented extension'' glossed \textsc{s} and the `onto' suffix as the ``role-changing extension'' glossed \textsc{on}. The abbreviations list also has \textsc{out} for ``Extension `out'{''} but this gloss does not appear elsewhere. Frajzyingier's descriptions of the meanings of each of these are difficult to diagnose and, in some cases, are contradicted by the data.} These suffixes have directional meaning with intransitive and transitive motion verbs, as shown with the inward suffix in (\ref{ex2:wand1}) and (\ref{ex2:wand2}).

\ea\label{ex2:wand1}
\langinfo{Wandala verb \textit{d-} `go' with inward suffix}{}{\cite[140]{frajzyngierGrammarWandala2012}} \\
\gll má də́-m tàtá-r kínì má nábà ɮàlà \\
\textsc{hypothetical} go-\gl{into} place-\gl{q} \gl{foc} 1\gl{incl} then go \\
\glt `No matter into what place, we can go there.'

\ex\label{ex2:wand2}
\langinfo{Wandala verb \textit{pu} `pour' with inward suffix}{}{\cite[269]{frajzyngierGrammarWandala2012}} \\
\gll pù-m-pwà \\
pour-\gl{into}-pour.\gl{tr}	\\
\glt `pour in'
\z 

At times, the same Wandala phrase is given different translations, and in a few cases it can be seen that at least some free translations of examples with directional suffixes are omitting directional information in the English translation. The same phrase, \textit{má də́m tàtàr}, that has an inward motion translation in (\ref{ex2:wand1}) also appears in (\ref{ex2:wand3}) but without an inward directional meaning in the translation. This raises the possibility that directional semantics are lost in translation elsewhere as well. 

\ea\label{ex2:wand3}
\langinfo{Wandala verb \textit{d-} `go' with inward suffix}{}{\cite[96]{frajzyngierGrammarWandala2012}} \\
\gll má də́-m tàtà-r kínì á wà-mì kà ɮyáwà \\
1\gl{incl} go-\gl{into} place-\gl{q} \gl{foc} 3\gl{sg} bite-1\gl{incl} \gl{neg} fear \\
\glt `Wherever we would go the fear will not bite us.'
\z 

In addition to the above suffixes, there is a suffix \textit{-t} called a ``target extension'' described as ``indicating movement toward a target'' and meaning ``that an object or subject has a new location as the result of the event'' \citep[16, 286]{frajzyngierGrammarWandala2012}. This analysis is not supported by the examples given of verbs with the suffix \textit{-t}. It is more likely that \textit{-t} is a cognate of the Malgwa upward suffix \textit{-tə́} \citep[153]{lohrSpracheMalgwaNara2002}. In the Wandala data, the \textit{-t} suffix is most frequently found with the verb \textit{tsə́} `rise, get up' and is described as having a ``perfectivizing effect'' similar to other upward suffixes \citep[285]{frajzyngierGrammarWandala2012}. Since no examples are given showing \textit{-t} with upward directional meaning in Wandala, it will be assumed that this suffix has grammaticalized away from directional meaning. 

\citet[248--256]{frajzyngierGrammarWandala2012} describes the suffix \textit{-w} as a ``ventive extension'' claiming that, in its locative function, the suffix can ``indicate movement toward the place of the speaker''. The only examples of this suffix with ventive directional meaning are when it is suffixed to the ventive directional verb \textit{s-} `come'. This suffix appears to actually be an aspectual marker, most likely related to the Malgwa ``totality'' suffix \textit{-úú} \citep[158]{lohrSpracheMalgwaNara2002}. 

\subparagraph{Caused accompanied motion (CAM)}

Directed caused accompanied motion (CAM) is expressed in transitivized forms of the basic directional verbs \textit{s-} `come' and \textit{d-} `go', as in (\ref{ex2:wand5}) and (\ref{ex2:wand6}). \citet[271--272]{frajzyngierGrammarWandala2012} calls these forms the ``goal extension'' glossed \textsc{go}, but that label is misleading as the forms are not related to locative arguments or anything specifically motion-related, but do appear to alter the valency of the verb root. For that reason Frajzyingier's ``goal marker'' is glossed as a transitive form.

\ea\label{ex2:wand5}
\langinfo{Wandala verb \textit{s-} `come' in transitive form}{}{\cite[78]{frajzyngierGrammarWandala2012}} \\
\gll tà s-á wè\\
3\gl{pl} come-\gl{tr} what\\
\glt `What did they bring?'

\ex\label{ex2:wand6}
\langinfo{Wandala verb \textit{d-} `go' in transitive form}{}{\cite[612]{frajzyngierGrammarWandala2012}} \\
\gll
tá d-á názù tá də́ ndàvà-ŋ-án ŋànè\\
3\gl{pl} go-\gl{tr} that 3\gl{sg} \gl{fut} ask-3\gl{sg}-\gl{asoc} 3\gl{sg} \\
\glt `They should bring that with which they are going to ask for her.'
\z 

These directed CAM verbs can further be modified by a directional suffix, as in the outward suffix \textit{-s} in (\ref{ex2:wand7}). 

\ea\label{ex2:wand7}
\langinfo{Wandala verb \textit{s-} `come' in transitive form with outward suffix}{}{\cite[272]{frajzyngierGrammarWandala2012}} \\
\gll tà də̀ sə̀-s-á dàdà á-m úvgè \\
3\gl{pl} \gl{seq} come-\gl{out}-\gl{tr} father \gl{prep}-into grave \\
\glt `They raised the father from the grave.'
\z

\subparagraph{Buy-sell verbs}

There are two distinct verb roots in Wandala: \textit{və̀lə̀} `sell' and \textit{ʃkwà} `buy', as in (\ref{ex2:wand4}). The verb \textit{və̀lə̀} `sell' is also said to mean `send' \citep[155]{frajzyngierGrammarWandala2012}.

\ea\label{ex2:wand4}
\langinfo{Wandala verbs \textit{və̀lə̀} `sell' and \textit{ʃkwà} `buy'}{}{\cite[156]{frajzyngierGrammarWandala2012}} \\
\gll ká və̀lə̀ mtú ká ʃkwà hè \\
2\gl{sg} sell or 2\gl{sg} buy \gl{q} \\
\glt `Are you selling or are you buying?'
\z 

\subsubsection{Mofuic (7/9)}

The Mofuic languages are subdivided into three groups. Of the three Meri languages, Merey and Zulgo-Gemzek are discussed below. No information is available about Dugwor (dugw1239). Of the two Mofu languages, Mofu-Gudur is discussed below, but there is no information about verbal morphology in North Mofu (nort3046). All four Tokombere languages are included: Mada, Moloko, Muyang and Wuzlam. 

\paragraph*{Merey (mere1246)} \label{mere}

Information on Merey grammar is primarily from a grammar sketch published as an online manuscript \citep{gravinaVerbPhraseMerey2007} and a lexicon \citep{gravinaMereyProvisionalLexicon2003}.

\newpage
\subparagraph{Directional verbs}~

\begin{table}[H]
	\caption{Merey directional verbs} 
	\label{tab:mereV}
	\begin{tabular}{lll} % add l for every additional column or remove as necessary
		\lsptoprule
		DIR & verb & gloss   \\ %table header
		\midrule
		MOTION & \textit{de} & `\textit{aller} [go]' \\
		UP & \textit{tsal} & `\textit{monter} [go up]' \\
		DOWN & \textit{mbəzla} & `\textit{descendre} [descend]' \\
		\lspbottomrule
	\end{tabular}
\end{table}

Merey directional verbs are from \citet{gravinaMereyProvisionalLexicon2003}. The translation of \textit{de} suggests that it is an itive verb, but since it is compatible with a ventive directional extensions (see below) and since there does not appear to be a ventive directional verb root (e.g., `come'), it is assumed that \textit{de} is actually a general motion verb unspecified for direction. 

\subparagraph{Directional extensions}~

\begin{table}[H]
	\caption{Merey directional morphology}
	\label{tab:mere}
	\begin{tabular}{llll} % add l for every additional column or remove as necessary
		\lsptoprule
		forms & source gloss & direction  & other functions  \\ %table header
		\midrule
		\textit{-aw}  &  \textsc{dir}  &  \textsc{vent} &   \textsc{subs.vent}  \\
		\lspbottomrule
	\end{tabular}
\end{table} 

\citet[14--15]{gravinaVerbPhraseMerey2007} describes one directional suffix \textit{-aw} with ventive meaning: ``The direction of the verb is always to ‘here’.'' Six examples are given of this ventive suffix. One example has a translation indicating ventive directional meaning with an intransitive motion verb, shown in (\ref{ex2:mere1}). 

\ea\label{ex2:mere1}
\langinfo{Merey verb \textit{y} `leave' with ventive suffix}{}{\cite[17]{gravinaVerbPhraseMerey2007}} \\
\gll na y-aw mə-tə-ɗe yam.	\\
1\gl{sg}.\gl{compl} leave-\gl{vent} \gl{inf}-draw-\gl{inf} water	\\
\glt `I came to draw water.'
\z 

In (\ref{ex2:mere2}), the verb form \textit{do} is glossed as a combination of the verb \textit{d} `go/come' with the ventive suffix \textit{-aw}. Presumably the form of the verb can be explained by a phonological process (/d-aw/ → [do]) but this is not explained and no morpheme breaks are given in the source text. Confusingly, the translation is itive rather than ventive direction, contradicting Gravina's analysis.

\ea\label{ex2:mere2}
\langinfo{Merey verb \textit{d} `go/come' in ventive form }{}{\cite[12]{gravinaVerbPhraseMerey2007}} \\
\gll d-o a gay i Maɗaf mә Gemzek.	\\
go/come-\gl{vent} to house of Maɗaf in Gemzek		\\
\glt `Go to Maɗaf’s house in Gemzek.'
\z 

With non-motion verbs, two examples have no motion-related translation, and in one, shown in (\ref{ex2:mere3}), there is a subsequent (caused accompanied) ventive motion translation.

\ea\label{ex2:mere3}
\langinfo{Merey verb \textit{səkəm} `buy' with ventive suffix}{}{\cite[15]{gravinaVerbPhraseMerey2007}} \\
\gll  na d-iye mata səkəm-aw wah ka təv i Mbəzum hay	\\
1\gl{sg}.\gl{incl} go-\gl{incl} for buy-\gl{vent} milk on place of Fulani \gl{pl}	\\
\glt `I am going in order to buy milk at the place where the Fulani are (and bring it back).'
\z 

\subparagraph{Caused accompanied motion (CAM)}

It appears (in one example) that ventive caused accompanied motion can be expressed by the use of the ventive suffix with the verb, as shown in (\ref{ex2:mere4}).

\ea\label{ex2:mere4}
\langinfo{Merey verb \textit{zla} `take' with ventive suffix}{}{\cite[18]{gravinaVerbPhraseMerey2007}} \\
\gll  ti ye faya wunaka a Matakam mə-zl-aw \\
3\gl{pl}.\gl{compl} go on.him ?? to Mafa \gl{inf}-take-\gl{vent} \\
\glt `They went to Mafa to bring back [a wife].'
\z 

\subparagraph{Buy-sell verbs}

The verb \textit{səkəm} means `buy' and maintains that meaning with a ventive suffix, as in (\ref{ex2:mere3}). The concept of `sell' is expressed by the verb \textit{kal} `throw' followed by a verbal particle \textit{ha} which has a transitivizing function \citep[14]{gravinaVerbPhraseMerey2007}.

\paragraph*{Zulgo-Gemzek (zulg1242)} \label{zulg}

Zulgo and Gemzek are listed together as a single language in the Ethnologue and Glottologue, but \citet[5]{gravinaOutlineSketchGemzek2005} considers them separate languages on phonological grounds. The two lects have slightly different forms of motion-related verbal morphology as well. They will be discussed side-by-side here as available descriptions of each lect are somewhat limited. 

Information on Zulgo-Gemzek grammar is from a grammar sketch in a journal article \citep{hallerVerbalComplexZulgo1981}, two sketches published as online manuscripts \citep{gravinaOutlineSketchGemzek2005,scherrerOverviewGemzekNarrative2012} and an unpublished lexicon \citep{amadouZulgoLexicon1986}.

\subparagraph{Directional verbs}~

\begin{table}[H]
	\caption{Zulgo-Gemzek directional verbs} 
	\label{tab:zulgV}
	\begin{tabular}{lll} % add l for every additional column or remove as necessary
		\lsptoprule
		DIR & verb & gloss   \\ %table header
		\midrule
		MOTION & \textit{d-} & `go' \\
		INTO & \textit{péts} & `enter' \\
		UP & \textit{tsə́l} &  `\textit{monter} [go up]' \\
		DOWN & \textit{tàr} & `\textit{descendre} [descend]'   \\
		\lspbottomrule
	\end{tabular}
\end{table}

In \citet{amadouZulgoLexicon1986}, subentries for the verb \textit{dá} `\textit{aller }[go]' include the expressions \textit{dá (ádəm)} `\textit{entrer} [enter]' and \textit{dá (iɗə́m á)} `\textit{sortir} [exit]'. In these cases, inward and outward motion are indicated by the apparent adverbial expressions \textit{áɗəm} `\textit{dedans} [inside]' and \textit{iɗə́m á} `\textit{dehors} [outside]'.

\subparagraph{Directional extensions}~

\begin{table}[H]
	\caption{Zulgo-Gemzek directional morphology}
	\label{tab:zulg}
	\begin{tabular}{llll} % add l for every additional column or remove as necessary
		\lsptoprule
		forms & source gloss & direction  & other functions  \\ %table header
		\midrule
		\textit{-ara} (Zulgo)  &  \textsc{egr}  &  \textsc{vent} &  \textsc{subs.vent} \\
		\textit{-aw, -ya} (Gemzek)  &  \textsc{dir}  &  \textsc{vent}  &  \textsc{subs.vent} \\
		\textit{aha}  &  \textsc{ingr}, \textsc{dir}  & \textsc{itive} & \textsc{subs.vent/itive} \\
		\lspbottomrule
	\end{tabular}
\end{table} 

The ventive form in Zulgo is \textit{-ára} and in Gemzek it is \textit{-aw} (which suppletes to \textit{-ya} following a habitual suffix).\footnote{\citet[46--48]{hallerVerbalComplexZulgo1981} use the term ``egressive'' (outward) for the ventive suffix. Their description of the meaning of motion suffixes refers to an undefined ``reference point'' other than the deictic center: ``the motion is away from the reference point and towards the location of the speaker or original observer.'' Following this pattern, they use ``ingressive'' to describe the itive marker in Zulgo.} In both lects, it can be seen that the verb \textit{d-} glossed `go' is a non-directional motion verb which may take an itive reading in the unmarked case, but can also have a ventive directional meaning, as in (\ref{ex2:zulg1}). In just a few examples of transitive motion verbs with a ventive suffix, there is no directional interpretation. 

\ea\label{ex2:zulg1}
\langinfo{Zulgo verb \textit{d-} `go/come' with ventive suffix}{}{\cite[48]{hallerVerbalComplexZulgo1981}} \\
\gll  á-dá-ára á gá-yá \\
he-go/come-\gl{vent} into house-here \\
\glt `He came into the house.'
\z 

With non-motion verbs, there can be a subsequent motion interpretation of the ventive suffix, as in (\ref{ex2:zulg2}). The translations of these examples typically include a caused accompanied motion meaning.

\ea\label{ex2:zulg2}
\langinfo{Gemzek verb \textit{sékém} `buy' with ventive suffix}{}{\cite[11]{gravinaOutlineSketchGemzek2005}} \\
\gll na sékém-aw sla ya\\
1s buy-\gl{vent} cow here	\\
\glt `I bought a cow (and brought it home).'
\z 

\citet{hallerVerbalComplexZulgo1981} also describe a round-trip interpretation of the ventive suffix, as in (\ref{ex2:zulg3}). However, the same verbs appear in \citet{scherrerOverviewGemzekNarrative2012} with a ventive suffix and no motion meaning in the translation. This could be a difference between Zulgo and Gemzek, an issue of context-dependent interpretation or information lost in the translation of the Gemzek text. 

\ea\label{ex2:zulg3}
\langinfo{Zulgo verb \textit{sékém} `buy' with ventive suffix}{}{\cite[48]{hallerVerbalComplexZulgo1981}} \\
\gll á-tə̀f-ára yam-á \\
she-drew-\gl{vent} water-here \\
\glt `She went, drew water and brought it here.'
\z 

In both Gemzek and Zulgo, there is an itive marker \textit{aha} which is described as a suffix in Zulgo and as a particle in Gemzek. It is shown to express itive directional meaning with intransitive and transitive motion verbs, as in (\ref{ex2:zulg4}) and (\ref{ex2:zulg5}).\footnote{When the marker \textit{aha} occurs the verb \textit{d-} `go', it usually  has an itive translation, but both \citet[47]{hallerVerbalComplexZulgo1981} and \citet[18]{gravinaOutlineSketchGemzek2005} include one example each of this form with ventive translation. I assume these are mistranslated examples.}

\ea\label{ex2:zulg4}
\langinfo{Zulgo verb \textit{val} `run' with itive suffix}{}{\cite[47]{hallerVerbalComplexZulgo1981}} \\
\gll á-val-áha \\
he-ran-\gl{itive} \\
\glt `He ran there.'

\ex\label{ex2:zulg5}
\langinfo{Gemzek verb \textit{slér} `send' with  itive marker}{}{\cite[11]{gravinaOutlineSketchGemzek2005}} \\
\gll Na slér aha kéla ga a tév kurom. \\
1\gl{sg} send \gl{itive} child my at place your	\\
\glt `I sent my child to your house.'
\z 

The itive marker can also have a subsequent (caused accompanied) motion interpretation, at least in Zulgo, as seen in (\ref{ex2:zulg6}) and (\ref{ex2:zulg7}). However, note that the translations inexplicably differ as to whether the direction of the subsequent motion is itive or ventive. When the itive marker appears with non-motion verbs in examples from \citet{scherrerOverviewGemzekNarrative2012}, it does not receive any motion interpretation in the translation. 

\ea\label{ex2:zulg6}
\langinfo{Zulgo verb \textit{zə́m} `eat' with itive suffix}{}{\cite[47]{hallerVerbalComplexZulgo1981}} \\
\gll á-zə́m-áha daf \\
he-ate-\gl{itive} fufu \\
\glt `He ate fufu and went there.'

\ex\label{ex2:zulg7}
\langinfo{Zulgo verb \textit{zə́m} `eat' with itive suffix}{}{\cite[47]{hallerVerbalComplexZulgo1981}} \\
\gll á-vak-áha mendzíkwir \\
she-roasted-\gl{itive} chicken \\
\glt `She roasted a chicken and brought it here.'
\z 

\subparagraph{Caused accompanied motion (CAM)}

In one example, (\ref{ex2:zulg8}), a transitivized form of the basic motion verb is combined with a ventive suffix to express ventive caused accompanied motion. 

\ea\label{ex2:zulg8}
\langinfo{Gemzek verb \textit{d-} `go/come' with ventive suffix and transitivizer}{}{\cite[11]{gravinaOutlineSketchGemzek2005}} \\
\gll ka d-aw di ya a ma ya \\
2\gl{sg} go/come-\gl{vent} \gl{tr} here to house here \\
\glt `You bring it to the house.'
\z 

Directed CAM can also be expressed by combining the ventive suffix with the verb \textit{zla} `take', as in (\ref{ex2:zulg9}).

\ea\label{ex2:zulg9}
\langinfo{Gemzek verb \textit{zla} `take' with ventive suffix}{}{\cite[23]{gravinaOutlineSketchGemzek2005}} \\
\gll kage ka yah ti kéver duñgway ti, ka méma, na ta zla-k-aw di ya \\
if 2 seek \gl{top} liver that \gl{top} 1\gl{pl} return-1\gl{du} 1\gl{sg} \gl{fut} take-2\gl{sg}.\gl{dat}-\gl{vent} \gl{tr} here	\\
\glt `If you are looking for that liver, we will return, I will bring you it.'
\z 

\subparagraph{Buy-sell verbs}

There are two separate verbs, \textit{sékém} `buy' (as in \ref{ex2:zulg2}) and \textit{kel} `sell' \citep[50]{hallerVerbalComplexZulgo1981}.

\paragraph*{Mofu-Gudur (mofu1248)} \label{mofu}

Information on Mofu-Gudur grammar is primarily from a grammar sketch and lexicon \citep{barreteauDescriptionMofuGudurLangue1983} with additional data from an unpublished grammar sketch \citep{hollingsworthTransitivityPragmaticsObject1995}. 

\subparagraph{Directional verbs}~

\begin{table}[H]
	\caption{Mofu-Gudur directional verbs} 
	\label{tab:mofuV}
	\begin{tabular}{lll} % add l for every additional column or remove as necessary
		\lsptoprule
		DIR & verb & gloss   \\ %table header
		\midrule
		ITIVE & \textit{daw} & `\textit{aller} [go]' \\
		VENT & \textit{sawa} & `\textit{venir} [come]' \\
		INTO& \textit{mbəz} & `\textit{entrer} [enter]' \\
		OUT & \textit{b} & `\textit{sortir}, \textit{quitter} (\textit{un endroit}) [exit, leave]' \\
		UP & \textit{təp} & `\textit{monter} [go up]' \\ 
		DOWN & \textit{bə́rŋg} & `\textit{descendre} [descend]' \\
		\lspbottomrule
	\end{tabular}
\end{table}

\subparagraph{Directional extensions}~

\begin{table}[H]
	\caption{Mofu-Gudur directional morphology}
	\label{tab:mofu}
	\begin{tabular}{llll} % add l for every additional column or remove as necessary
		\lsptoprule
		forms & source gloss & direction  & other functions  \\ %table header
		\midrule
		\textit{-wa}  &  \textit{rapprochement}, \textsc{dir}  &  \textsc{vent} &  \textsc{subs.vent}, \textsc{ben}  \\
		\lspbottomrule
	\end{tabular}
\end{table} 

The suffix \textit{-wa} is described as an \textit{extension de rapprochement} expressing centripetal (or ventive) motion \citep[424]{barreteauDescriptionMofuGudurLangue1983}. Ventive directional meaning can be seen in at least a small number of examples of this suffix, as in (\ref{ex2:mofu1}) and (\ref{ex2:mofu2}). 

\ea\label{ex2:mofu1}
\langinfo{Mofu verb \textit{há} `run' with ventive suffix}{}{\cite[70]{barreteauDescriptionMofuGudurLangue1983}} \\
\gll há-wa \\
run-\gl{vent} \\
\glt `Come here quickly!' [French: `\textit{Viens-vite ici} !']

\ex\label{ex2:mofu2}
\langinfo{Mofu verb \textit{bəl} `chase' with ventive suffix}{}{\cite[9]{hollingsworthTransitivityPragmaticsObject1995}} \\
\gll ta', a bəl-tər-wa \\
then 3 exit-\gl{obj}.\gl{pl}-\gl{vent} \\
\glt `Then, he chased them (sheep) back.'
\z 

There is also one example, shown in (\ref{ex2:mofu5}), with a subsequent (caused accompanied) motion interpretation that also triggers a possible round-trip interpretation.

\ea\label{ex2:mofu5}
\langinfo{Mofu verb \textit{səp} `search' with ventive suffix}{}{\cite[9]{hollingsworthTransitivityPragmaticsObject1995}} \\
\gll  a səp-ər-wa ɗakw aŋga	\\
3 search-\gl{obj}-\gl{vent} goat his		\\
\glt `(He goes and) looks for his goat (to bring it back).'
\z 

However, in the majority of cases, the suffix \textit{-wa} appears to have grammaticalized away from a directional function. \citet[8, 10]{hollingsworthTransitivityPragmaticsObject1995} describes a beneficiary meaning when the ventive co-occurs with a first- or second-person object suffix (compare \ref{ex2:mofu3} and \ref{ex2:mofu4}) and an ``individuated'' meaning when co-occurring with a third-person object suffix. 

\ea\label{ex2:mofu3}
\langinfo{Mofu verb \textit{jaw} `tie up' with  object suffix}{}{\cite[10]{hollingsworthTransitivityPragmaticsObject1995}} \\
\gll ya jaw-ka \\
1 tie.up-2\gl{sg}.\gl{obj}	\\
\glt `I tie you up.'

\ex\label{ex2:mofu4}
\langinfo{Mofu verb \textit{jaw} `tie up' with object suffix and ventive suffix}{}{\cite[10]{hollingsworthTransitivityPragmaticsObject1995}} \\
\gll ya jaw-ka-wa \\
1 tie.up-2\gl{sg}.\gl{obj}-\gl{vent} \\
\glt `I tie for you.'
\z 

%One peculiarity of the distribution of the suffix \textit{-wa} is that it has become part of the verb root \textit{sawa} `come' \citep[155]{barreteauDescriptionMofuGudurLangue1983}.

\citet{barreteauDescriptionMofuGudurLangue1983} also describes a second directional verbal suffix. The suffix \textit{-fá} is said to be quite rare and is identical in form to the preposition \textit{fá} `on (top of)'. Its function is described as orienting an event onto something,\footnote{{}``\textit{Le locatif -fá} (\textit{beaucoup plus rare}) \textit{oriente le procès exprimé par le verbe sur quelque chose} [The locative -\textit{fá} (much rarer) orients the process expressed by the verb on top of something]'' \citep[424]{barreteauDescriptionMofuGudurLangue1983}.} however the two examples found of this suffix do not have motion-related translations. Since there is no evidence available that the suffix maintains a directional function, it is not included in the discussion here. 

\subparagraph{Caused accompanied motion (CAM)}

Many examples show a verb \textit{la} glossed `\textit{prendre} [take]' or `\textit{emporter} [bring]' with a ventive suffix expressing CAM, as in (\ref{ex2:mofu6}). In most cases the translation does not make explicit whether the motion is ventive or not. 

\ea\label{ex2:mofu6}
\langinfo{Mofu verb \textit{la} `take' with ventive suffix}{}{\cite[459]{barreteauDescriptionMofuGudurLangue1983}} \\
\gll ka da la-wa ɗáf bá \\
you \gl{fut} bring-\gl{vent} boule \gl{neg}	\\
\glt `Do not bring boule.' [French: `\textit{N’emporte pas de boule de mil}.']
\z 

\subparagraph{Buy-sell verbs}

The verb \textit{heɗk} `buy' changes meaning in its causative form to \textit{heɗkad} `sell' \citep[67]{barreteauDescriptionMofuGudurLangue1983}.

\paragraph*{Mada (mada1293)} \label{mada}

Information on Mada grammar is from a short grammar sketch available as an online manuscript \citep{ernstkurdiTenseAspectMood2016} and a lexicon \citep{barreteau2000dictionnaire}.

\subparagraph{Directional verbs}~

\begin{table}[H]
	\caption{Mada directional verbs} 
	\label{tab:madaV}
	\begin{tabular}{lll} % add l for every additional column or remove as necessary
		\lsptoprule
		DIR & verb & gloss   \\ %table header
		\midrule
		MOTION & \textit{ódoro}/\textit{ádara} & `\textit{aller} [go]'/`\textit{venir} [come]' \\ 
		ITIVE & \textit{áŋga} & `\textit{aller} [go]' \\
		VENT & \textit{áwá} & `\textit{venir}, \textit{arriver} [come, arrive]' \\
		UP & \textit{ácala} & `\textit{monter} [go up]' \\
		\lspbottomrule
	\end{tabular}
\end{table}

One motion verb has itive meaning in its labialized form \textit{ódoro} and ventive meaning in its palatal form \textit{ádara}. This directional use of palatalization and labialization is also found with the verb \textit{ɬá} `go home' and the verb \textit{grá} `do' (with an altrilocative interpretation) \citep[29]{ernstkurdiTenseAspectMood2016}.

\largerpage[2]
\subparagraph{Directional extensions}~

\begin{table}[H]
	\caption{Mada directional morphology}
	\label{tab:mada}
	\begin{tabular}{llll} % add l for every additional column or remove as necessary
		\lsptoprule
		forms & source gloss & direction  & other functions  \\ %table header
		\midrule
		\textit{-re}, [palatal]  &  \textsc{vnt}, \textsc{loc}  &  \textsc{vent} &  \textsc{loc}  \\
		\textit{-ro}, [labial]  &  \textsc{itv}  & \textsc{itive}  &    \\
		\textit{-ra}  &  \textsc{egr}  &   \textsc{down}  &    \\
		\textit{-va}  &  \textsc{ing}  &    \textsc{into} &   \\
		\textit{-a}  &  \textsc{abl}  &   \textsc{into}   &    \\	
		\textit{-ka}  &  \textsc{dir}  &   \textsc{onto}  &   lexicalized  \\		
		\textit{-fa}  &  \textsc{ade}  &   \textsc{onto}  &   \textsc{loc}  \\	
%		\textit{-la}  &  \textsc{dir}  &   \textsc{along}  &  NA  \\	
		\lspbottomrule
	\end{tabular}
\end{table} 

\citet{ernstkurdiTenseAspectMood2016} describes Mada verbal morphology with a focus on the aspectual system, but also includes examples of some directional suffixes. Verbal morphology shown to have a directional function is included in Table \ref{tab:mada}. \citet[49]{barreteau2000dictionnaire} also list a number of ``extensions'' but without giving examples or including relevant morpheme breaks and glossing in the example sentences included in their dictionary. There are other potentially motion-related (or locative) suffixes not included here due to a lack of examples such as \textit{-ábà} `inside' and \textit{-hà} `under'.

In addition to the labialized/palatalized pairs discussed above, directional meaning can also be expressed with verbal suffixes including a ventive suffix \textit{-re} and an itive suffix \textit{-ro}. These suffixes are only shown to have directional meaning with transitive motion verbs, as in (\ref{ex2:mada1}).\footnote{In the one example of the ventive suffix \textit{-re} with an intransitive motion verb it occurs with the verb \textit{áŋga} `go' and the translation remains itive rather than ventive \citep[74]{ernstkurdiTenseAspectMood2016}. This inconsistency is unexplained.}

\ea\label{ex2:mada1}
\langinfo{Mada verb \textit{ple} `throw' with ventive suffix}{}{\cite[45]{ernstkurdiTenseAspectMood2016}} \\
\gll ē-ké-plē-ré \\
3\gl{sg}-\gl{prog}-throw-\gl{vent}	\\
\glt `He is throwing (it) (towards speaker).'
\z 

In the one example in which the ventive suffix occurs with a non-motion verb, shown in (\ref{ex2:mada2}), the translation is altrilocative. 

\ea\label{ex2:mada2}
\langinfo{Mada verb \textit{zlèc} `hit' with ventive suffix}{}{\cite[46]{ernstkurdiTenseAspectMood2016}} \\
\gll  zlèc-ēré \\
hit.\gl{imp}-\gl{vent}		\\
\glt `Hit (in another place)!'
\z 

Examples of the other directional suffixes are somewhat limited, most often being shown to have a directional function with transitive motion verbs. This may be in addition to other possible meanings, as in the case of the suffix \textit{-fa} which can have a locational meaning of `besides, next to' with non-motion verbs, as in \textit{ckàɗa-fá-ŋ} [sit.\gl{imp}-\gl{onto}-\gl{tr}] `sit / sit next to (it)' \citep[46]{ernstkurdiTenseAspectMood2016}.

\subparagraph{Caused accompanied motion (CAM)}

The verbs \textit{áhhal} `\textit{porter} [carry]' and \textit{ázá} `\textit{prendre}, \textit{emporter} [take (away)]' express caused accompanied motion and can be modified by a directional suffix, as in (\ref{ex2:mada4}). 

\ea\label{ex2:mada4}
\langinfo{Mada verb \textit{hàlà} `take' with downward suffix}{}{\cite[46]{ernstkurdiTenseAspectMood2016}} \\
\gll hàlà-rā	\\
carry.\gl{imp}-\gl{down} \\
\glt `take down / take out'
\z 

\subparagraph{Buy-sell verbs}

\citet[240]{barreteau2000dictionnaire} list \textit{óskwom} `buy' and the itive form \textit{óskwom-oro} `sell'. However, \citet{ernstkurdiTenseAspectMood2016} gives examples showing that the verb stem \textit{skóm} can mean `buy' or `sell' without a directional suffix, as in (\ref{ex2:mada5}), so it appears to be ambiguous between both meanings, but specified as `sell' with an itive suffix.

\ea\label{ex2:mada5}
\langinfo{Mada prior motion SVC?}{}{\cite[88]{ernstkurdiTenseAspectMood2016}} \\
\gll kwadzaŋ afalaŋa kà-wá ma nò-skóm-á awak nehe \\
tomorrow when 2\gl{sg}.\gl{irr}-come \gl{top} 1\gl{sg}.\gl{irr}-sell-\gl{pfv} goat \gl{det} \\
\glt `Tomorrow when you come, I will have sold this goat.'
\z 

\paragraph*{Moloko (molo1266)} \label{molo}

Moloko is described in a grammar by \citet{friesenGrammarMoloko2017}. 

\subparagraph{Directional verbs}~

\begin{table}[H]
	\caption{Moloko directional verbs} 
	\label{tab:moloV}
	\begin{tabular}{lll} % add l for every additional column or remove as necessary
		\lsptoprule
		DIR & verb & gloss   \\ %table header
		\midrule
		MOTION & \textit{lo} & `go, come' \\
		INTO& \textit{tar} & `enter' \\
		OUT & \textit{zləray} & `exit, go out' \\
		UP & \textit{ɓoroy} & `go up, climb' \\
		DOWN & \textit{fatay} & `descend' \\
		\lspbottomrule
	\end{tabular}
\end{table}

\largerpage[2]
The basic motion verb \textit{lo} in Moloko is not inherently directional. It can be given directional meaning by adding a ventive or itive suffix, as discussed below. 

\subparagraph{Directional extensions}~

\begin{table}[H]
	\caption{Moloko directional extensions}
	\label{tab:molo}
	\begin{tabular}{llll} % add l for every additional column or remove as necessary
		\lsptoprule
		forms & source gloss & direction  & other functions  \\ %table header
		\midrule
		\textit{ala}  &  \textsc{to}  &   \textsc{vent} &     \\
		\textit{alaj}  &  \textsc{away}  & \textsc{itive}  &  lexicalized  \\
		\lspbottomrule
	\end{tabular}
\end{table} 

There are two directional enclitics: ventive and itive. \citet[239]{friesenGrammarMoloko2017} also describes the verbal plural/pluractional \textit{aya} as a directional, but it does not contribute information about the path of motion. There are also two “adpositional enclitics”: \textit{aka} ‘on top’ and \textit{ava} ‘in’. These sometimes co-occur with directionals extensions. Like directionals, their semantic contribution is not always clear and may be lexicalized in some cases.

The verb \textit{lo} ‘go’ is a non-deictic motion verb that can occur on its own without a directional particle or with either the ventive or itive directional particle, as in (\ref{ex2:molo1}) and (\ref{ex2:molo2}). The examples in \citet{friesenGrammarMoloko2017} suggest that the ventive is used more frequently used with this verb, and that the unmarked form is frequently used in the context of itive motion.

\ea\label{ex2:molo1}
\langinfo{Moloko verb \textit{lo} `go/come' with ventive marker}{}{\cite[304]{friesenGrammarMoloko2017}} \\
\gll tə́-l=ala	a	həlaŋ	ga	ava	ɛʃɛ	tú-wuɗak=ala	har	a	mʊsijɔŋ	ava\\
3\gl{pl}-go/come.\gl{ipfv}=\gl{vent}	at	back	\gl{adj}	in	again	3\gl{pl}-divide.\gl{ipfv}=\gl{vent}	body at	mission	in \\
\glt `They come [home] again, they disperse after church.'

\ex\label{ex2:molo2}
\langinfo{Moloko verb \textit{lo} `go/come' with itive marker}{}{\cite[6]{friesenGrammarMoloko2017}} \\
\gll a-l=ala	na	gʷɔgʷɔlvaŋ	na	ɔ̀-lɔ=alaj\\
3\gl{sg}-go/come=\gl{vent}	\gl{top}	snake	\gl{top}	3\gl{sg}.\gl{pfv}-go/come=\gl{itive} \\
\glt `Some time later, the snake went.'
\z 

With non-motion verbs, there is just one example, shown in (\ref{ex2:molo3}), in which a ventive directional has an interpretation similar to subsequent ventive motion, but which \citet[242]{friesenGrammarMoloko2017} describes as indicating ``that the person has already eaten, but at some other location, since the directional gives the idea that food has come to the speaker.''

\ea\label{ex2:molo3}
\langinfo{Moloko verb \textit{zəm} `eat' with ventive marker }{}{\cite[242]{friesenGrammarMoloko2017}} \\
\gll nə̀-zəm=ala	tɔhʷɔ \\
1\gl{sg}.\gl{pfv}-eat=\gl{vent}	\gl{dem} \\
\glt `I already ate over there (some other person’s house – before I arrived here).'
\z 

In (\ref{ex2:molo4}), the verb `drink' appears with the same ventive directional with a perfective interpretation rather than subsequent associated motion. This may be an artifact of the translation process, or it is possible that “the directional =\textit{ala} may be in the process of becoming grammaticalised for past tense or a subtype of Perfect” \citep[243]{friesenGrammarMoloko2017}. 

\ea\label{ex2:molo4}
\langinfo{Moloko verb \textit{s} `eat' with ventive marker }{}{\cite[188]{friesenGrammarMoloko2017}} \\
\gll nà-s=ala \\
1\gl{sg}.\gl{pfv}-drink=\gl{vent} \\
\glt `I drank already. (lit., I drank towards)'
\z 

In many other cases where a directional particle combines with a non-motion verb, there is no indication what the function of the directional particle is, as in (\ref{ex2:molo5}) and (\ref{ex2:molo6}).

\ea\label{ex2:molo5}
\langinfo{Moloko verb \textit{mət} `die' with itive marker }{}{\cite[176]{friesenGrammarMoloko2017}} \\
\gll mə-mət=ava=alaj	a	vɛr	ava \\
\gl{inf}-die=in=\gl{itive}	at	room	in \\
\glt `She died in the room.'

\ex\label{ex2:molo6}
\langinfo{Moloko verb \textit{təkam} `taste' with itive marker }{}{\cite[344]{friesenGrammarMoloko2017}} \\
\gll Hʊrmbʊlɔm na ɓərav=ahaŋ a-təkam=alaj na a-vah-aj ɛlɛ=ahaŋ bɔtɔt\\
God \gl{top} heart=3\gl{sg}.\gl{poss} 3\gl{sg}-taste=\gl{itive} \gl{top} 3\gl{sg}-fly-?? thing=3\gl{sg}.\gl{poss} \gl{ideo}	\\
\glt `God (for his part) got angry; [and so] he went away. (lit., God, he tasted his heart, he flew his thing.)'
\z 

Finally, it is worth noting that the directional particles in Moloko do not only occur with verbs. There are also two examples where a directional particle follows an ideophone where the ideophones appears to have a predicative function, as in (\ref{ex2:molo7}).

\ea\label{ex2:molo7}
\langinfo{Moloko ideophone \textit{ʃɛŋ} with ventive marker }{}{\cite[241]{friesenGrammarMoloko2017}} \\
\gll ʃɛŋ=ala \\
\gl{ideo}=\gl{vent} \\
\glt `Going, [he] came [to the chief’s house].'
\z 

In other cases, directional particles appear in phrases with a verb, but follow pronouns or other non-predicate function words. This may be related to a discourse cohesion function of directionals \citep[241]{friesenGrammarMoloko2017} or the distinct distribution of these forms may merit a separate analysis as homophonous adverbs. 

\subparagraph{Caused accompanied motion (CAM)}

Caused accompanied motion (CAM) can be expressed by a the combination of the verb `take' with the ventive and itive directional enclitics, as in (\ref{ex2:molo10}) and (\ref{ex2:molo11}). 

\ea\label{ex2:molo10}
\langinfo{Moloko verb \textit{zaɗ} `take' with ventive marker }{}{\cite[241]{friesenGrammarMoloko2017}} \\
\gll  zaɗ=ala ɛtɛmɛ \\
take.2\gl{sg}.\gl{imp}=\gl{vent} onion	\\
\glt `Bring the onion (to me).'

\ex\label{ex2:molo11}
\langinfo{Moloko verb \textit{zaɗ} `take' with ventive marker }{}{\cite[241]{friesenGrammarMoloko2017}} \\
\gll  zaɗ=alaj ɛtɛmɛ	\\
take.2\gl{sg}.\gl{imp}=\gl{itive} onion	\\
\glt `Take the onion away (from me).'
\z 

Note that the directional enclitic does not need to be adjacent to the verb, as in (\ref{ex2:molo12}).

\ea\label{ex2:molo12}
\langinfo{Moloko verb \textit{zaɗ} `take' with ventive marker }{}{\cite[309]{friesenGrammarMoloko2017}} \\
\gll z=aw na=ala \\
carry.2\gl{sg}.\gl{imp}=1\gl{sg}.\gl{dat} 3\gl{sg}.\gl{obj}=\gl{vent} \\
\glt `Bring it to me!'
\z

\subparagraph{Buy-sell verbs}

The verb \textit{skom} can mean `buy' or `sell'	\citep[292]{friesenGrammarMoloko2017}. The verb can occur with or without a directional marker. When occurring with a ventive marker, as in (\ref{ex2:molo13}), the verb has a `buy' interpretation. When occurring with an itive marker, as in (\ref{ex2:molo14}), the verb has a `sell' interpretation. 

\ea\label{ex2:molo13}
\langinfo{Moloko verb \textit{skom} `buy/sell' with ventive marker }{}{\cite[240]{friesenGrammarMoloko2017}} \\
\gll nə̀-sʊkʷɔm=ala awak\\
1\gl{sg}.\gl{pfv}-buy/sell=\gl{vent} goat	\\
\glt `I bought a goat.'

\ex\label{ex2:molo14}
\langinfo{Moloko verb \textit{skom} `buy/sell' with itive marker }{}{\cite[240]{friesenGrammarMoloko2017}} \\
\gll nə̀-sʊkʷɔm=alaj awak\\
1\gl{sg}.\gl{pfv}-buy/sell=\gl{itive} goat	\\
\glt `I sold my goat.'
\z

\paragraph*{Muyang (muya1243)} \label{muya}

Information on Muyang grammar is from manuscripts of two grammar sketches \citep{smithMuyangVerbPhrase2002,smithDefinitenessTopicalisationTheme2003}. These two papers are by the same author but present slightly different analyses of the verbal morphology which complicates the discussion below. A more recently published dictionary was not accessible \citep{smithDictionnaireMuyangfrancaisanglaisAvec2017}.

\subparagraph{Directional verbs}~

\begin{table}[H]
	\caption{Muyang directional verbs} 
	\label{tab:muyaV}
	\begin{tabular}{lll} % add l for every additional column or remove as necessary
		\lsptoprule
		DIR & verb & gloss   \\ %table header
		\midrule
		MOTION & \textit{hər} & `jump, go, come' \\
%		ITIVE & \textit{ru} & `go' \\
%		VENT & \textit{ra} & `come' \\
		UP & \textit{cəl} & `climb, go up' \\
		DOWN & \textit{ndahàr} & `descend' \\
		\lspbottomrule
	\end{tabular}
\end{table}

The verb root \textit{hər} is variously glossed `jump', `go' or `come', and appears in the lexicon with the glosses `jump, fly, go' \citep[38]{smithMuyangProvisionalLexicon2003}. However, the examples that include this verb are always translated as `go' or `come' with the exception of one example in the lexiconː `the bird flys off into the sky'. The evidence that \textit{hər} is a manner-of-motion verb is very weak, and it is more likely that Muyang follows the pattern of several closely related languages in having a general motion verb not specified for direction (or manner). 

The lexicon also includes the entries \textit{ra} `come' and \textit{ru} `go'  \citep[76--77]{smithMuyangProvisionalLexicon2003}. The entry for \textit{ra} includes several examples sentences (not parsed or glossed) expressing ventive motion, itive motion, future tense, past habitual and another non-motion meaning. The status of these words remains unclear.

\subparagraph{Directional extensions}~

\begin{table}[H]
	\caption{Muyang directional morphology}
	\label{tab:muya}
	\begin{tabular}{llll} % add l for every additional column or remove as necessary
		\lsptoprule
		forms & source gloss & direction  & other functions  \\ %table header
		\midrule
		\textit{-biyu}  &  \textsc{home(ward), hither}  &  \textsc{vent} &   \textsc{subs.vent}, \textsc{loc} \\
		\textit{-iyu}  &  \textsc{into}  & \textsc{into}  &    \\
		\textit{-aya}  &  \textsc{out}  & \textsc{out}  &    \\
		\lspbottomrule
	\end{tabular}
\end{table} 

There are at least three verbal suffixes in Muyang which combine with motion verbs to contribute path-related information. The directional meaning of the suffixes appears most frequently in caused accompanied motion expressions (discussed below), but there are also a number of examples with motion verbs. The ventive suffix \textit{-biya} can appear with an (intransitive) motion verb, as in (\ref{ex2:muya1}). 

\ea\label{ex2:muya1}
\langinfo{Muyang verb \textit{slək} `return' with ventive suffix}{}{\cite[19]{smithDefinitenessTopicalisationTheme2003}} \\
\gll ezʉwi ga-ni am-a-sləka-biyu va do aw  \\
fly \gl{rel}-\gl{def} \gl{pot}-3\gl{sg}-return-\gl{vent} again \gl{neg} \gl{q} \\
\glt `The fly that did this will come back, won't he?'
\z 

In some examples, the ventive suffix is also shown to be able to express a subsequent ventive motion (or subsequent ventive CAM), as in (\ref{ex2:muya4}). However, note that the same verb is shown with a ventive suffix in (\ref{ex2:muya2}) where there is no explicit subsequent motion in the translation.

\ea\label{ex2:muya4}
\langinfo{Muyang verb \textit{səkum} `buy' with ventive suffix}{}{\cite[16]{smithMuyangVerbPhrase2002}} \\
\gll a-səkum-ləŋ-biyu səmu ana goro \\
3\gl{sg}-buy-behind-\gl{vent} cement to \gl{rel}-1s	\\
\glt `He buys cement to come and add to mine.'
\z 

There is also one example, shown in (\ref{ex2:muya5}), where the ventive suffix appears to have an altrilocative interpretation. It is not clear in what contexts these various interpretations of the ventive suffix arise. 

\ea\label{ex2:muya5}
\langinfo{Muyang verb \textit{gr} `make [greeting]' with ventive suffix}{}{\cite[15]{smithMuyangVerbPhrase2002}} \\
\gll a-gr-uk-biyu sa \\
3\gl{sg}-make-2\gl{sg}.\gl{dat}-\gl{vent} greeting	\\
\glt `He greets you from a distance.'
\z 

The suffix \textit{-iyu} has an inward directional function with intransitive and transitive motion verbs, as in (\ref{ex2:muya2}) and (\ref{ex2:muya3}). Notice in (\ref{ex2:muya3}) that the inward meaning is also expressed by an adposition \textit{vu}. This dual expression of directional or locational meaning through both a preposition and a verbal suffix appears to be relatively common in Muyang. 

\ea\label{ex2:muya2}
\langinfo{Muyang verb \textit{hər} `go/come' with inward suffix}{}{\cite[27]{smithDefinitenessTopicalisationTheme2003}} \\
\gll À-hər-iyu à-səkum-biyu sla ga-yaŋ, sisi \\
3\gl{sg}-go/come-\gl{into} 3\gl{sg}-buy-\gl{vent} cow \gl{rel}-3\gl{sg} five.francs	\\
\glt `He went in and bought himself a cow, for five francs.'

\ex\label{ex2:muya3}
\langinfo{Muyang verb \textit{b} `pour' with inward suffix}{}{\cite[14]{smithMuyangVerbPhrase2002}} \\
\gll  kâ-ra kə̂-b-iyu e ɗidʉ vu \\
2\gl{sg}.\gl{sbj}-come 2\gl{sg}.\gl{sbj}-pour-\gl{into} at pot into	\\
\glt `...then come and pour it into the pot.'
\z 

The suffix \textit{-aya} can have an an outward directional function with (intransitive) motion verbs, as in (\ref{ex2:muya6}). 

\ea\label{ex2:muya6}
\langinfo{Muyang verb \textit{dəg} `fall' with outward suffix}{}{\cite[14]{smithDefinitenessTopicalisationTheme2003}} \\
\gll ba hay a-dəg-aya ciliŋ, ba hay a-dəg-aya ciliŋ \\
only millet 3\gl{sg}-fall-\gl{out} only only millet 3\gl{sg}-fall-\gl{out} only \\
\glt `...and millet just poured out.'
\z 

Two examples of the suffix \textit{-aya}, including (\ref{ex2:muya10}), show several possible directional interpretations including ventive and itive, leaving it unclear how the interpretation of this suffix is determined.

\ea\label{ex2:muya10}
\langinfo{Muyang verb \textit{hər} `go/come' with outward suffix}{}{\cite[17]{smithMuyangVerbPhrase2002}} \\
\gll hər-aya	\\
go/come-\gl{out}		\\
\glt `come in, come out, come down...'
\z 

\citet{smithMuyangVerbPhrase2002} describes other verbal morphology as ``directional'' but without examples that indicate the suffix is giving information about the directionality of a figure on a path of motion. Some of the morphology can be described as locational, and in these cases the locational verbal morphology often co-occurs with a particular adposition, as in the case of \textit{-ki} `on'. Other suffixes are described as motion-related but the examples do not support the analysis. For example, the suffix \textit{-aba} is described as meaning that a ``direct object is removed from the ground'' but based on the available examples it appears to be aspectual in meaning \citep[19]{smithMuyangVerbPhrase2002}. 

\citet{smithMuyangVerbPhrase2002} also suggests that many of these verbal endings can be decomposed into single segments, e.g., ``\textit{-biyu} (= \textit{b} + \textit{y} + \textit{u})'' but the combinations of consonants generally do not break down into natural sub-components of meaning. \citet{smithMuyangVerbPhrase2002} proposes that the verb-final vowels \textit{a} and \textit{u} are related to ``separation'', however, \citet[26]{smithDefinitenessTopicalisationTheme2003} re-analyzes the final \textit{a} as a TAM marker. 

Finally, the suffix \textit{-oro} is described as a ``continuation marker'' \citep[22]{smithMuyangVerbPhrase2002}, but it does not appear to have an aspectual meaning. It is worth noting that this suffix is cognate with the Mada itive suffix \textit{-oro} and the Wuzlam itive suffix \textit{-ərə}. The examples available do not make clear whether this suffix has a directional function in Muyang, but it is possible that it is an (erstwhile) itive suffix. 

\subparagraph{Caused accompanied motion (CAM)}

Directed caused accompanied motion can be expressed by the combination of the verb \textit{z-} `take' or \textit{həl-} `gather' with a directional suffix, as in (\ref{ex2:muya8}).

\ea\label{ex2:muya8}
\langinfo{Muyang verb \textit{həl} `gather' with outward suffix}{}{\cite[4]{smithDefinitenessTopicalisationTheme2003}} \\
\gll a-həl-i-ya ana zal ga-ni na ti, wiyaŋ \\
3\gl{sg}-gather-\gl{dat}-\gl{out} for man \gl{rel}-\gl{def} \gl{def} \gl{top} sand \\
\glt `The thing that she brought out for her husband was sand.'
\z 

The suffix \textit{-oro} is also found with the verb \textit{z} `take' to express directed CAM, as in (\ref{ex2:muya9}), providing further support to the idea that it may be an itive directional suffix. 

\ea\label{ex2:muya9}
\langinfo{Muyang verb \textit{z} `take' with \textit{-oro} suffix}{}{\cite[22]{smithMuyangVerbPhrase2002}} \\
\gll a-z-oro zlam a gosku vu \\
3\gl{sg}-take-\gl{cont} thing at market into \\
\glt `He takes something away to market.'
\z 

\subparagraph{Buy-sell verbs}

The verb \textit{səkum} `buy' with the suffix \textit{-oro} becomes \textit{səkumoro} `sell' \citep[79]{smithMuyangProvisionalLexicon2003}. As mentioned above, the suffix \textit{-oro} is described as a ``continuation'' marker by \citet{smithMuyangVerbPhrase2002}, but it may be an itive suffix.

\paragraph*{Wuzlam (wuzl1236)} \label{wuzl}

Information on the grammar of Wulzam (also known as Ouldemé) is from a grammar sketch and lexicon \citep{colombelLangueOuldemeNordCameroun2005}.

\subparagraph{Directional verbs}~

\begin{table}[H]
	\caption{Wuzlam directional verbs} 
	\label{tab:wuzlV}
	\begin{tabular}{lll} % add l for every additional column or remove as necessary
		\lsptoprule
		DIR & verb & gloss   \\ %table header
		\midrule
		MOTION & \textit{w} & `\textit{aller}, \textit{venir} [go, come]' \\
		UP & \textit{cəl} & `\textit{monter} [ascend]' \\
		\lspbottomrule
	\end{tabular}
\end{table}

\subparagraph{Directional extensions}~

\begin{table}[H]
	\caption{Wuzlam directional morphology}
	\label{tab:wuzl}
	\begin{tabular}{llll} % add l for every additional column or remove as necessary
		\lsptoprule
		forms & source gloss & direction  & other functions  \\ %table header
		\midrule
		\textit{-ārá}  &  \textit{centrip.}  &  \textsc{vent}  &    \\
		\textit{-ārāy, ə̄rə̄}  &  \textit{centrif.}  & \textsc{itive}   &    \\
		\lspbottomrule
	\end{tabular}
\end{table} 

Wuzlam (or Ouldemé) has a ventive suffix and an itive suffix. Examples of these suffixes expressing ventive and itive directionality are mostly with the verb root \textit{w-} which is a generic motion verb that can mean either `go' or `come' depending on the suffix, as in (\ref{ex2:wuzl1}) and (\ref{ex2:wuzl2}). These suffixes are not found with any transitive motion verbs. When occurring with non-motion verbs, there was no indication of their function, other than in the cases of caused accompanied motion, discussed below. 

\ea\label{ex2:wuzl1}
\langinfo{Wuzlam verb \textit{w} `go/come' with itive suffix}{}{\cite[20]{colombelLangueOuldemeNordCameroun2005}} \\
\gll t-ə̄-w-ə̄rə̄ á yàm àgè \\
they-\gl{real}-go/come-\gl{itive} \gl{prep} water in \\
\glt `They go to look for water.' [French: `\textit{Ils vont chercher de l'eau.}']

\ex\label{ex2:wuzl2}
\langinfo{Wuzlam verb \textit{w} `go/come' with ventive suffix}{}{\cite[57]{colombelLangueOuldemeNordCameroun2005}} \\
\gll ā-w-ārá ánàwə́ \\
he.\gl{real}-go/come-\gl{vent} yesterday \\
\glt `He came yesterday.' [French: `\textit{Il est venu hier.}']
\z	

Several other verbal suffixes are described as being of the ``\textit{cadre spatio-temporel}'', for example, the suffix \textit{-ākā} is described as meaning ``\textit{destination avec lieu où rester, où laisser [destination with a place to rest, or leave]}'' \citep[54]{colombelLangueOuldemeNordCameroun2005} but this suffix is not found in examples expressing path of motion, but rather in static locational contexts. 

\subparagraph{Caused accompanied motion (CAM)}

The ventive and itive suffix can  be used with a verb of transfer of possession to express directed caused accompanied motion, as in (\ref{ex2:wuzl4}). 

\ea\label{ex2:wuzl4}
\langinfo{Wuzlam verb \textit{zàk} `take' with itive suffix}{}{\cite[40]{colombelLangueOuldemeNordCameroun2005}} \\
\gll n-ə́-zàk(w)à-àkw-ø̄-ārāy \\
1\gl{sg}-\gl{real}-take-2\gl{sg}.\gl{dat}-3\gl{obj}-\gl{itive} \\
\glt `I will bring it to you.' [French: `\textit{Je te le rapporterai.}']
\z 

\subparagraph{Buy-sell verbs}

The verb \textit{sə̄kə̄m} `buy' with the itive suffix becomes \textit{sə̄kə̄m-arə} `sell' \citep[162]{colombelLangueOuldemeNordCameroun2005}. 

\subsection{Maroua (1/3)}

Only one language from the Maruou group, North Giziga, has been covered by any detailed morphosyntactic analysis. No description is available for South Giziga (sout3051) or Baldemu (bald1241). Note that there is some debate whether North and South Giziga are best classified as separate languages or as dialects of the same language \citep[2--3]{shayGrammarGizigaChadic2021}.  

\paragraph*{North Giziga (nort3047)} \label{gizi}

Information on the grammar of North Giziga is from a reference grammar \citep{shayGrammarGizigaChadic2021}.

\subparagraph{Directional verbs}~

\begin{table}[H]
	\caption{North Giziga directional verbs} 
	\label{tab:giziV}
	\begin{tabular}{lll} % add l for every additional column or remove as necessary
		\lsptoprule
		DIR & verb & gloss   \\ %table header
		\midrule
		ITIVE & \textit{r-} & `go' \\
		VENT & \textit{s-áwà} & `come' \\
		INTO& \textit{cí} & `enter' \\
		OUT &\textit{b-} & `get out' \\
		\lspbottomrule
	\end{tabular}
\end{table}

\subparagraph{Directional extensions}~

\begin{table}[H]
	\caption{North Giziga directional morphology}
	\label{tab:gizi}
	\begin{tabular}{llll} % add l for every additional column or remove as necessary
		\lsptoprule
		forms & source gloss & direction  & other functions  \\ %table header
		\midrule
		\textit{-áwà, -óò, -o}  &  \textsc{vent}  &  \textsc{vent}  &  \textsc{ben}  \\
		\lspbottomrule
	\end{tabular}
\end{table} 

The ventive suffix in Giziga is shown to have a directional function with intransitive motion motion verbs, as in (\ref{ex2:gizi1}). No examples were found with transitive motion verbs. 

\ea\label{ex2:gizi1}
\langinfo{Giziga verb \textit{kl} `run' with ventive suffix}{}{\cite[205]{shayGrammarGizigaChadic2021}} \\
\gll kl-ó \\
run-\gl{vent} \\
\glt `Come running (over here)!'
\z 

The ventive suffix also appears with non-motion verbs, but in these cases there is no consistent indication of what the function of the ventive suffix may be. In (\ref{ex2:gizi2}) there appears to be a benefactive interpretation. 

\ea\label{ex2:gizi2}
\langinfo{Giziga prior associated motion SVC}{}{\cite[216]{shayGrammarGizigaChadic2021}} \\
\gll à=r(ú) à=ɬ-ó hùmàcà \\
3=go 3=cut-\gl{vent} hay \\
\glt `He goes and cuts hay for himself.'
\z 

\citet[203]{shayGrammarGizigaChadic2021} notes that although the verb root \textit{s} `come' never appears without the ventive suffix, this is not completely redundant as it serves to distinguish the verb from the future tense marker \textit{sà}, presumably derived from the motion verb. 

\subparagraph{Caused accompanied motion (CAM)}

Directed caused accompanied motion can be expressed by the verb \textit{zɓ} `take' with a ventive directional suffix or by the use of an associative prepositional phrase with a motion verb, as in (\ref{ex2:gizi3}). However, \citet[200]{shayGrammarGizigaChadic2021} notes that when the item being caused to move is pronominalized, it is no longer preceded by a preposition, but rather appears as a direct object. 

\ea\label{ex2:gizi3}
\langinfo{Giziga verb \textit{r} `go' with associative prepositional phrase}{}{\cite[71]{shayGrammarGizigaChadic2021}} \\
\gll  '1= r-am ti tapa ti aw(d) \\
1= go-\gl{pl} \gl{asoc} tobacco \gl{asoc} goat \\
\glt `We brought (lit. ‘went with’) tobacco (and) a goat.'
\z 

\subparagraph{Buy-sell verbs}

The verb \textit{hìɗík} can mean `trade, sell, buy', but in its ventive form, \textit{hìɗk-ò}, it explicitly means `buy' \citep[150]{shayGrammarGizigaChadic2021}. 

\subsection{Musguic (2/3)}

Two of the three Musguic languages are included here. The third language, Muskum (musk1256), is reportedly dormant and without any description of its morphosyntax \citep{tourneuxLangueTchadiqueDisparue1977}.

\paragraph*{Mbara (mbar1260)} \label{mbar}

Information on the grammar of Mbara is from a book containing a grammar sketch, lexicon and texts \citep{tourneuxMbaraLeurLangue1986}.

\newpage
\subparagraph{Directional verbs}~

\begin{table}[H]
	\caption{Mbara directional verbs} 
	\label{tab:mbarV}
	\begin{tabular}{lll} % add l for every additional column or remove as necessary
		\lsptoprule
		DIR & verb & gloss   \\ %table header
		\midrule
		MOTION & \textit{hílí, kál} & `go' \\
		OUT & \textit{dày} & `go out' \\
		UP & \textit{rìː} & `go up' \\
		\lspbottomrule
	\end{tabular}
\end{table}

The lexicon in \citet{tourneuxMbaraLeurLangue1986} provides an imperative form of a ventive motion verb, \textit{ʔàgú-sì} `come!', but in texts ventive motion is expressed by the verb \textit{cóː} `arrive' or \textit{hílí} `go' with the ventive suffix.

\subparagraph{Directional extensions}~


\begin{table}[H]
	\caption{Mbara directional morphology}
	\label{tab:mbar}
	\begin{tabular}{llll} % add l for every additional column or remove as necessary
		\lsptoprule
		forms & source gloss & direction  & other functions  \\ %table header
		\midrule
		\textit{-sì}  &  \textsc{direct}  &  \textsc{vent} &     \\
		\lspbottomrule
	\end{tabular}
\end{table} 

\citet[189]{tourneuxMbaraLeurLangue1986} describe a ventive suffix \textit{-sì} in Mbara ``\textit{pour indiquer que le procès se fait ou doit se faire en direction de celui qui parle} (\textit{directionnel centripète}) [to indicate that the process is done or should be done in the direction of the one who speaks (centripetal directional)].'' The ventive suffix is shown with both intranstive and transitive motion verbs, as in (\ref{ex2:mbar1}) and (\ref{ex2:mbar2}), but not found in any examples with non-motion verbs. 

\ea\label{ex2:mbar1}
\langinfo{Mbara verb \textit{cóː} `arrive' with ventive suffix}{}{\cite[190]{tourneuxMbaraLeurLangue1986}} \\
\gll kì cóː-kí-sì \\
2\gl{pl}.\gl{pfv} arrive-2\gl{pl}-\gl{vent} \\
\glt `You came.' [French: `\textit{Vous êtes venus}.']

\ex\label{ex2:mbar2}
\langinfo{Mbara verb \textit{hà} `give' with ventive suffix}{}{\cite[190]{tourneuxMbaraLeurLangue1986}} \\
\gll hà-ú-sì mùgù'dòy \\
give-me-\gl{vent} knife \\
\glt `Give me the knife!' [French: `\textit{Donne-moi le couteau} !']
\z 

\subparagraph{Caused accompanied motion (CAM)}

There are few examples of caused accompanied motion available, however, in (\ref{ex2:mbar3}), the verb \textit{tùrkú} `remove' with a ventive suffix appears to express ventive caused accompanied motion. The verb \textit{yàm} `take' is not shown with the ventive suffix.

\ea\label{ex2:mbar3}
\langinfo{Mbara verb \textit{tùrkú} `remove' with ventive suffix}{}{\cite[190]{tourneuxMbaraLeurLangue1986}} \\
\gll ì tùrkú-sì mòmóy-kà	\\
3\gl{pl}.\gl{pfv} remove-\gl{vent} honey-\gl{pfv} \\
\glt `They came back with their harvest of honey.' \newline [French: `\textit{Ils sont revenus avec leur récolte de miel}.']
\z 


\subparagraph{Buy-sell verbs}

\citet{tourneuxMbaraLeurLangue1986} list the verb \textit{lùf} `buy', but do not list any forms meaning `sell'. 

\paragraph*{Musgu (musg1254)} \label{musg}

Information on the grammar of Musgu is from a reference grammar \citep{meyer-bahlburgStudienZurMorphologie1972} and a lexicon \citep{tourneuxLexiquePratiqueMunjuk1991}. Note that example sentences in the reference grammar are not glossed, so glossing is added to examples below where possible.

\subparagraph{Directional verbs}~

\begin{table}[H]
	\caption{Musgu directional verbs} 
	\label{tab:musgV}
	\begin{tabular}{lll} % add l for every additional column or remove as necessary
		\lsptoprule
		DIR & verb & gloss   \\ %table header
		\midrule
		MOTION & \textit{hili, hala} & `\textit{aller} [go]' \\
		INTO/OUT & \textit{niyi} & `\textit{entrer ou sortir} [enter or exit]' \\
		\lspbottomrule
	\end{tabular}
\end{table}

The verbs in Table \ref{tab:musgV} are from \citet{tourneuxLexiquePratiqueMunjuk1991}. The verb \textit{hili, hala} is glossed `go' but it can occur with the ventive directional suffix indicating that it is not an itive directional verb. 

\subparagraph{Directional extensions}~

\begin{table}[H]
	\caption{Musgu directional morphology}
	\label{tab:musg}
	\begin{tabular}{llll} % add l for every additional column or remove as necessary
		\lsptoprule
		forms &  gloss & direction  & other functions  \\ %table header
		\midrule
		\textit{sí}  &  \textsc{vent}  &  \textsc{vent} &     \\
		\textit{di}  &  \textsc{itive}  & \textsc{itive}   &    \\
		\lspbottomrule
	\end{tabular}
\end{table} 

\citet{meyer-bahlburgStudienZurMorphologie1972} describes two directional extensions. The ventive extension \textit{sí}, labeled \textit{Entfernungsmorphem} [distance morpheme], has ventive directional meaning with motion verbs, as in (\ref{ex2:musg1}) and (\ref{ex2:musg5}).

\ea\label{ex2:musg1}
\langinfo{Musgu verb \textit{nìyà} `enter/exit' with ventive suffix}{}{\cite[132]{meyer-bahlburgStudienZurMorphologie1972}} \\
\gll  nìyà s àmláy! \\
enter/exit \gl{vent} ?? \\
\glt `Come in!' [German: `\textit{Tritt herein}!']

\ex\label{ex2:musg5}
\langinfo{Musgu verb \textit{hili} `go/come' with ventive suffix}{}{\cite[132]{meyer-bahlburgStudienZurMorphologie1972}} \\
\gll hili si \\
go/come \gl{vent}	\\
\glt `come' [German: `\textit{Kommen}']
\z 

In the few cases of the ventive suffix with non-motion verbs the function is not clear. \citet[133]{meyer-bahlburgStudienZurMorphologie1972} discusses the use of the ventive suffix with the verb \textit{na} `to be' to ``express distance'' [German: `\textit{zum Ausdruck der Entfernung}'], but it is not obvious from the examples what this means. There is one example of a resultative motion interpretation, shown in (\ref{ex2:musg2}).

\ea\label{ex2:musg2}
\langinfo{Musgu verb \textit{yika} `call' with ventive suffix}{}{\cite[132]{meyer-bahlburgStudienZurMorphologie1972}} \\
\gll  a yika si arogwi n agwi si-ni	\\
?? call \gl{vent} ?? ?? ?? ??-?? \\
\glt `He called a little child to him.' [German: `\textit{Er rief ein kleines Kind zu sich.}']
\z 

\citet[133--134]{meyer-bahlburgStudienZurMorphologie1972} mentions another extension \textit{di} which appears to have an itive function, at least with intranstive motion verbs, as in (\ref{ex2:musg3}). They note that this extension is likely derived from the adverbial \textit{edi} `out' [German: `\textit{hinaus}'].

\ea\label{ex2:musg3}
\langinfo{Musgu verb \textit{hala} `go/come' with itive suffix}{}{\cite[134]{meyer-bahlburgStudienZurMorphologie1972}} \\
\gll Yesu a hala di enye \\
Jesus ?? go/come \gl{itive} ?? \\
\glt `Jesus went away.' [German: `\textit{Jesus ging fort}.']
\z 

\subparagraph{Caused accompanied motion (CAM)}

The ventive suffix with the verb \textit{ká} `take' expresses ventive caused accompanied motion, as in (\ref{ex2:musg4}). 

\ea\label{ex2:musg4}
\langinfo{Musgu verb \textit{ká} `take' with ventive suffix}{}{\cite[132]{meyer-bahlburgStudienZurMorphologie1972}} \\
\gll ká sí wàlá símí ! \\
take \gl{vent} ?? ??	\\
\glt `Bring me something to eat!' [German: `\textit{Bring mir zu essen}!']
\z 

\subparagraph{Buy-sell verbs}

There are two verbs meaning `buy', \textit{livi} and \textit{wiri}. With the adverbial \textit{àdí}, the expression is translated `sell', as in \textit{lìvì ... àdí} `sell' \citep[240]{tourneuxMulwiOuVulum1978} and \textit{wiri adi} `sell' \citep[124]{tourneuxLexiquePratiqueMunjuk1991}. The form \textit{àdí} is glossed as `\textit{en brousse, hors du village};\textit{ vers l'extérieur} [in the bush, out of the village; toward the outside]' \citep[71]{tourneuxLexiquePratiqueMunjuk1991}. The adverbial \textit{àdí} is said to be the source that the itive verbal extension \textit{di} is derived from \citep[133--134]{meyer-bahlburgStudienZurMorphologie1972}.

\section{South Biu-Mandara (13/27)}

The 27 South Biu-Mandara languages are further divided into five groups. For the largest group, Bataic, descriptions are only available for two of 12 languages. The Dabaic group stands out in having descriptions available for all six languages in the group. 

\subsection{Bataic (2/12)}

The 12 Bataic languages are mostly spoken in an area that has been difficult for foreigners to access in the eastern part of Nigeria along the border with Cameroon, so it is not surprising that morphosyntactic descriptions are only available for two of them: Bacama and Gude. There is also a description of Nzanyi (nzan1240), but its analysis of directionals is too incomplete to be included here \citep[44--46]{bensonNzanyiGrammarWriteup2014}. Nzanyi is one of eight languages in a subgroup with Gude, the other six being: Fali (fali1285), Jimi (jimi1254), Zizilivakan (zizi1238), Sharwa (shar1249), Tsuvant (tsuv1243), and one extinct language, Holma (holm1250). The other Bataic languages are
Bara (bata1314),
Gudu (gudu1250) and
Ngwaba (ngwa\-1251). Note that Bacama-Yimbiru (baca1245) was reclassified as a dialect of Bacama in Glottolog 5.0.

\paragraph*{Bacama (baca1246)} \label{baca}

Information on the grammar of Bacama is from a grammar sketch published as a journal article \citep{carnochanCategoriesVerbalPiece1970} and a dictionary \citep{pweddonBwatiyeEnglishDictionary2001}. Note that the grammar sketch does not include interlinear glossing so glosses have been added to examples below where possible. 

\subparagraph{Directional verbs}~

\begin{table}[H]
	\caption{Bacama directional verbs} 
	\label{tab:bacaV}
	\begin{tabular}{lll} % add l for every additional column or remove as necessary
		\lsptoprule
		DIR & verb & gloss   \\ %table header
		\midrule
		ITIVE & \textit{wudo, zo} & `go' \\
		VENT & \textit{ʃi} & `come' \\
		INTO& \textit{púr} & `enter, pass by' \\
		OUT & \textit{düm-} & `exit' \\
		DOWN & \textit{jang-} & `descend' \\		
		\lspbottomrule
	\end{tabular}
\end{table}

Verbs in Table \ref{tab:bacaV} are from \citet{pweddonBwatiyeEnglishDictionary2001}.

\subparagraph{Directional extensions}~

\begin{table}[H]
	\caption{Bacama directional morphology}
	\label{tab:baca}
	\begin{tabular}{llll} % add l for every additional column or remove as necessary
		\lsptoprule
		forms & source gloss & direction  & other functions  \\ %table header
		\midrule
		\textit{-á}  &  \textsc{adessive}  &  \textsc{vent} &   \textsc{subs.vent}, \textsc{prior.itive}  \\
		\lspbottomrule
	\end{tabular}
\end{table} 

\citet{carnochanCategoriesVerbalPiece1970} describes an ``adessive'' suffix \textit{-á} which has a ventive directional meaning with intranstive and transitive motion verbs, as in (\ref{ex2:baca1}) and (\ref{ex2:baca2}).

\ea\label{ex2:baca1}
\langinfo{Bacama verb \textit{dùm} `exit' with ventive suffix}{}{\cite[93]{carnochanCategoriesVerbalPiece1970}} \\
\gll Hómon a dùm-a \\
chief ?? exit-\gl{vent}	\\
\glt `The chief came out.'

\ex\label{ex2:baca2}
\langinfo{Bacama verb \textit{ŋgàl} `pull' with ventive suffix}{}{\cite[91]{carnochanCategoriesVerbalPiece1970}} \\
\gll  Ndā ŋgàl-áa-dò nzēy sàlakēy \\
3\gl{sg}.\gl{m} pull-\gl{vent}-\gl{caus} boy rope	\\
\glt `He made the boy pull the rope towards (him).'
\z

With non-motion verbs, there are several possible interpretations of the ventive suffix, without any discussion of the contexts in which the various meanings are found. One possible meaning is subsequent ventive motion, as in (\ref{ex2:baca3}). 

\ea\label{ex2:baca3}
\langinfo{Bacama verb \textit{ɗàw} `cut' with ventive suffix}{}{\cite[92]{carnochanCategoriesVerbalPiece1970}} \\
\gll Ndā ɗàw-á Victòr kadā	\\
3\gl{sg}.\gl{m} cut-\gl{vent} Victor tree		\\
\glt `He cut down the tree for Victor and returned.'
\z 

However, in another example of the ventive suffix with the same verb root, shown in (\ref{ex2:baca4}), the interpretation is prior itive motion. The main difference between the two examples is the use of the ``deprivative'' suffix in the second.

\ea\label{ex2:baca4}
\langinfo{Bacama verb \textit{ɗàw} `cut' with ventive suffix}{}{\cite[93]{carnochanCategoriesVerbalPiece1970}} \\
\gll Ndā ɗàw-á-go Pwèddon kadā	\\
3\gl{sg}.\gl{m} cut-\gl{vent}-\textsc{deprivative} Pweddon tree		\\
\glt `He went and cut down the tree without Pweddon's knowledge.'
\z 

Another possible interpretation of the ventive suffix with non-motion verbs is subsequent caused accompanied motion, as in (\ref{ex2:baca5}). 

\ea\label{ex2:baca5}
\langinfo{Bacama verb \textit{ɓay} `break' with ventive suffix}{}{\cite[102]{carnochanCategoriesVerbalPiece1970}} \\
\gll Hūn ɓay-a kàdē	\\
1\gl{sg} break-\gl{vent} sticks		\\
\glt `I have broken the sticks and brought them.'
\z 

Other examples of the ventive suffix with non-motion verbs do not have a motion-related interpretation, as in (\ref{ex2:baca6}), which uses the same verb root as (\ref{ex2:baca5}) but appears to have an interpretation related to intentionality. 

\ea\label{ex2:baca6}
\langinfo{Bacama verb \textit{ɓiy} `break' with ventive suffix}{}{\cite[102]{carnochanCategoriesVerbalPiece1970}} \\
\gll Hūn ɓiy-a kàdā	\\
1\gl{sg} break-\gl{vent} stick		\\
\glt `I have broken the stick on purpose.'
\z 

\subparagraph{Caused accompanied motion (CAM)}

Directed caused accompanied motion can be expressed through the causative form of motion verbs, as in (\ref{ex2:baca7}) and (\ref{ex2:baca8}).

\ea\label{ex2:baca7}
\langinfo{Bacama verb \textit{dùm} `exit' with causative suffix}{}{\cite[91]{carnochanCategoriesVerbalPiece1970}} \\
\gll Ndá dùm-do nzey	\\
3\gl{sg}.\gl{m} exit-\gl{caus} boy		\\
\glt `He caused the boy to go out. / He took the boy out.'

\ex\label{ex2:baca8}
\langinfo{Bacama verb \textit{dùm} `exit' with causative suffix}{}{\cite[91]{carnochanCategoriesVerbalPiece1970}} \\
\gll Nda dum-d-i	\\
3\gl{sg}.\gl{m} exit-\gl{caus}-1\gl{sg}		\\
\glt `He took me out.'
\z 

\subparagraph{Buy-sell verbs}

The verb \textit{ɗörö} `buy' in its causative form is \textit{ɗör-d-} `sell' \citep[162]{pweddonBwatiyeEnglishDictionary2001}.

\paragraph*{Gude (gude1246)} \label{gude}

Gude is described by \citet{hoskinsonGrammarDictionaryGude1983}. This work does not include interlinearized glossing making it difficult to find examples of the use of motion-related suffixes outside of the sections where they are explicitly discussed in the grammar sketch. \citet[353–355, 359–361]{schuhChadicCornucopia2017} gives some additional examples of Gude verbal morphology based on elicitation. 

\subparagraph{Directional verbs}~

\begin{table}[H]
	\caption{Gude directional verbs} 
	\label{tab:gude}
	\begin{tabular}{lll} % add l for every additional column or remove as necessary
		\lsptoprule
		DIR & verb & gloss   \\ %table header
		\midrule
		ITIVE & \textit{dzə} & `go' \\
		VENT & \textit{shi} & `come' \\
		INTO/OUT & \textit{dəmə} & `go in or out' \\
		DOWN & \textit{jim} & `descend' \\
		\lspbottomrule
	\end{tabular}
\end{table}

\subparagraph{Directional extensions}~

\begin{table}[H]
	\caption{Gude directional morphology}
	\label{tab:gudeV}
	\begin{tabular}{llll} % add l for every additional column or remove as necessary
		\lsptoprule
		forms & source gloss & direction  & other functions  \\ %table header
		\midrule
		\textit{-ʸà}  &   & \textsc{vent}   &   \textsc{subs.vent}, lexicalized  \\
		\textit{-gərə}  &    &   \textsc{into/down} &   \\
		\textit{-gi}  &   &   \textsc{out/up}  &  \textsc{compl}   \\
		\textit{-paa}  &   & \textsc{down}   &  \textsc{subs.down}, lexicalized  \\
		%\textit{-və}  &  \textsc{allative}  &   lexicalized &     \\
		\lspbottomrule
	\end{tabular}
\end{table} 

There are three primary directional suffixes in Gude, and one suffix \textit{-paa} which can have a directional interpretation with particular verbs.\footnote{In addition, \citet[99]{hoskinsonGrammarDictionaryGude1983} describes the suffix \textit{-və} as ``motion to a specific place'' but in no examples does this suffix add information about the direction of the path of motion \citep[cf.][360--361]{schuhChadicCornucopia2017}.} The ventive suffix is interpreted as a directional with motion verbs, as in (\ref{ex2:gude1}).\footnote{One notable exception is the example \textit{hwa-ya} `run-\gl{vent}' translated as `grow horizontally (of vine)' \citep[203]{hoskinsonGrammarDictionaryGude1983}.} 

\ea\label{ex2:gude1}
\langinfo{Gude verb \textit{gim} `enter/exit' with ventive suffix }{}{\cite[100]{hoskinsonGrammarDictionaryGude1983}} \\
\gll gim-a	\\
enter/exit-\gl{vent}		\\
\glt `come in or out (toward speaker)'
\z 

However, with the motion verb \textit{kuly} `fall', the translation in (\ref{ex2:gude2}) suggests a locative interpretation of the ventive suffix.

\ea\label{ex2:gude2}
\langinfo{Gude verb \textit{kuly} `fall' with ventive suffix }{}{\cite[98]{hoskinsonGrammarDictionaryGude1983}} \\
\gll kuly-a \\
fall-\gl{vent} \\
\glt `fall here'
\z 

With non-motion verbs, the ventive can have a subsequent ventive motion interpretation, as in (\ref{ex2:gude4}) and (\ref{ex2:gude5}). The example in (\ref{ex2:gude5}) is the one case where the interpretation includes caused accompanied motion.

\ea\label{ex2:gude4}
\langinfo{Gude verb \textit{ly} `cut' with ventive suffix }{}{\cite[98]{hoskinsonGrammarDictionaryGude1983}} \\
\gll ly-a \\
cut-\gl{vent} \\
\glt `cut and come'

\ex\label{ex2:gude5}
\langinfo{Gude verb \textit{yir} `buy' with ventive suffix }{}{\cite[354]{schuhChadicCornucopia2017}} \\
\gll yir-a-n \\
buy-\gl{vent}-\gl{inf} \\
\glt `buy something and bring it'
\z

In (\ref{ex2:gude3}) the ventive meaning most likely refers to the path of motion of the projectile being shot, which is not shown to be an overt argument of the verb.

\ea\label{ex2:gude3}
\langinfo{Gude verb \textit{hy} `shoot' with ventive suffix }{}{\cite[98]{hoskinsonGrammarDictionaryGude1983}} \\
\gll hy-a	\\
shoot-\gl{vent}	\\
\glt `shoot here'
\z 

The suffix \textit{-gərə} can express either inward or downward directional motion when combined with a motion verb, as in (\ref{ex2:gude6}) and (\ref{ex2:gude7}), and its counterpart \textit{-gi} can express outward or upward directional motion, as in (\ref{ex2:gude8}) and (\ref{ex2:gude9}). It is not clear whether the interpretation is determined by context or other factors such as the semantics of the main verb.

\ea\label{ex2:gude6}
\langinfo{Gude verb \textit{dəmə} `enter/exit' with inward/downward suffix }{}{\cite[100]{hoskinsonGrammarDictionaryGude1983}} \\
\gll dəmə-gərə	\\
enter/exit-\gl{into}/\gl{down} \\
\glt `go in'

\ex\label{ex2:gude7}
\langinfo{Gude verb \textit{dzə} `go' with inward/downward suffix }{}{\cite[25]{hoskinsonGrammarDictionaryGude1983}} \\
\gll dzə-gərə	\\
go-\gl{into}/\gl{down}	\\
\glt `go down'

\ex\label{ex2:gude8}
\langinfo{Gude verb \textit{dəmə} `enter/exit' with outward/upward suffix }{}{\cite[100]{hoskinsonGrammarDictionaryGude1983}} \\
\gll dəmə-gi \\
enter/exit-\gl{out}/\gl{up}	\\
\glt `go out'

\ex\label{ex2:gude9}
\langinfo{Gude verb \textit{ndərə} `climb' with outward/upward suffix }{}{\cite[25]{hoskinsonGrammarDictionaryGude1983}} \\
\gll ndərə-gi	\\
climb-\gl{out}/\gl{up} \\
\glt `climb up'
\z 

The data contains just one example of each of these suffixes where the patient is a figure on the path of motion. In (\ref{ex2:gude10}), the inward/downward suffix \textit{-gərə} has a downward interpretation. 

\ea\label{ex2:gude10}
\langinfo{Gude verb \textit{ka} `throw' with inward/downward suffix }{}{\cite[92]{hoskinsonGrammarDictionaryGude1983}} \\
\gll ka-gərə	\\
throw-\gl{down}/\gl{into}	\\	
\glt `throw down'
\z 

In (\ref{ex2:gude11}), the outward/upward suffix \textit{-gi} has an unexpected itive translation, which is more likely an example of the ``homophonous'' totality (``do completely'') suffix \citep[101]{hoskinsonGrammarDictionaryGude1983}, as also seen in (\ref{ex2:gude12}). 

\ea\label{ex2:gude11}
\langinfo{Gude verb \textit{ka} `throw' with outward/upward suffix }{}{\cite[92]{hoskinsonGrammarDictionaryGude1983}} \\
\gll ka-gi \\
throw-\gl{up}/\gl{out}	\\
\glt `throw away'

\ex\label{ex2:gude12}
\langinfo{Gude verb \textit{adə} `eat' with outward/upward suffix }{}{\cite[92]{hoskinsonGrammarDictionaryGude1983}} \\
\gll adə-gi	\\
eat-\gl{up}/\gl{out} \\
\glt `eat (it) up'
\z 

With other non-motion verbs, these suffixes do not have motion interpretations, as seen in (\ref{ex2:gude13}).

\ea\label{ex2:gude13}
\langinfo{Gude verb \textit{tsa} `tear' with outward/upward suffix }{}{\cite[360]{schuhChadicCornucopia2017}} \\
\gll tsa-gə̀r\\
tear/split-\gl{down}/\gl{into}	\\
\glt `slit, cut a slit'
\z 

The other motion-related suffix in Gude, \textit{-paa} ``is not productive in modern Gude and occurs only with a few words'' \citep[101]{hoskinsonGrammarDictionaryGude1983}. Hoskinson describes the meaning of the suffix as ``do completely''. \citet[359]{schuhChadicCornucopia2017} notes that historically the suffix ``derives from the locative adverb \textit{pà} ‘at the base of, down, under’ and a few of the verbs that bear this extension show a ``downness'' relation.'' In most cases, the suffix does not have a motion meaning, but two examples indicate that it can express downward motion of the patient argument, as in (\ref{ex2:gude14}). In one other example, shown in (\ref{ex2:gude15}), the translation indicates that there is also subsequent motion (of the patient).

\ea\label{ex2:gude14}
\langinfo{Gude verb \textit{uurə} `throw' with \textit{-paa} suffix }{}{\cite[102]{hoskinsonGrammarDictionaryGude1983}} \\
\gll uurə-paa	\\
throw-\gl{down} \\
\glt `throw down'

\ex\label{ex2:gude15}
\langinfo{Gude verb \textit{'anyə} `tie' with \textit{-paa} suffix }{}{\cite[362]{schuhChadicCornucopia2017}} \\
\gll  nyi 'anyə-pà \\
I tie-\gl{down}	\\
\glt `I tied (it) and laid (it) on the ground.'
\z 

The ventive is the only motion-related suffix that can productively co-occur with any other motion-related suffix \citep[95]{hoskinsonGrammarDictionaryGude1983}, as in (\ref{ex2:gude16}) and (\ref{ex2:gude17}).

\ea\label{ex2:gude16}
\langinfo{Gude verb \textit{ndir} `climb' with ventive and inward/downward suffix }{}{\cite[100]{hoskinsonGrammarDictionaryGude1983}} \\
\gll ndir-a-gərə	\\
climb-\gl{vent}-\gl{down}/\gl{into}	\\
\glt `climb down (toward speaker)'

\ex\label{ex2:gude17}
\langinfo{Gude verb \textit{ndir} `climb' with ventive and outward/upward suffix }{}{\cite[100]{hoskinsonGrammarDictionaryGude1983}} \\
\gll ndir-a-gi \\
climb-\gl{vent}-\gl{up}/\gl{out}	\\
\glt `climb up (toward speaker)'

%gim-a-gərə	enter/exit-\gl{vent}-\gl{down}/\gl{into}		come in (toward speaker)	hoskinsonGrammarDictionaryGude1983	100
%gim-a-gi	enter/exit-\gl{vent}-\gl{up}/\gl{out}		come out (toward speaker)	hoskinsonGrammarDictionaryGude1983	100
\z 

Where other motion-related suffixes co-occur the combination is restricted to a particular verb stem \citep[364]{schuhChadicCornucopia2017}.  Also note that \citet[359--360]{schuhChadicCornucopia2017} proposes various adverbial locatives as the etymological source of all of the motion-related suffixes except for the ventive.

\subparagraph{Caused accompanied motion (CAM)}

The verb \textit{kərə} is defined as `carry away, take away' but with a ventive suffix, \textit{kira}, it means `bring' \citep[209]{hoskinsonGrammarDictionaryGude1983}.

\subparagraph{Buy-sell verbs}

The verb \textit{ɗərə} `buy, trade' means `sell' when used with an upward/outward suffix: \textit{ɗərəgi} \citep[182]{hoskinsonGrammarDictionaryGude1983}. Note that the ventive suffix with the same verb has a subsequent motion interpretation, as in (\ref{ex2:gude5}) above. In that example, the verb is also interpreted as `buy', but it is not clear if that is due to the ventive suffix or other context.

\subsection{Dabaic (6/6)}

Morphosyntactic descriptions are available for all six Dabaic languages. Within the Daba group, Buwal and Gavar are closely related, as are Daba and Mazagway. All six languages have a ventive suffix which can be reconstructed as \textit{*aha}. In all languages the ventive suffix (or suffixes) have multiple functions. In general, with motion verbs the suffix is a ventive directional, and with non-motion verbs the suffix can be analyzed as a subsequent motion marker. However, in Mina it can function as an altrilocative marker, and in Mbudum as a marker of prior itive motion. Other possible uses of ventive suffixes include temporal-aspectual and lexicalized uses. In addition to the \textit{*aha} suffix that can express ventive directional meaning, Daba and Buwal also have an itive marker, not found in the other Dabaic languages.

\paragraph*{Buwal (buwa1243)} \label{buwa}

The grammar of Buwal is described in detail in a doctoral dissertation \citep{viljoenGrammaticalDescriptionBuwal2013}. 

\subparagraph{Directional verbs}~

\begin{table}[H]
	\caption{Buwal directional verbs} 
	\label{tab:buwaV}
	\begin{tabular}{lll} % add l for every additional column or remove as necessary
		\lsptoprule
		DIR & verb & gloss   \\ %table header
		\midrule
		MOTION & \textit{nda} & `go, come' \\
		UP & \textit{tèv} & `ascend' \\
		DOWN & \textit{tàdàkʷ} & `descend' \\
		INTO/OUT & \textit{dàm} & `enter, exit' \\
		\lspbottomrule
	\end{tabular}
\end{table}

The verb \textit{nda} can mean `come' or `go' as seen in (\ref{ex2:buwa1}). Likewise, \textit{dàm} can mean both `enter' and `exit'.

\ea\label{ex2:buwa1}
\langinfo{Buwal verb \textit{nda} `go/come' }{}{\cite[365]{viljoenGrammaticalDescriptionBuwal2013}} \\
\gll  màdā māwàl ká-ndā āzá nèné-ná-ndā á eɡljz \\
if husband \gl{pfv}-go/come \gl{compl} 1\gl{excl}.\gl{sbj}-\gl{fut}-go/come \gl{prep} church(Fr.) \\
\glt `If my husband has come, we will go to church.'
\z

\subparagraph{Directional extensions}~

\begin{table}[H]
	\caption{Buwal directional morphology}
	\label{tab:buwa}
	\begin{tabular}{llll} % add l for every additional column or remove as necessary
		\lsptoprule
		forms & source gloss & direction  & other functions  \\ %table header
		\midrule
		\textit{-ā}  &  \textsc{vnt.prox}  &  \textsc{vent} &   \textsc{subs.vent}, \textsc{tam}  \\
		\textit{-xā}  &  \textsc{vnt.dist}  & \textsc{vent} &   \textsc{subs.vent}, \textsc{tam} \\
		\textit{āzà}  &  \textsc{it}  &   \textsc{itive}   &  \textsc{subs.itive} \\
		\lspbottomrule
	\end{tabular}
\end{table} 
\largerpage[-1]

Two ventive suffixes (`proximal' and `distal') and an itive post-verbal particle are described in detail by \citet{viljoenGrammaticalDescriptionBuwal2013}. A short discussion is also included in \citet{viljoenBuwalNarrativeDiscourse2015}. The proximal ventive suffix, \textit{-ā}, can have ventive directional meaning with intransitive and transitive motion verbs, as in (\ref{ex2:buwa01}) and (\ref{ex2:buwa5}).\footnote{The example in (\ref{ex2:buwa01}) is a bit unusual in that has both a ventive and itive marker -- an issue discussed below. No other examples of the proximal ventive suffix with intransitive motion verbs were found in the data, but this is an accidental gap as this suffix can regularly occur with intransitive motion verbs (Melanie Viljoen, personal communication).}

\ea\label{ex2:buwa01}
\langinfo{Buwal verb \textit{zèn} `return' with proximal ventive suffix }{}{\cite[378]{viljoenGrammaticalDescriptionBuwal2013}} \\
\gll  
ā-zèn-ā āzà ŋ́ ŋɡɮē ŋ́ ndā ŋ́ ɓàɬ dvàr	\\
3\gl{sg}.\gl{sbj}-return-\gl{vent}.\gl{prox} \gl{itive} \gl{prep} forge \gl{inf} go \gl{inf} forge hoe	\\
\glt `...he comes back from there into the forge (lit. to go) to forge a hoe.'

\ex\label{ex2:buwa5}
\langinfo{Buwal verb \textit{ndzēw} `pull' with proximal ventive suffix }{}{\cite[373]{viljoenGrammaticalDescriptionBuwal2013}} \\
\gll  ā-ndzēw-ā rā xèdzè fāgʷālākʷ	\\
\gl{sbj}.3\gl{sg}-pull-\gl{vent}.\gl{prox} hand person leper		\\
\glt `...he pulled out (toward him) the hand of the leper person.'
\z 

%\ea\label{ex2:buwa6}
%\langinfo{Buwal verb \textit{làɓ} `send' with proximal ventive suffix }{}{\cite[191]{viljoenGrammaticalDescriptionBuwal2013}} \\
%\gll  ɗālā má=ɬáɓ xʷā-làɓ-ā-zā ārá \\
%someone \gl{rel}=ready 2\gl{sg}.\gl{sbj}-send-\gl{vent}.\gl{prox}-\gl{tr} SIM \\
%\glt `Someone who is ready, you (should) send him on the way.'
%\z 

With non-motion verbs, the proximal ventive suffix can have a ventive subsequent motion interpretation, as in (\ref{ex2:buwa7}). Like the distal ventive suffix, it is also said to have temporal-aspectual functions, including completive aspect, as in (\ref{ex2:buwa8}). In several examples, the function of the suffix in unclear, as in (\ref{ex2:buwa9}).

\ea\label{ex2:buwa7}
\langinfo{Buwal verb \textit{ɮàk} `sow' with proximal ventive suffix }{}{\cite[374]{viljoenGrammaticalDescriptionBuwal2013}} \\
\gll   sā-ɮàk-ā lā nākā nɣā	 \\
\gl{sbj}.1\gl{sg}-sow-\gl{vent}.\gl{prox} field 1\gl{sg}.\gl{poss} \gl{dem}.\gl{prox}		\\
\glt `I sowed this field of mine (and returned). (The field is visible.)'

\ex\label{ex2:buwa8}
\langinfo{Buwal verb \textit{mār} `begin' with proximal ventive suffix }{}{\cite[376]{viljoenGrammaticalDescriptionBuwal2013}} \\
\gll xājāk bwāl ā-mār-ā á tā blàkʷ téŋɡʷlèŋ á témérè nfáɗ	\\
country Buwal \gl{sbj}.3\gl{sg}-begin-\gl{vent}.\gl{prox} \gl{prep} on thousand one \gl{prep} hundred four	\\
\glt `The Buwal country began in 1400. (The Buwal country still continues to this day.)'

\ex\label{ex2:buwa9}
\langinfo{Buwal verb \textit{ɮàr} `open' with proximal ventive suffix }{}{\cite[407]{viljoenGrammaticalDescriptionBuwal2013}} \\
\gll nà tsá xèdzè-jé má=á bwāl=éɡē j́-ká-ɮàr-ā ndzé āzá tsékʷɗē màvdāj nāsārā=éɡē	\\
now \gl{top} person-\gl{pl} \gl{rel}=\gl{prep} Buwal=\gl{pl} 3\gl{pl}.\gl{sbj}-\gl{pfv}-open-\gl{vent}.\gl{prox} eye \gl{compl} a.little because white.man(Ful.)=\gl{pl}		\\
\glt `Now people who are in Buwal have opened their eyes a little, because of white men...'
\z 

The distal ventive suffix, \textit{-xā}, has directional meaning with motion verbs, as in (\ref{ex2:buwa100}), and subsequent motion with non-motion verbs, as in (\ref{ex2:buwa2}). 

\ea\label{ex2:buwa100}
\langinfo{Buwal verb \textit{dám} `enter/exit' with distal ventive suffix }{}{\cite[372]{viljoenGrammaticalDescriptionBuwal2013}} \\
\gll   dám-xā á bzā	\\
come.out-\gl{vent}.\gl{dist} \gl{prep} outside		\\
\glt `Come outside. (Speaker is outside.)'

\ex\label{ex2:buwa2}
\langinfo{Buwal verb \textit{ɮàk} `sow' with distal ventive suffix }{}{\cite[374]{viljoenGrammaticalDescriptionBuwal2013}} \\
\gll  sā-ɮàk-xā	lā	nākā	á	dámāw \\
\gl{sbj}.1\gl{sg}-sow-\gl{vent}.\gl{dist}	field	1\gl{sg}.\gl{poss}	\gl{prep}	bush \\
\glt `I sowed my field in the bush (and returned). (The field is not visible.)'
\z 

Other examples of the distal ventive suffix do not appear to have any motion-related meaning, as in (\ref{ex2:buwa3}) and (\ref{ex2:buwa4}). The ventive distal suffix is said to have temporal-aspectual meanings including inceptive aspect (in the present and past), completive aspect (in the past) and prospective aspect (in the future), as in (\ref{ex2:buwa4}) \citep[375--376]{viljoenGrammaticalDescriptionBuwal2013}.

\ea\label{ex2:buwa3}
\langinfo{Buwal verb \textit{skām} `buy' with distal ventive suffix }{}{\cite[211]{viljoenGrammaticalDescriptionBuwal2013}} \\
\gll ājāw sā-ká-skām-xā nklèf ārá ákʷāw séj kʷésē	 \\
yes 1\gl{sg}.\gl{sbj}-\gl{pfv}-buy-\gl{vent}.\gl{dist} fish \textsc{simultaneity} \gl{neg} only(Ful.) doughnut	\\
\glt `Yes, I didn’t buy any fish at the same time, only doughnuts.'

\ex\label{ex2:buwa4}
\langinfo{Buwal verb \textit{ɡʷàrzàm} `get up' with distal ventive suffix }{}{\cite[375]{viljoenGrammaticalDescriptionBuwal2013}} \\
\gll ā-ná-ɡʷàrzàm-xā	vāɡʷmtáɗ \\ 
\gl{sbj}.3\gl{sg}-\gl{fut}-get.up-\gl{vent}.\gl{dist}	day.after.tomorrow \\
\glt `He will leave tomorrow.'
\z 

The itive enclitic \textit{āzà} can have an itive directional interpretation with both intransitive and transitive motion verbs, as in (\ref{ex2:buwa10}) and (\ref{ex2:buwa11}).

\ea\label{ex2:buwa10}
\langinfo{Buwal verb \textit{āntā} `run' with itive enclitic}{}{\cite[280]{viljoenGrammaticalDescriptionBuwal2013}} \\
\gll ā-xēj āntā āzà skʷá	\\
\gl{sbj}.3\gl{sg}-ran 3\gl{sg}.\gl{poss} \gl{itive} \gl{q}	\\
\glt `…did he run away?'

\ex\label{ex2:buwa11}
\langinfo{Buwal verb \textit{zmbàr} `throw far' with itive enclitic}{}{\cite[377]{viljoenGrammaticalDescriptionBuwal2013}} \\
\gll  vàkʷtāŋ ā-zmbàr āzà rā má-ràɮ-ā-ràɮ wēsé	\\
throw.far 3\gl{sg}.\gl{sbj}-throw.far \gl{itive} hand \gl{inf}-cut-\gl{nmlz}-cut \gl{dem}.\gl{dist}		\\
\glt `He threw away that cut off hand.'
\z

With non-motion verbs, the itive enclitic is said to ``imply that there is movement away from the place where the activity has been conducted'' \citep[378]{viljoenGrammaticalDescriptionBuwal2013}, as in (\ref{ex2:buwa12}) and (\ref{ex2:buwa13}). This can be interpreted as a type of itive subsequent motion, though this is not always made explicit in the translations. 

\ea\label{ex2:buwa12}
\langinfo{Buwal verb \textit{mbāl} `catch' with itive enclitic}{}{\cite[378]{viljoenGrammaticalDescriptionBuwal2013}} \\
\gll mbāl āzà ɡāmtāk	\\
catch \gl{itive} chicken	\\
\glt `Catch the chicken from there! (The addressee is near the chicken. He catches it and brings it away with him.)'

\ex\label{ex2:buwa13}
\langinfo{Buwal verb \textit{skèn} `grind' with itive enclitic}{}{\cite[378]{viljoenGrammaticalDescriptionBuwal2013}} \\
\gll  sā-skèn āzà ndrèj mávāw á màsēn	\\
1\gl{sg}.\gl{sbj}-grind \gl{itive} sorghum beer \gl{prep} mill(Fr.)		\\
\glt `I grind beer sorghum from there at the mill. (The sorghum is brought away from there. Said in some other place.)'
\z 

The itive form does not have a TAM interpretation, but another post-verbal clitic with the same segments but distinct tone and syntactic distribution, \textit{āzá}, is described as a completive marker. It is possible that this marker is grammaticalized from the itive. There is also a hortative form, \textit{āzā} ‘go.\textsc{hort}’, which might also be etymologically related to the itive marker \citep[471--472]{viljoenGrammaticalDescriptionBuwal2013}. 

The proximal and distal ventive suffixes may co-occur. Examples of their co-occurrence are all with non-motion verbs. \citet[377]{viljoenGrammaticalDescriptionBuwal2013} gives two examples, including (\ref{ex2:buwa14}), in which the interpretation is that some time passes between the main activity and the subsequent motion.\footnote{Note that in (\ref{ex2:buwa14}) the translation makes explicit that the subsequent motion is expected but not entailed in the same way as the main activity. This would actually exclude this use of the suffix from (at least some definitions of) the category of associated motion.} Other examples of this combination of suffixes, as in (\ref{ex2:buwa15}), do not make this temporal aspect of the meaning explicit. 

\largerpage
\ea\label{ex2:buwa14}
\langinfo{Buwal verb \textit{skām} `buy' with proximal and distal ventive suffixes}{}{\cite[377]{viljoenGrammaticalDescriptionBuwal2013}} \\
\gll á-kā-skām-ā-xā nklèf á lẃmà	\\
\gl{sbj}.3\gl{sg}-\gl{ipfv}-buy-\gl{vent}.\gl{prox}-\gl{vent}.\gl{dist} fish \gl{prep} market(Ful.)		\\
\glt `He is buying fish from the market (and will return). (He left a few days ago and hasn't come back yet.)'

\ex\label{ex2:buwa15}
\langinfo{Buwal verb \textit{zàm} `eat' with proximal and distal ventive suffixes}{}{\cite[386]{viljoenGrammaticalDescriptionBuwal2013}} \\
\gll	sā-zàm-ā-xā wdā āká á wātā \\
1\gl{sg}.\gl{sbj}-eat-\gl{vent}.\gl{prox}-\gl{vent}.\gl{dist} food ?? \gl{prep} home \\
\glt `I just came from eating at home.'
\z 

\citet[378]{viljoenGrammaticalDescriptionBuwal2013} gives one example of the proximal ventive suffix co-occur\-ring with the itive enclitic, shown in (\ref{ex2:buwa16}), and one example of the distal ventive suffix co-occurring with the itive enclitic, shown in (\ref{ex2:buwa17}). In both cases, the verb is an intransitive motion verb. The explanation for the co-occurrence of particles of seemingly contradictory meaning is that the ``the deictic reference point of \textit{āzà} is always some place other than the speaker of the location... the itive marker referring to the location where the movement begins, and the ventive suffix referring to the location where the movement ends'' \citep[378]{viljoenGrammaticalDescriptionBuwal2013}. In other words, the movement is `towards here' and `away from there'.

\ea\label{ex2:buwa16}
\langinfo{Buwal verb \textit{tàdàkʷ} `descend' with itive and distal ventive extensions}{}{\cite[378]{viljoenGrammaticalDescriptionBuwal2013}} \\
\gll  xʷné-tàdàkʷ-xā āzà á tā xājāk \\
2\gl{pl}.\gl{sbj}-descend-\gl{vent}.\gl{dist} \gl{itive} \gl{prep} on ground \\
\glt `You come down from there onto the ground! (Speaker on the ground.)'

\ex\label{ex2:buwa17}
\langinfo{Buwal verb \textit{zèn} `return' with itive and proximal ventive extensions}{}{\cite[378]{viljoenGrammaticalDescriptionBuwal2013}} \\
\gll ā-zèn-ā āzà ŋ́ ŋɡɮē ŋ́ ndā ŋ́ ɓàɬ dvàr	\\
3\gl{sg}.\gl{sbj}-return-\gl{vent}.\gl{prox} \gl{itive} \gl{prep} forge \gl{inf} go \gl{inf} forge hoe	\\
\glt `…he comes back from there into the forge (lit. to go) to forge a hoe.'
\z 


\subparagraph{Caused accompanied motion (CAM)}

The verb \textit{dā} `bring' can take the distal ventive suffix, as in (\ref{ex2:buwa18}). The verb \textit{ɮāɗ} `take away, pick up, collect' is found with the proximal ventive suffix and with the itive enclitic \citep[24, 151, 168, 179]{viljoenBuwalNarrativeDiscourse2015}, but without any directional meaning in the translations (i.e., not directed CAM). 

\ea\label{ex2:buwa18}
\langinfo{Buwal verb \textit{dā} `bring' with distal ventive suffix}{}{\cite[70]{viljoenBuwalNarrativeDiscourse2015}} \\
\gll  ājā á twsé j́-dā-xā kàn ɗērēwēl má=á tā māj-ātā ézē \\
then \gl{prep} there 3\gl{pl}.\gl{sbj}-bring-\gl{vent}.\gl{dist} thing paper(\gl{Ful.}) \gl{rel}=\gl{prep} on choose-3\gl{pl}.\gl{obj} therefore \\
\glt `So then at that time they brought the voting slips (lit. the paper for choosing them).'
\z 

\subparagraph{Buy-sell verbs}

The verb \textit{skām} can mean `buy' or `sell' \citep[379]{viljoenGrammaticalDescriptionBuwal2013}. With the ventive proximal suffix, the verb maintains its ambiguity \citep[330, 374, 571]{viljoenGrammaticalDescriptionBuwal2013}. The only examples of the verb with the ventive distal suffix are translated `buy', but this may be coincidental. The verb \textit{bér} also means `sell' though Buwal speakers report this to be a borrowed word (Melanie Viljoen, personal communication). 

\paragraph*{Gavar (gava1241)} \label{gava}

\subparagraph{Directional verbs}~

Information on Gavar grammar is from an extended grammar sketch published as an online manuscript \citep{viljoenGavarVerbPhrase2017}.

\begin{table}[H]
	\caption{Gavar directional verbs} 
	\label{tab:gavaV}
	\begin{tabular}{lll} % add l for every additional column or remove as necessary
		\lsptoprule
		DIR & verb & gloss   \\ %table header
		\midrule
		MOTION & \textit{nda} & `go' \\
		UP & \textit{taf} & `ascend, climb' \\
		DOWN & \textit{taduk}, \textit{srahw} & `descend' \\
		\lspbottomrule
	\end{tabular}
\end{table}

Gavar has a general motion verb, \textit{nda}, not specified for direction. A ventive interpretation (`come') can be enforced by the use of the ventive suffix.

\subparagraph{Directional extensions}~

\begin{table}[H]
	\caption{Gavar directional morphology}
	\label{tab:gava}
	\begin{tabular}{llll} % add l for every additional column or remove as necessary
		\lsptoprule
		forms & source gloss & direction  & other functions  \\ %table header
		\midrule
		\textit{-(h)à}  &  \textsc{vnt}  &   \textsc{vent} &  \textsc{subs.vent}  \\
		\lspbottomrule
	\end{tabular}
\end{table} 

Gavar has a single directional marker, the ventive suffix -(h)à \citep{viljoenGavarVerbPhrase2017}. The ventive combines with motion verbs to indicate that the path of motion is towards the deictic center, as in (\ref{ex2:gava1}) and (\ref{ex2:gava2}). 

\ea\label{ex2:gava1}
\langinfo{Gavar verb \textit{nda} `go/come' with ventive suffix}{}{\cite[23]{viljoenGavarVerbPhrase2017}} \\
\gll à:-nda-xà	á	sa	kʷə̀nə̀	 \\
3\gl{pl}.\gl{sbj}-go/come-\gl{vent}	\gl{inf}	drink	millet.drink \\
\glt `They come to drink the millet drink.'

\ex\label{ex2:gava2}
\langinfo{Gavar verb \textit{ʒin} `return' with ventive suffix}{}{\cite[35]{viljoenGavarVerbPhrase2017}} \\
\gll  séj	à:-má-vax-xà	dàɡálá	ábə̀	àːɡí-áː-ʒin-xà	á	xàjàk	ɡàvàr \\
except(Ful.)	3\gl{pl}.\gl{sbj}-\gl{sbjv}-spend.time-\gl{vent}	a.lot	with	3\gl{pl}.\gl{sbj}-\gl{fut}-return-\gl{vent}	to	country	Gavar \\
\glt `They must spend a lot of time (away), before they (will) return to the land of Gavar.'
\z 

When attached to a non-motion verb, the ventive can have a subsequent motion interpretation, as in (\ref{ex2:gava3}) and (\ref{ex2:gava4}).

\ea\label{ex2:gava3}
\langinfo{Gavar verb \textit{nkɪʃ} `urinate' with ventive suffix}{}{\cite[36]{viljoenGavarVerbPhrase2017}} \\
\gll  sà-nkɪʃ-à-ɓà	bə̀zà \\
1\gl{sg}.\gl{sbj}-urinate-\gl{vent}-\gl{ben}	\gl{compl} \\
\glt `I will urinate there (before coming back) first.'

\ex\label{ex2:gava4}
\langinfo{Gavar verb \textit{ʔʷəs} `cultivate' with ventive suffix}{}{\cite[37]{viljoenGavarVerbPhrase2017}} \\
\gll  xʷà-ʔʷəs-xɐ̀ \\
2\gl{sg}.\gl{sbj}-cultivate-\gl{vent} \\
\glt `You cultivate (there before returning home).'
\z 

There are also examples in which the translation does not make explicit what the semantic contribution of the ventive suffix is. In (\ref{ex2:gava5}), there is no indication of a change of location between the activity expressed by the ventive-marked verb and the activity expressed by the following predicate. The meaning of the ventive suffix in (\ref{ex2:gava6}) is uncertain.

\ea\label{ex2:gava5}
\langinfo{Gavar verb \textit{mbla} `pick' with ventive suffix}{}{\cite[33]{viljoenGavarVerbPhrase2017}} \\
\gll tá-mbla-xà	wə̀rì	bə̀zà	ɬetʃ	á-ɬetʃ	bə̀zà \\
\gl{pfv}-pick-\gl{vent}	vegetables	\gl{compl}	pull.off.stems	3\gl{sg}.\gl{sbj}-pull.off.stems	\gl{compl} \\
\glt `(When) she has picked the vegetables, she pulls off the stems.'

\ex\label{ex2:gava6}
\langinfo{Gavar verb \textit{ɗaː} `rain' with ventive suffix}{}{\cite[73]{viljoenGavarVerbPhrase2017}} \\
\gll  mànávàn	à-ká-diw	á	ɗaː-xà	tʃí ʃà	mà	ká-ɗaː	á	xà	tə̀ \\
rain	3\gl{sg}.\gl{sbj}-\gl{ipfv}-begin	\gl{inf}	rain-\gl{vent}	\gl{top} 3\gl{sg}	\gl{rel}.\gl{sbj}	\gl{ipfv}-rain	with	head	3\gl{sg}.\gl{poss} \\
\glt `When it begins to rain, it rains by itself.'
\z

\subparagraph{Caused accompanied motion (CAM)}

Verbs expressing caused accompanied motion include the itive \textit{dʒaɓ} `take along' and the ventive \textit{daa} `bring'. The latter can be combined with the ventive suffix, as in (\ref{ex2:gava9}).

\ea\label{ex2:gava9}
\langinfo{Gavar verb \textit{daː} `bring' with ventive suffix}{}{\cite[59]{viljoenGavarVerbPhrase2017}} \\
\gll ʔʷən à:-ʔʷən lékʷál àká à:-ká-daː-xà àlà  mà srək-ànə̀ skə̀n á wə̀ʒí mà ntáɡʷə̀lèŋ\\
build 3\gl{pl}.\gl{sbj}-build school(Fr.) now 3\gl{pl}.\gl{sbj}-\gl{ipfv}-bring-\gl{vent} someone \gl{rel}.\gl{sbj} teach-3\gl{sg}.\gl{dat} thing to children \gl{rel}.\gl{sbj} alone \\
\glt `(When) they have just built a school, they bring only one teacher.'
\z 

\subparagraph{Buy-sell verbs}

The verb \textit{skəm} is glossed `buy' and is found twice with a ventive suffix, as in (\ref{ex2:gava10}). The verb \textit{ber} means `sell'.

\ea\label{ex2:gava10}
\langinfo{Gavar verb \textit{skəm} `buy' with ventive suffix}{}{\cite[55]{viljoenGavarVerbPhrase2017}} \\
\gll tá-skəm-à skə̀n àːrá ká \\
\gl{pfv}-buy-\gl{vent} thing \textsc{simultaneity} \gl{purp} \\
\glt `He bought something on the way...'
\z 

\paragraph*{Daba (nucl1683)} \label{daba}

The description of Daba morphosyntax is spread across several papers focused on particular topics \citep{lienhardDabaGrammarClause1978,lienhardVerbeDaba1980,lienhardModaliteVerbeDaba1986,lienhardDocumentReferencePour2009}. There is also an online dictionary \citep{webonaryDictionnaireDaba2019}. Note that \citet{schuhChadicCornucopia2017} conflates Daba (nucl1683) with the closely related Mazagway language (maza1304).



\subparagraph{Directional verbs}~

\begin{table}[H]
\caption{Daba directional verbs} 
\label{tab:dabaV}
\begin{tabular}{lll} % add l for every additional column or remove as necessary
	\lsptoprule
	DIR & verb & gloss   \\ %table header
	\midrule
	ITIVE & \textit{vá} & `\textit{aller} [go]' \\
	VENT & \textit{ya} & `\textit{venir} [come]' \\
	UP & \textit{tàp} & `\textit{monter} [go up]' \\
	DOWN & \textit{jèm} &`\textit{descendre} [go down]' \\
	INTO& \textit{tūh} & `\textit{entrer} [enter]' \\
	\lspbottomrule
\end{tabular}
\end{table}

\subparagraph{Directional extensions}~

\begin{table}[H]
\caption{Daba directional morphology}
\label{tab:daba}
\begin{tabular}{llll} % add l for every additional column or remove as necessary
	\lsptoprule
	forms & source gloss & direction  & other functions  \\ %table header
	\midrule
	\textit{-a/-(á)ha}  &  \textit{vers ici} &   \textsc{vent} &   \textsc{subs.vent}  \\
	\textit{=ay}  &  \textit{vers toi} & \textsc{itive}  &    \\ 	
	\lspbottomrule
\end{tabular}
\end{table} 

The ventive and itive markers are only shown with a limited number of verbs in the data. With the verb  \textit{ya} `come', it can be assumed that the ventive suffix redundantly expresses ventive direction, as in (\ref{ex2:daba1}), but unfortunately no examples were found with other motion verbs to confirm this assumption.

\ea\label{ex2:daba1}
\langinfo{Daba verb \textit{ya} `come' with ventive suffix }{}{\cite[3]{lienhardVerbeDaba1980}} \\
\gll tâ yà-há	\\
1\gl{sg} come-\gl{vent} \\
\glt `I come here.' [French: `\textit{Je viens ici.}']
\z 

With the non-motion verbs \textit{gə̀r} `find' and \textit{bld} `cut', as in (\ref{ex2:daba2}), there is a subsequent ventive motion interpretation. With the verb \textit{wət} `take', the interpretation is caused accompanied motion, discussed below.

\ea\label{ex2:daba2}
\langinfo{Daba verb \textit{bdl} `cut' with ventive suffix }{}{\cite[45]{lienhardModaliteVerbeDaba1986}} \\
\gll ni bdl-a pay \\
1\gl{pl}.\gl{excl} cut-\gl{vent} wood \\
\glt `We are coming (towards here) from cutting the wood. \newline
[French: \textit{Nous venons (vers ici) de couper du bois.}']
\z 

The itive suffix is found with only two verb stems: \textit{vá} `go' and \textit{tap} `ascend'. The directional suffix is apparently redundant in the first case, but with \textit{tap} `ascend' the itive function is more clearly itive, as in (\ref{ex2:daba3}). 

\ea\label{ex2:daba3}
\langinfo{Daba verb \textit{tap} `ascend' with itive suffix }{}{\cite[46]{lienhardModaliteVerbeDaba1986}} \\
\gll tâ va ka tap=ay a ki \\
1\gl{sg} \gl{fut} ?? go.up=\gl{itive} \gl{prep} where	\\
\glt `Which way do I go up (to arrive at your place)?' \newline
[French: `\textit{Par où est-ce que je peux monter} (\textit{pour arriver chez toi}) ?']
\z 

There is one example, shown in (\ref{ex2:daba6}), in which a manner of motion verb, \textit{si} ‘run’, is followed by a directional motion verb. This combination of verbs is found frequently in directional serial verb constructions in other languages, but there is not enough data to determine how it is used in Daba. 

\ea\label{ex2:daba6}
\langinfo{Possible Daba multiverb construction }{}{\cite[3]{lienhardDabaGrammarClause1978}} \\
\gll kàtá	sì	và	wə̀t	tèp-èn-u	ítá \\
1\gl{sg}	run	go	carry	bring-him-it	house \\
\glt `I am going to carry it to him at the house.' \newline [French: `\textit{Je vais le lui apporter à la maison}.']
\z 

\subparagraph{Caused accompanied motion (CAM)}

Ventive caused accompanied motion can be expressed by the use of the ventive suffix with the verb \textit{wə̀t} `take' (also glossed `\textit{apporter}, \textit{porter}' [`bring, carry']), as in (\ref{ex2:daba7}). The verb \textit{da} `carry' is also used to express (non-directed) CAM \citep[45]{lienhardVerbeDaba1980}.

\ea\label{ex2:daba7}
\langinfo{Daba verb \textit{wət} `take' with ventive suffix }{}{\cite[3]{lienhardVerbeDaba1980}} \\
\gll  tâ wət-a pay \\
1\gl{sg} bring-\gl{vent} wood \\
\glt `I bring the wood here.' [French: `\textit{J'apporte le bois ici.}']
\z 

\subparagraph{Buy-sell verbs}

The verb \textit{səkəm} can be translated `buy' or `sell'. No morphosyntactic means of distinguishing these two meanings was found in the available data. There is also a verb \textit{ber} `sell' \citep{webonaryDictionnaireDaba2019}.

\paragraph*{Mazagway (maza1304)} \label{maza}

The information available on Mazagway grammar is from a grammar sketch and lexicon \citep{mouchetParlerDabaEsquisse1966}. 

\subparagraph{Directional verbs}~

\begin{table}[H]
\caption{Mazagway directional verbs} 
\label{tab:mazaV}
\begin{tabular}{lll} % add l for every additional column or remove as necessary
	\lsptoprule
	DIR & verb & gloss   \\ %table header
	\midrule
	ITIVE & \textit{va} & `\textit{aller} [go]' \\
	VENT & \textit{ya} & `\textit{venir} [come]' \\
	UP & \textit{tap} &`\textit{monter} [go up]' \\
	INTO& \textit{put} & `\textit{entrer} [enter]' \\
	\lspbottomrule
\end{tabular}
\end{table}

There is no lexical form given for outward motion; rather, one of the basic motion verbs is used with the word \textit{cadak} `outside' \citep[220]{mouchetParlerDabaEsquisse1966}. 

\subparagraph{Directional extensions}~

\begin{table}[H]
\caption{Mazagway directional morphology}
\label{tab:maza}
\begin{tabular}{llll} % add l for every additional column or remove as necessary
	\lsptoprule
	forms & source gloss & direction  & other functions  \\ %table header
	\midrule
	\textit{-aha}  &  \textit{directif}  &   \textsc{vent} &   \textsc{subs.vent}   \\
	\lspbottomrule
\end{tabular}
\end{table} 

The only directional form \citet{mouchetParlerDabaEsquisse1966} describes is a ventive suffix \textit{-aha} (called \textit{directif}), as in (\ref{ex2:maza1}). 

\ea\label{ex2:maza1}
\langinfo{Mazagway verb \textit{nə} `enter' with ventive suffix}{}{\cite[85]{mouchetParlerDabaEsquisse1966}} \\
\gll  nə	put	aha \\
1\gl{sg}	enter	\gl{vent} \\
\glt `Enter; the speaker is in the hut, so there is movement toward him.' \newline [French: `\textit{Entrez} ; \textit{le parlant est dans la case}, \textit{il y a donc mouvement vers lui}.']
\z 

\citet[85]{mouchetParlerDabaEsquisse1966} lists a few examples of the ventive suffix used with non-motion verbs for which the translation suggests a prior motion interpretation, for example: \textit{təvkadar aha} `I will go dance.' (French: `\textit{J'irai danser}.'), \textit{nəvkižəgaha} `We will go sow (seeds).' (French: `\textit{Nous irons semer}.'). According to a Mazagway speaker, Alim Tekilem Abednego (personal communication), those translations are not accurate. The correct translation should be of subsequent motion, as in \textit{tə ham-aha} ‘I ate and came here.’ Mouchet‘s translations lead \citet[313]{schuhChadicCornucopia2017} to label the ventive suffix \textit{-aha} as an itive directional. \citet[313]{schuhChadicCornucopia2017} also follows \citet{newmanChadicExtensionsPredative1977} in making a tenuous etymological connection between ventive directionals and the (non-motion) suffix \textit{-eŋ} which \citet[84]{mouchetParlerDabaEsquisse1966} labels \textit{destinatif}. %There is no evidence that this or was a motion-related suffix. It is more likely a benefactive suffix etymologically related to an indirect object pronominal. 

\subparagraph{Caused accompanied motion (CAM)}

The verb \textit{wut} ‘take’ can be combined with the ventive suffix to express directed caused accompanied motion, as in (\ref{ex2:maza11}) \citep[cf.][170]{mouchetParlerDabaEsquisse1966}. The verb \textit{daha}, glossed `\textit{apporter} [carry]', also expresses CAM without any indication of inherent directional semantics.

\ea\label{ex2:maza11}
\langinfo{Mazagway verb \textit{wut} `take' with ventive suffix}{}{\cite[73]{mouchetParlerDabaEsquisse1966}} \\
\gll tu	wut	aha \\
3\gl{sg}	take	\gl{vent} \\	
\glt `He brought (it) here (towards the speaker).'
\z 

\subparagraph{Buy-sell verbs}

The verb \textit{skəm} can mean `buy' or `sell' \citep[150]{mouchetParlerDabaEsquisse1966}. No examples are given of this verb with a ventive suffix. According to Alim Tekilem Abednego (personal communication), the verb means `buy' and can take a ventive suffix without changing its meaning. The verb \textit{ber} means `sell'.

\paragraph*{Mbudum (mbed1242)} \label{mbud}

Information on Mbudum grammar is from a doctoral dissertation \citep[]{fobakayyangPhonologieMorphologieSyntaxe2021}.

\subparagraph{Directional verbs}~

\begin{table}[H]
\caption{Mbudum directional verbs} 
\label{tab:mbudV}
\begin{tabular}{lll} % add l for every additional column or remove as necessary
	\lsptoprule
	DIR & verb & gloss   \\ %table header
	\midrule
	MOTION & \textit{nda} & `\textit{aller} [go]' \\
	DOWN &  \textit{brɛ̀ŋ} & `\textit{descendre} [go down]' \\
	INTO/OUT & \textit{təɗ} & `\textit{entrer}/\textit{sortir} [enter/exit]' \\
	\lspbottomrule
\end{tabular}
\end{table}

The verb \textit{nda} is a general verb of motion can be specified for ventive directionality through the verbal morphology. The verb \textit{təɗ} is glossed `enter' unless it takes the ventive suffix \textit{-àhà} in which case it is glossed `exit' \citep[290]{fobakayyangPhonologieMorphologieSyntaxe2021}. The examples given do not specify the deictic orientation, but it is possible, as in related languages, that deixis is how the `enter/exit' interpretation is determined. 

\subparagraph{Directional extensions}~

\begin{table}[H]
\caption{Mbudum directional morphology}
\label{tab:mbud}
\begin{tabular}{llll} % add l for every additional column or remove as necessary
	\lsptoprule
	forms & source gloss & direction  & other functions  \\ %table header
	\midrule
	\textit{-áhà}  &  \textsc{ma}  &   \textsc{vent} &   \textsc{prior.itive}  \\
	\lspbottomrule
\end{tabular}
\end{table} 

\noindent
\citet{fobakayyangPhonologieMorphologieSyntaxe2021} describes the suffix \textit{-áhà} as ``associated movement'' (French: \textit{mouvement associé}). The suffix \textit{-àhà} is said to have an allophone \textit{-à} when the verb is followed by a direct object \citep[292]{fobakayyangPhonologieMorphologieSyntaxe2021}, as in the alternation seen in two instances of the verb \textit{kərət} ‘boil’ in (\ref{ex2:mbud8}).\footnote{However, there are some possible exceptions to this where the form \textit{-àhà} appears on a verb followed by a direct object, as in (\ref{ex2:mbud7}) \citep[see also][175, 238]{fobakayyangPhonologieMorphologieSyntaxe2021} and at least one example of the \textit{-à} suffix not followed by a direct object \citep[265]{fobakayyangPhonologieMorphologieSyntaxe2021}.}

\ea\label{ex2:mbud8}
\langinfo{Mbudum verb \textit{kərət} `boil' with ventive suffix}{}{\cite[436]{fobakayyangPhonologieMorphologieSyntaxe2021}} \\
\ea \label{ex2:mbud8a}
\gll  hwa	bas-aŋ	kahaw	a	kərət-a	məsləpic \\
2\gl{sg}.\gl{sbj}	light-3\gl{sg}.\gl{obj}	fire	3\gl{sg}.\gl{sbj}	boil-\gl{vent}	waste \\
\glt `You light a fire and boil away the chaff (from the millet).'

\ex \label{ex2:mbud8b}
\gll  hwa	bət	məsləpic	za,	a	kərət-aha	baha \\
2\gl{sg}.\gl{sbj}	take	waste	\textsc{nprs}	3\gl{sg}.\gl{sbj}	boil-\gl{vent}	again \\
\glt `You take out the chaff and boil (the millet) again.'
\z 
\z 

As in the other Dabaic languages, the suffix has a ventive directional meaning with verbs of motion, as in (\ref{ex2:mbud1}) and (\ref{ex2:mbud2}). 

\ea\label{ex2:mbud1}
\langinfo{Mbudum verb \textit{zə̀ŋgʲ} `return' with ventive suffix}{}{\cite[36]{fobakayyangPhonologieMorphologieSyntaxe2021}} \\
\gll  zə̀ŋgʲ-àhà \\
return-\gl{vent} \\
\glt `Come back here.'

\ex\label{ex2:mbud2}
\langinfo{Mbudum verb \textit{bə̀ràŋgʲ} `descend' with ventive suffix}{}{\cite[291]{fobakayyangPhonologieMorphologieSyntaxe2021}} \\
\gll bə̀ràŋgʲ-àhà \\
descend-\gl{vent} \\
\glt `Come down here.'
\z 

The general motion verb \textit{nda} is not specified for deictic orientation, but is interpreted as a expressing a ventive path when the directional suffix \textit{-áhà} is attached, as in (\ref{ex2:mbud3}).

\ea\label{ex2:mbud3}
\langinfo{Mbudum verb \textit{ndà} `go/come' with ventive suffix}{}{\cite[290]{fobakayyangPhonologieMorphologieSyntaxe2021}} \\
\gll sà=ndà-àhà	ə́ŋg	wə́tà \\
1\gl{sg}.\gl{sbj}=go/come-\gl{vent}	\gl{prep}	house \\
\glt `I return to the house.' [French: `\textit{Je rentre à la maison}.']
\z 

With non-motion verbs, the suffix \textit{-àhà} can have an itive prior motion interpretation, as in (\ref{ex2:mbud4}) and (\ref{ex2:mbud14}).\footnote{Note that example (\ref{ex2:mbud14}) shows a prior motion interpretation of the suffix \textit{-àhà} attached to the verb \textit{sə̀kə̀m} `buy', but elsewhere the same verb+suffix combination does not have a motion-related translation \citep[143, 176, 181, 261, 343, 435]{fobakayyangPhonologieMorphologieSyntaxe2021}.} This use may be related to the prior motion multiverb construction discussed below. 

\ea\label{ex2:mbud4}
\langinfo{Mbudum verb \textit{ʣàw} `attach' with ventive suffix}{}{\cite[289]{fobakayyangPhonologieMorphologieSyntaxe2021}} \\
\gll sà=ʣàw	ɮə`j=ə́jə̀	ə́j=ɗə̀s-àhà	(ə́ŋg	dámə̀w) \\
1\gl{sg}.\gl{sbj}=attach	cattle=\gl{pl}	\gl{inf}=cultivate-\gl{vent}	\gl{prep}	bush \\
\glt `I attach the cattle to go cultivate (the fields).' \newline [French: `\textit{J'attache les boeufs pour aller cultiver (au champs)}.']

\ex\label{ex2:mbud14}
\langinfo{Mbudum verb \textit{sə̀kə̀m} `buy' with ventive suffix}{}{\cite[289]{fobakayyangPhonologieMorphologieSyntaxe2021}} \\
\gll à=sə̀kə̀m-àhà ndə̀rə̀j ə́ŋg lə́mà hwàsàm	 \\
3\gl{sg}.\gl{sbj}=buy-\gl{vent} millet \gl{prep} market Hosom	\\
\glt `He goes and buys millet at the market in Hosom.' \newline
[French: `\textit{Il part achèter le mil au marché de Hosom}.']
\z

In several examples, the suffix \textit{-àhà} does not appear to have any motion-related meaning, as in (\ref{ex2:mbud7}).  \citet[291]{fobakayyangPhonologieMorphologieSyntaxe2021} describes some uses of \textit{-àhà} as occurring in the context of a ``\textit{changement de processus}'' or a change of state.

\ea\label{ex2:mbud7}
\langinfo{Mbudum verb \textit{mbə̀w} `give birth' with ventive suffix}{}{\cite[233]{fobakayyangPhonologieMorphologieSyntaxe2021}} \\
\gll  wàlà	wà	kə̀=mbə̀w-àhà	dàrlàj=ə̀jə̀=zà	dàp	\\
woman	\gl{dem}.\gl{prox}	\textsc{nfut}=give.birth-\gl{vent}	daughter=\gl{pl}=\textsc{nprs}	only	 \\
\glt `This woman only gave birth to daughters.' \newline [French: `\textit{Cette femme a accouché seulement les filles.}']
\z 

The prior motion interpretation of the suffix \textit{-àhà} may be related to a prior motion construction where the verb \textit{nda} expresses motion away from the deictic center, as in (\ref{ex2:mbud5}) and (\ref{ex2:mbud6}). In these examples, the non-motion verb takes the suffix \textit {-àhà}. Without the suffix \textit {-àhà} on the second verb, this construction is said to express future tense \citep[220; but see 271 for counterexamples in the imperative]{fobakayyangPhonologieMorphologieSyntaxe2021}. This is a context in which a semantic shift to an itive prior motion meaning might occur. 

\ea\label{ex2:mbud5}
\langinfo{Mbudum verb \textit{ɗə̀s} `cultivate' with ventive suffix in multiverb construction}{}{\cite[255]{fobakayyangPhonologieMorphologieSyntaxe2021}} \\
\gll mbə̀w	à=ndà	ə́j=ɗə̀s-àhà	ábə̀	ɮà=ə́jə̀ \\
child	3\gl{sg}.\gl{sbj}=go	\gl{inf}=cultivate-\gl{vent}	with	cattle=\gl{pl} \\
\glt `The child goes and cultivates with the cattle.' \newline [French: `\textit{L’enfant part cultiver avec les boeufs.}']

\ex\label{ex2:mbud6}
\langinfo{Mbudum verb \textit{və̀ŋg} `vomit' with ventive suffix  in purposive construction}{}{\cite[292]{fobakayyangPhonologieMorphologieSyntaxe2021}} \\
\gll mádàj	tə̀kwà	à=ndà	ə́j=və̀ŋg-àhà	ə́ŋg	kwàdàm	tà \\
friend	\gl{poss}.2\gl{sg}	3\gl{sg}.\gl{sbj}=move	\gl{inf}=vomit-\gl{vent}	in	plate	\gl{det}.\gl{def} \\
\glt `Your friend is going to vomit in that plate.' \newline [French: `\textit{Ton ami va vomir dans l'assiete là.}']
\z 

\subparagraph{Caused accompanied motion (CAM)}

Caused accompanied motion (CAM) can be expressed by two verbs, neither of which are found with the ventive suffix. The verb \textit{ɗɛb} `bring, take' (French: \textit{amener}) can be used with any deictic orientation. The verb \textit{daha} `bring (here)' (French: \textit{apporter}, \textit{emmener}) only expresses ventive CAM (Foba K'ayyang, personal communication) and likely has an etymological link to the ventive suffix. Ventive CAM can also be expressed by the use of the ventive suffix with the verb \textit{ɓət} `take', as in \textit{ɓət-aha-w} `bring it here' (French: `\textit{prendre pour venir}') (Foba K'ayyang, personal communication). 

\subparagraph{Buy-sell verbs}

In Mbudum, the verb \textit{sə̀kə̀m}  means `buy' and the verb \textit{bɛ̀r} means `sell'. 

\paragraph*{Mina (mina1276)} \label{mina}

Information on the grammar of Mina is from a reference grammar \citep{frajzyngierGrammarMina2005}.

\subparagraph{Directional verbs}~

\begin{table}[H]
\caption{Mina directional verbs} 
\label{tab:minaV}
\begin{tabular}{lll} % add l for every additional column or remove as necessary
	\lsptoprule
	DIR & verb & gloss   \\ %table header
	\midrule
	MOTION & \textit{nd} & `go/come' \\
	DOWN & \textit{kùrə̀k} & `descend' \\
	INTO& \textit{nástə̀} & `enter (Fulfulde)' \\
	\lspbottomrule
\end{tabular}
\end{table}

The verb \textit{tíl} is variously glossed `go', `leave', `depart' or `enter' leaving its actual meaning (or range of meaning) unclear. 

\subparagraph{Directional extensions}~

\begin{table}[H]
\caption{Mina directional morphology}
\label{tab:mina}
\begin{tabular}{llll} % add l for every additional column or remove as necessary
	\lsptoprule
	forms & source gloss & direction  & other functions  \\ %table header
	\midrule
	\textit{-á(hà)}  &  \textsc{go}  &   \textsc{vent} &   \textsc{loc}  \\
	\lspbottomrule
\end{tabular}
\end{table} 

\citet[171]{frajzyngierGrammarMina2005} label the suffix \textit{-á(hà)} ``goal-orientation extension''. It will be referred to here as a ventive suffix. The suffix is said to have the form \textit{-á} in ``phrase-internal'' position and \textit{-áhà} in ``phrase-final'' position. The relevant context appears to be phonological, as in (\ref{ex2:mina1}) where the longer form is used only if a pause (represented by an ellipsis) is inserted between the verb and the following prepositional phrase. This alternation is said to be part of a general phonological rule whereby word-final vowels are deleted phrase-medially, although in this case the consonant \textit{h} is also deleted. However, other variations such as \textit{-áh}, \textit{-h}, \textit{-há} are found in a few transcriptions \citep[112, 316, 341]{frajzyngierGrammarMina2005}. Forms that include \textit{h} are much rarer than the single vowel suffix, appearing with only eight verb roots in the entire grammar.

\ea\label{ex2:mina1}
\langinfo{Mina verb \textit{nda} `go/come' with ventive suffix}{}{\cite[172]{frajzyngierGrammarMina2005}} \\
\gll í	n	kə̀ ták-á-k	kə̀ kè	nd-á(hà…)	á	pàt \\
3\gl{pl}	\gl{prep}	\gl{inf}	forbid-\gl{tr}-1\gl{sg}	??	\gl{inf}	go/come-\gl{vent}	\gl{prep} next.day \\
\glt `They will forbid me to come tomorrow.'
\z 

A form \textit{-á} also regularly occurs before pronominal objects (except 3\gl{pl}), as with the verb \textit{ták} `forbid' in (\ref{ex2:mina1}). \citet[97]{frajzyngierGrammarMina2005} also label this \textit{-á} ``goal-orientation'' even though this form in this position has no functional similarity with the \textit{-á(hà)} suffix.\footnote{In a critique of an analysis of Mina adpositions, \citet[455]{hellwigAdpositionsTheirDistribution2022} raises a question about the function of the \textit{-á} verbal suffix, to which \citet[461]{frajzyngierMeaningEncodedGrammatical2022} replies: ``The function of the verbal suffix \textit{-a} in Mina, glossed as `goal', indeed merits a revised analysis.''} The only case where a suffix glossed ‘goal orientation’ preceding an object suffix has an explicit motion-related interpretation is in (\ref{ex2:mina2}) where the full form of the directional \textit{-áhà} is used. This contrasts with the short form \textit{-á} further demonstrating that the form \textit{-á} before an object pronoun is not a motion-related morpheme. To avoid confusion, the pre-pronominal vowel is glossed here as transitive (\gl{tr}).

\ea\label{ex2:mina2}
\langinfo{Mina verb \textit{y} `call' with ventive suffix}{}{\cite[174]{frajzyngierGrammarMina2005}} \\
\gll  ndə̀	y-áhà-w \\
go	call-\gl{vent}-3\gl{sg} \\
\glt `Go call him here.'
\z 

The meaning of the suffix \textit{-á(hà)} is described as a ventive directional marker with motion verbs and as an altrilocative marker with non-motion verbs. As in other Dabaic languages, the verb \textit{nd} is a general verb of motion which receives an itive interpretation by default, as in (\ref{ex2:mina3}),\footnote{There is an error in the original gloss which marks the verb in (\ref{ex2:mina3}) as having the directional suffix.} and a ventive interpretation with the directional suffix, as in (\ref{ex2:mina4}). Unfortunately, this analysis is somewhat obscured by at least four examples of the motion verb \textit{nd-á} glossed as containing the ventive suffix (\textsc{go}) but with a translation that is either not explicit about the deictic orientation, or suggest an itive path, as in (\ref{ex2:mina14}).

\ea\label{ex2:mina3}
\langinfo{Mina verb \textit{nd} `go/come' without ventive suffix}{}{\cite[174]{frajzyngierGrammarMina2005}} \\
\gll  à	ndə̀	zə́	vù \\
3\gl{sg}	go/come	\textsc{end.of.event}	\gl{q} \\
\glt `Will he go?'

\ex\label{ex2:mina4}
\langinfo{Mina verb \textit{nd} `go/come' with ventive suffix}{}{\cite[174]{frajzyngierGrammarMina2005}} \\
\gll  à	nd-á	zə́	vù \\
3\gl{sg}	go/come-\gl{vent}	\textsc{end.of.event}	\gl{q} \\
\glt `Will he come?'

\ex\label{ex2:mina14}
\langinfo{Mina verb \textit{nd} `go/come' with ventive suffix?}{}{\cite[90]{frajzyngierGrammarMina2005}} \\
\gll káyéfì í-bə̀ nd-á tə̀tàŋ	\\
strange(Ful.) \gl{pl}-\gl{asoc} go/come-\gl{vent} 3\gl{pl}	\\
\glt `Never seen before, they up and left [the room].'
\z

The ventive interpretation can also be seen when the suffix attaches to a manner{\slash}of{\slash}motion verb, as in (\ref{ex2:mina5}), but note that there is no ventive meaning in the translation of the verb \textit{tál} `walk' with the ventive suffix in (\ref{ex2:mina15}).

\ea\label{ex2:mina5}
\langinfo{Mina verb \textit{sí} `run' with ventive suffix}{}{\cite[20]{frajzyngierGrammarMina2005}} \\
\gll  í	mə̀	 sí-á-yí	tə̀tə̀	zá \\
3\gl{pl}	\gl{rel}	run-\gl{vent}-\gl{stative}	3\gl{pl}	\textsc{end.of.event} \\
\glt `They have returned running.'

\ex\label{ex2:mina15}
\langinfo{Mina verb \textit{tál} `walk' with ventive suffix}{}{\cite[314]{frajzyngierGrammarMina2005}} \\
\gll  sə̀	tál-áhà	... \\
1\gl{sg}	walk-\gl{vent}	 \\
\glt `I took walks...'
\z

\citet[176]{frajzyngierGrammarMina2005} describe the other major function of the suffix as altrilocative: ``With non-movement verbs, the goal orientation indicates that the event happened at a place other than the place of speech.'' This interpretation is made explicit in the translations in (\ref{ex2:mina6}), (\ref{ex2:mina7}) and (\ref{ex2:mina8}). In some cases the translations are possibly compatible with a subsequent motion interpretation (e.g., re-translate (\ref{ex2:mina6})  as ‘He cooked it (elsewhere) then came here'). However, the habitual aspect in (\ref{ex2:mina7}) makes it unlikely that subsequent motion is a part of the semantics in this context. 

\ea\label{ex2:mina6}
\langinfo{Mina verb \textit{d} `cook' with ventive suffix}{}{\cite[172]{frajzyngierGrammarMina2005}} \\
\gll  kə́	d-á	zà \\
\gl{inf}	cook-\gl{vent}	\textsc{end.of.event} \\
\glt `He cooked it (not in the place of speech).'

\ex\label{ex2:mina7}
\langinfo{Mina verb \textit{màr} `shepherd' with ventive suffix}{}{\cite[175]{frajzyngierGrammarMina2005}} \\
\gll í	ndí	màr-áhà		\\
3\gl{pl}	\gl{hab}	shepherd-\gl{vent}	 \\
\glt `They have the habit of pasturing somewhere else.'

\ex\label{ex2:mina8}
\langinfo{Mina verb \textit{nz} `stay' with ventive suffix}{}{\cite[176]{frajzyngierGrammarMina2005}} \\
\gll sə̀	nz-á	màrɓák	\\
1\gl{sg}	stay-\gl{vent}	Marbak	 \\
\glt `I was in Marbak (said in Maroua).'
\z 

\largerpage
Many examples of \textit{-á(hà)} with non-motion verbs do not give any explicit information about motion or deixis, as in (\ref{ex2:mina9}). This could simply be an omission of information that would render the free translation less natural, or it could be the case that the suffix has other non-motion/locational meanings that have not yet been explored.

\ea\label{ex2:mina9}
\langinfo{Mina verb \textit{zə̀m} `eat' with ventive suffix}{}{\cite[27]{frajzyngierGrammarMina2005}} \\
\gll á zə̀m-á wùdá \\	
3\gl{sg} eat-\gl{vent} food	\\
\glt `He ate the food.'
\z

\subparagraph{Caused accompanied motion (CAM)}

The verbs \textit{déɓ} `bring' and \textit{da} `bring' with the ventive suffix \textit{-á(hà)} express ventive CAM. Ventive CAM can also be expressed by the verb \textit{ɓət} `take, get' with the ventive suffix, as in (\ref{ex2:mina13}). 

\ea\label{ex2:mina13}
\langinfo{Mina verb \textit{bə́t} `take' with ventive suffix}{}{\cite[398]{frajzyngierGrammarMina2005}} \\
\gll séy	í	lím	kə́	bə́t-á	yə̀m	zá	kə̀	mə̀ts	kúhú \\
so	3\gl{pl}	see	\gl{inf}	take-\gl{vent}	water	\textsc{end.of.event}	\gl{inf}	die	fire \\
\glt `They saw them fetch water to extinguish the fire.'
\z 

\subparagraph{Buy-sell verbs}

In Mina, the verb \textit{skəm} means `buy' and \textit{bèr} means `sell'. Both verbs can appear with the ventive marker. 

\subsection{Matakam (2/3)}

There are descriptions available of the morphosyntax of Cuvok and Mafa, but not of Mefele (mefe1242).

\paragraph*{Cuvok (cuvo1236)} \label{cuvo}

Information on the grammar of Cuvok is from a doctoral dissertation \citep{dadakGrammaireCuvokLangue2021}.

\largerpage
\subparagraph{Directional verbs}~

\begin{table}[H]
\caption{Cuvok directional verbs} 
\label{tab:cuvoV}
\begin{tabular}{lll} % add l for every additional column or remove as necessary
	\lsptoprule
	DIR & verb & gloss   \\ %table header
	\midrule
	MOTION & \textit{d} & `\textit{aller} [go]' \\
	UP & \textit{tàp} & `\textit{escalader}, \textit{monter} [climb, go up]' \\
	DOWN & \textit{bə̀ràŋ} & `\textit{descendre} [go down]' \\
	INTO& \textit{ʃɛ̀hw} & `\textit{enrter} [enter]' \\
	OUT & \textit{báɬ} & `\textit{sortir} [exit]' \\
	\lspbottomrule
\end{tabular}
\end{table}

\subparagraph{Directional extensions}~

\begin{table}[H]
\caption{Cuvok directional morphology}
\label{tab:cuvo}
\begin{tabular}{llll} % add l for every additional column or remove as necessary
	\lsptoprule
	forms & source gloss & direction  & other functions  \\ %table header
	\midrule
	\textit{-ék}  &  \textsc{cpt}  &   \textsc{vent} &   \textsc{subs.vent}  \\
	\textit{-áɗ}  &  \textsc{cfg}  & \textsc{itive}  &   \textsc{subs.itive}, lexicalized \\
	\lspbottomrule
\end{tabular}
\end{table} 

Cuvok has two directional suffixes, ventive and itive. The suffixes have a directional function with the general verb of motion, as in (\ref{ex2:cuvo1}) and (\ref{ex2:cuvo2}), with manner-of-motion verbs, as in (\ref{ex2:cuvo3}) and (\ref{ex2:cuvo4}).

\ea\label{ex2:cuvo1}
\langinfo{Cuvok verb \textit{d} `go/come' with ventive suffix}{}{\cite[303]{dadakGrammaireCuvokLangue2021}} \\
\gll d-ɛ́k	\\
go/come-\gl{vent} \\
\glt `Come here.'

\ex\label{ex2:cuvo2}
\langinfo{Cuvok verb \textit{d} `go/come' with itive suffix}{}{\cite[308]{dadakGrammaireCuvokLangue2021}} \\
\gll  d-áɗ \\
go/come-\gl{itive} \\
\glt `Go over there.'

\ex\label{ex2:cuvo3}
\langinfo{Cuvok verb \textit{hw} `run' with ventive suffix}{}{\cite[304]{dadakGrammaireCuvokLangue2021}} \\
\gll ndàh	má-ʦə̀n-á	mà	àná	háj	sə̀, ɛ̀ʧ-kɛ̀-hw-ɛ́k-ɛ́j	kà-tá-wùzàɗ	ʤə̀	màsá	kɛ̀-ŋg-ɛ́j \\
people	\gl{inf}-listen-\gl{intr}	mouth	\gl{def}	\gl{pl}	\gl{top} 3\gl{pl}.\gl{sbj}-P-run-\gl{vent}-\gl{intr}	\gl{purp}-\gl{fut}-see	thing	\gl{rel}	\gl{pst}.3\gl{sg}.\gl{sbj}-do-\gl{intr} \\
\glt `The people who heard those words ran here to see what happened.'

\ex\label{ex2:cuvo4}
\langinfo{Cuvok verb \textit{hw} `run' with itive suffix}{}{\cite[323]{dadakGrammaireCuvokLangue2021}} \\
\gll hw-áɗ	kàd	á-wùnàm	bə̀ná	gálàŋ	tá-ká	kà-mbə̀ɗ-àtá \\
run-\gl{itive}	toward	\gl{loc}-house	because	wall	\gl{rel}-2\gl{sg}	\gl{pst}.3\gl{sg}.\gl{sbj}-destroy-\gl{intr} \\
\glt `Go home because your wall fell down.'
\z 

With transitive verbs the ventive suffix can have a directional meaning, as in (\ref{ex2:cuvo5}) or a resultative motion meaning, as in (\ref{ex2:cuvo6}).

\ea\label{ex2:cuvo5}
\langinfo{Cuvok verb \textit{kɛ̀ɮ} `throw' with ventive suffix}{}{\cite[303]{dadakGrammaireCuvokLangue2021}} \\
\gll kɛ̀ɮ-ɛ́k \\
throw-\gl{vent} \\
\glt `Throw it here.'

\ex\label{ex2:cuvo6}
\langinfo{Cuvok verb \textit{kə̀ɗ} `hit' with itive suffix}{}{\cite[311]{dadakGrammaireCuvokLangue2021}} \\
\gll Kàdàmà	á-kə̀ɗ-áɗ	bàlɔ̀ŋ \\
Kadama	3\gl{sg}.\gl{sbj}-hit-\gl{itive}	ball \\
\glt `Kadama sends away the ball by hitting it.'
\z 

Both suffixes can also express subsequent motion when used with non-motion verbs, as in (\ref{ex2:cuvo7}), (\ref{ex2:cuvo8}) and (\ref{ex2:cuvo9}).

\ea\label{ex2:cuvo7}
\langinfo{Cuvok verb \textit{nd} `eat' with ventive suffix}{}{\cite[303]{dadakGrammaireCuvokLangue2021}} \\
\gll  nd-ɛ́k \\
eat-\gl{vent} \\
\glt `Eat and come here.'

\ex\label{ex2:cuvo8}
\langinfo{Cuvok verb \textit{hə̀v} `cultivate' with ventive suffix}{}{\cite[305]{dadakGrammaireCuvokLangue2021}} \\
\gll  ŋgwázàh	ɛ̀ʧ-kɛ̀-hə̀v-ɛ́k	lɛ̀j \\
women	3\gl{pl}.\gl{sbj}-\gl{pst}-cultivate-\gl{vent}	field \\
\glt `The women worked the field and have returned.'

\ex\label{ex2:cuvo9}
\langinfo{Cuvok verb \textit{hə̀v} `cultivate' with itive suffix}{}{\cite[308]{dadakGrammaireCuvokLangue2021}} \\
\gll  hə̀v-áɗ \\
cultivate-\gl{itive} \\
\glt `Cultivate and go there.'
\z 

However, (\ref{ex2:cuvo9}) is the only clear example of itive subsequent motion, and it appears to be elicited. Compare (\ref{ex2:cuvo7}) where the ventive suffix combines with the verb \textit{nd} ‘eat’ is translated as subsequent motion with (\ref{ex2:cuvo10}) where the itive suffix combines with the same verb, but there is no subsequent motion in the translation, nor is there any other indication of what the semantic contribution of the itive suffix is in this context. 

\ea\label{ex2:cuvo10}
\langinfo{Cuvok verb \textit{nd} `eat' with itive suffix}{}{\cite[195]{dadakGrammaireCuvokLangue2021}} \\
\gll Kàdìbàj	ɛ́-nd-ɛ́k	ɬàw	tá-zàŋgwà \\
Kadibay	3\gl{sg}.\gl{sbj}-eat-\gl{itive}	meat	\gl{rel}-donkey \\
\glt `Kadibay ate donkey meat.'
\z 

In many examples the translations do not explicitly indicate the function of the ventive/itive suffix. This could be due to the unnaturalness of making the semantics of the suffix explicit in the target language or it could indicate a grammaticalization or lexicalization of the suffix, but more research would be needed to be able to discern how the suffixes function in these cases. \citet[303]{dadakGrammaireCuvokLangue2021} suggests visibility may be a relevant part of the semantics, as in example (\ref{ex2:cuvo11}) in which the child being born becomes visible.

\ea\label{ex2:cuvo11}
\langinfo{Cuvok verb \textit{hɛ̀j} `give birth' with itive suffix}{}{\cite[231]{dadakGrammaireCuvokLangue2021}} \\
\gll kàɬ-pá	mɛ́-hɛ̀j-ɛ́k	tá-wàt	àná \\
\gl{loc}-\gl{loc}	\gl{inf}-give.birth-\gl{itive}	\gl{rel}-child	\gl{def} \\
\glt `until the birth of this child'
\z 

However, the use of the itive suffix with the verb \textit{mə̀ʦ} `die' in (\ref{ex2:cuvo12}) is a case where neither visibility nor motion seem to be likely explanation for the use of the itive suffix. 

\ea\label{ex2:cuvo12}
\langinfo{Cuvok verb \textit{mə̀ʦ} `die' with itive suffix}{}{\cite[396]{dadakGrammaireCuvokLangue2021}} \\
\gll màdùwáŋ	ɛ́-tɛ́-d-ɛ́k,	á-gwàɗ-á	kàd	də̀m	ʃɛ́j	 mə̀ʦ-áɗ-áj \\
rat	3\gl{sg}.\gl{sbj}-\gl{fut}-go/come-\gl{vent}	3\gl{sg}.\gl{sbj}-speak-3\gl{sg}.\gl{obj}	\gl{loc}	girl	\gl{dem}	 2\gl{sg}.\gl{imp}.die-\gl{itive}-\gl{intr} \\
\glt `The rat left and said to the girl: die!'
\z 

\subparagraph{Caused accompanied motion (CAM)}

The verb roots \textit{ks} (singular object) and \textit{hl} (plural object) are glossed `take' (French: \textit{prendre}), but with the ventive suffix the predicate means `bring', as seen in (\ref{ex2:cuvo14}). 

\ea\label{ex2:cuvo14}
\langinfo{Cuvok verb \textit{kə̀s} `take' with ventive suffix}{}{\cite[303]{dadakGrammaireCuvokLangue2021}} \\
\gll  kə̀s-ɛ́k	\\
take-\gl{vent}	\\
\glt `Bring it here.'
\z 

\subparagraph{Buy-sell verbs}

The verb \textit{husàm} can mean both `buy' and `sell'. The use of the `causative' suffix \textit{-dá} with this verb enforces the meaning `sell' [French: `\textit{faire acheter} (\textit{vendre})'] \citep[315]{dadakGrammaireCuvokLangue2021}. One example of this verb with a ventive was found, in which case it is translated `buy', as shown in (\ref{ex2:cuvo13}). 

\ea\label{ex2:cuvo13}
\langinfo{Cuvok verb \textit{husàm} `buy/sell' with ventive suffix}{}{\cite[191]{dadakGrammaireCuvokLangue2021}} \\
\gll dà-zàmàj sə̀, Kàdámà kɛ́-hùʃɛ̀m-ɛ́k ɬàw tá-tə̀màk \\ 
\gl{loc}-Zamay \gl{top} Kadama \gl{pst}.3\gl{sg}.\gl{sbj}-buy/sell--\gl{vent} meat \gl{rel}-sheep	\\
\glt `At Zamay, Kadama bought mutton.'
\z 

\paragraph*{Mafa (mafa1239)} \label{mafa}

Information on the grammar of Mafa is from a grammar sketch and lexicon \citep{barreteauLexiqueMafaLangue1990}. 

\subparagraph{Directional verbs}~

\begin{table}[H]
\caption{Mafa directional verbs} 
\label{tab:mafaV}
\begin{tabular}{lll} % add l for every additional column or remove as necessary
	\lsptoprule
	DIR & verb & gloss   \\ %table header
	\midrule
	ITIVE & \textit{d-} & `\textit{aller}, \textit{partir} [go, leave]' \\
	VENT & \textit{shik-} & `\textit{venir} [come]' \\
	UP & \textit{tə́v-} & `\textit{monter} [go up]' \\
	DOWN & \textit{súlóm-} & `\textit{descendre} [go down]' \\
	\lspbottomrule
\end{tabular}
\end{table}

The adverbial \textit{ápa} `inside there (French:\textit{là-dedans})' is used for inward motion and \textit{âwuda} `outside (French: \textit{dehors})' for outward motion \citep[78, 81]{barreteauLexiqueMafaLangue1990}. There is also a verb \textit{mutuhw-kâda} `\textit{se sortir de} [exit from]' which appears to obligatorily include the ventive suffix.

\subparagraph{Directional extensions}~


\begin{table}[H]
\caption{Mafa directional morphology}
\label{tab:mafa}
\begin{tabular}{llll} % add l for every additional column or remove as necessary
	\lsptoprule
	forms & source gloss & direction  & other functions  \\ %table header
	\midrule
	\textit{-ká(ɗá)}  &  \textit{rappr.} &  \textsc{vent} &  \textsc{subs.vent}  \\
	\lspbottomrule
\end{tabular}
\end{table} 

\citet[45]{barreteauLexiqueMafaLangue1990} call the ventive in Mafa ``\textit{le directionnel de rapprochement}''.\footnote{There is also a locative suffix \textit{-da} which indicates that the activity predicated by the verb took place elsewhere than the place of speech (``\textit{l'action s'est déroulée} (\textit{ou se déroule habituellement}) \textit{dans un lieu où le sujet ne se trouve pas au moment de l'énonciation}'') \citep[45]{barreteauLexiqueMafaLangue1990}. The locative suffix is not said to have any directional functions.} With intransitive motion verbs, the suffix has a ventive directional meaning, as in (\ref{ex2:mafa1}) and (\ref{ex2:mafa2}).

\ea\label{ex2:mafa1}
\langinfo{Mafa verb \textit{mbálə} `chase' with ventive suffix}{}{\cite[46]{barreteauLexiqueMafaLangue1990}} \\
\gll á mbálə-ká kəda	 \\
3\gl{sg}.\gl{m}.\gl{ipfv}  chase-\gl{vent} dog \\
\glt `He runs after the dog towards us.' \newline [French: `\textit{il court après le chien vers nous}']

\ex\label{ex2:mafa2}
\langinfo{Mafa verb \textit{ha} `run' with ventive suffix}{}{\cite[188]{barreteauLexiqueMafaLangue1990}} \\
\gll n ha-káɗíy á ngwáy	\\
3\gl{sg} run-\gl{vent} \gl{prep} house \\ 
\glt `He fled to the house (coming towards us).' \newline [French: `\textit{Il s'est enfui à la maison} (\textit{en venant vers nous}).']
\z 

With most non-motion verbs, the translations given of ventive forms do not indicate any motion-related meaning, except two examples which give a ventive subsequent motion interpretation, shown in (\ref{ex2:mafa3}) and (\ref{ex2:mafa4}). 

\ea\label{ex2:mafa3}
\langinfo{Mafa verb \textit{ta} `cook' with ventive suffix}{}{\cite[46]{barreteauLexiqueMafaLangue1990}} \\
\gll  n ta-ká mávə́r tə gáy	\\
3\gl{sg}.\gl{pfv} cook-\gl{vent} boule in hut \\
\glt `She cooked the boule in the hut and brought it.  \newline
[French: \textit{Elle a cuit la boule de mil dans la case et l'a apportée}.']

\ex\label{ex2:mafa4}
\langinfo{Mafa verb \textit{pár} `gather' with ventive suffix}{}{\cite[188]{barreteauLexiqueMafaLangue1990}} \\
\gll  m párə-káɗá	\\
3\gl{sg}.\gl{pfv} gather-\gl{vent}\\
\glt `He gathered and he came.' [French: `\textit{Il l'a cueilli et est venu}.']
\z 

\subparagraph{Caused accompanied motion (CAM)}

There are several CAM verbs in Mafa including \textit{rudungw-} `\textit{porter une grande quantité, dans un récipient} [carry a large quantity, in a container]' and \textit{ndzâp-} `\textit{porter plusieurs bottes, plusieurs fagots} [carry several bundles (of twigs)]' \citep[291, 314]{barreteauLexiqueMafaLangue1990}. These are not shown with a ventive suffix. The ventive suffix does occur with the verb stem \textit{tsu-} `\textit{prendre} [take]' resulting in the ventive CAM predicate \textit{tsu-kwáɗá} `\textit{apporter, amener} [bring]' \citep[359]{barreteauLexiqueMafaLangue1990}.

\largerpage[-2]
In addition, the causative suffix can be combined with motion verbs to express CAM. This can include manner-of-motion verbs; for example, the causative of \textit{ha-} `\textit{courir} [run]' is \textit{had-} `\textit{porter en courant, conduire vite, faire courir} [carry running, run fast, make run]' and the causative of \textit{zukw-} `\textit{faire semblant de} [pretend]' is \textit{zukud-} `\textit{faire semblant de porter} [pretend to carry]'  \citep[44, 165, 390]{barreteauLexiqueMafaLangue1990}. A similar pattern applies to the verb stem \textit{d-} `\textit{aller} [go]' except that in this case no causative morphology is used. The transitive use of this basic motion verb has the meaning `\textit{porter, amener, conduire, transporter} [carry, bring, drive, transport]' \citep[106]{barreteauLexiqueMafaLangue1990}.

\subparagraph{Buy-sell verbs}

There appear to be two separate verbs for buying and selling in Mafa: \textit{súm-} `\textit{acheter} [buy]' and \textit{pə̀r-} `\textit{délier, vendre} [untie, sell]' \citep[308, 323]{barreteauLexiqueMafaLangue1990}.

\subsection{Sakun (Sukur) (1/1)} \label{saku}

\paragraph*{Sakun (suku1272)}

Sakun (exonym: Sukur) is a South Biu-Mandara language that does not show evidence of being particularly closely related to any other of the other South Biu-Mandara languages \citep[50]{gravinaPhonologyProtoCentralChadic2014}. The language is described in the dissertation of \citet{thomasGrammarSakunSukur2014} and was also the subject of a language documentation project \citep{thomasSakunSukurLanguage2014}. Note that the two hours of annotated video in the documentation project do not appear to have morpheme-by-morpheme glossing, making it impractical to use that corpus to study the distribution of directional morphology. 

\subparagraph{Directional verbs}~

\begin{table}[H]
\caption{Sakun directional verbs} 
\label{tab:sakuV}
\begin{tabular}{lll} % add l for every additional column or remove as necessary
	\lsptoprule
	DIR & verb & gloss   \\ %table header
	\midrule
	ITIVE & \textit{dza} (\textsc{pfv}: \textit{ra}) & `go' \\
	VENT & \textit{já} & `come' \\
	\lspbottomrule
\end{tabular}
\end{table}

Sakun does not appear to have verb roots expressing motion in upward{\slash}downward or inward/outward directions. These directions are rather expressed through verbal extensions.

\newpage
\subparagraph{Directional extensions}~

\begin{table}[H]
\caption{Sakun directional morphology}
\label{tab:saku}
\begin{tabular}{llll} % add l for every additional column or remove as necessary
	\lsptoprule
	forms & source gloss & direction  & other functions  \\ %table header
	\midrule
%	\textit{-kə́}  &  \textsc{vent}  &   ??  &   ??  \\
%	\textit{-rá}  &  \textsc{cent}  & \textsc{vent/up}  &   ?? \\
	\textit{-má}  &  \textsc{up}  &    \textsc{up} &   \\
	\textit{-xá}  &  \textsc{down}  &   \textsc{down}  &    \\
	\textit{-ɣə}  &  \textsc{into}  &   \textsc{into}   &     \\
	\textit{-va}  &  \textsc{out}  &   \textsc{out}   &   lexicalized  \\
	\lspbottomrule
\end{tabular}
\end{table} 

\citet[120]{thomasGrammarSakunSukur2014} lists 17 verbal extensions. Of those, the four suffixes that express directional meaning are shown in Table \ref{tab:saku}. The most straightforward examples of directional verbal morphology in Sakun are the upward, downward and outward suffixes. These combine with motion verbs (intransitive and transitive) to express directional meaning, as exemplified by their use with the verb \textit{ja} `come' in (\ref{ex2:saku3}), (\ref{ex2:saku4}) and (\ref{ex2:saku5}).

\ea\label{ex2:saku3}
\langinfo{Sakun verb \textit{ja} `come' with upward suffix}{}{\cite[286]{thomasGrammarSakunSukur2014}} \\
\gll  ja-má xa sakún ŋí, dza ŋí dza da xaʔí \\
come-\gl{up} \gl{down} Sakun 1\gl{excl} go 1\gl{excl} go to here \\
\glt `We came up from Sakun, before going up here [Damai].'

\ex\label{ex2:saku4}
\langinfo{Sakun verb \textit{ja} `come' with downward suffix}{}{\cite[26]{thomasGrammarSakunSukur2014}} \\
\gll ka yá-xa	\\
\gl{purp} come-\gl{down} \\
\glt `In order to come down…'

\ex\label{ex2:saku5}
\langinfo{Sakun verb \textit{ja} `come' with outward suffix}{}{\cite[126]{thomasGrammarSakunSukur2014}} \\
\gll má ja-vá-j ja-vá kwá ná … \\
when come-\gl{out}-\gl{prog} come-\gl{out} 2\gl{sg} \gl{top} \\
\glt `Immediately when you come out…'
\z 

The inward suffix is only attested twice in \citet{thomasGrammarSakunSukur2014}, both times with the verb \textit{dza} `go' (or its perfective form \textit{ra}), as in (\ref{ex2:saku6}).

\ea\label{ex2:saku6}
\langinfo{Sakun verb \textit{ra} `go.\textsc{pfv}' with inward suffix}{}{\cite[214]{thomasGrammarSakunSukur2014}} \\
\gll  a sə́ kwá ɓá tə́ kə́ ⁿda=j rá-ɣə=ĵ \\
\gl{pfv} know 2\gl{sg} name 3\gl{m}.\gl{poss} \gl{prep} person=\gl{rel} go.\gl{pfv}-\gl{into}=\gl{q} \\
\glt `Do you know the name of the person who entered?'
\z 

With non-motion verbs, the upward, downward and outward suffixes do not have any directional or motion-related interpretation, as in (\ref{ex2:saku7}) and (\ref{ex2:saku8}). It is not clear what the functions of the suffixes may be in this context. 

\ea\label{ex2:saku7}
\langinfo{Sakun verb \textit{bɮk} `burn' with upward suffix}{}{\cite[27]{thomasGrammarSakunSukur2014}} \\
\gll a ⁿda bɮk-má n ku	\\
\gl{seq} person burn-\gl{up} with fire		\\
\glt `One partially burns it with fire.'

\ex\label{ex2:saku8}
\langinfo{Sakun verb \textit{ji} `give birth' with downward suffix}{}{\cite[129]{thomasGrammarSakunSukur2014}} \\
\gll  a	ji-xá	ka	rwi \\
\gl{pfv}	bear-\gl{down}	3\gl{f}	child \\
\glt `She bore a child.'
\z 

\citet[119]{thomasGrammarSakunSukur2014} gives one example of a lexicalized use of a directional suffix. The verb root \textit{tʃíɗ} ‘leak’ with the downward suffix becomes \textit{tʃíɗ-xá} ‘filter’ and with the upward suffix becomes \textit{tʃíɗ-má} ‘coil’. In addition, \citet[121]{thomasGrammarSakunSukur2014} mentions that ``Verbal extensions tend to exhibit a correlation with the telicity of events, or the viewing of the event as bounded or completed'' though this is said to be an ``indirect'' effect. The general pattern in the data presented in \citet{thomasGrammarSakunSukur2014} suggests a high degree of lexicalized usage of this morphology.

The motion verbs \textit{ja} `come' and \textit{dza} `go' are attested with the suffix \textit{-rá} which Thomas calls ``centrifugal'' and glosses \textsc{cent} as in (\ref{ex2:saku9}), (\ref{ex2:saku10}) and (\ref{ex2:saku11}). Examples (\ref{ex2:saku10}) and (\ref{ex2:saku11}) make clear that the suffix does not express motion away from the deictic center as this would conflict with the meaning of the verb as well as the adverb \textit{xaʔí} `here'. It remains unclear what the semantic contribution of this suffix is, but one possibility is that like the itive marker in Buwal, it has a meaning of `away from a location other than the deictic center', as indicated by the adverb `away' in the translation of (\ref{ex2:saku11}).

\ea\label{ex2:saku9}
\langinfo{Sakun verb \textit{dza} `go' with \textsc{cent} suffix}{}{\cite[193]{thomasGrammarSakunSukur2014}} \\
\gll má	dza-ra-j	dza-ra	ka	ná,	a	dza	ka	da=j	ná, \\
\textsc{hypothetical}	go-\textsc{cent}-\gl{prog}	go-\textsc{cent}	3\gl{f}	\gl{top}	\gl{pfv}	go	3\gl{f}	thing=\gl{det}	\gl{top} \\
\glt `Immediately when it went home, it went and (did) this thing…'

\ex\label{ex2:saku10}
\langinfo{Sakun verb \textit{ja} `come' with \textsc{cent} suffix}{}{\cite[203]{thomasGrammarSakunSukur2014}} \\
\gll já-rá	xaʔí \\
come-\textsc{cent}	here \\
\glt `Come here!'

\ex\label{ex2:saku11}
\langinfo{Sakun verb \textit{ja} `come' with \textsc{cent} suffix}{}{\cite[104]{thomasGrammarSakunSukur2014}} \\
\gll  náx	a	ja-rá	və́	xátʃju	 \\
when	\gl{pfv}	come-\textsc{cent}	\gl{asoc}	Hacu \\
\glt `When Hacu’s people came away…'
\z 

With non-motion verbs, the function of the ``centrifugal'' suffix is obscure, as in (\ref{ex2:saku12}) and (\ref{ex2:saku13}). 

\ea\label{ex2:saku12}
\langinfo{Sakun verb \textit{də́} `cook' with \textsc{cent} suffix}{}{\cite[97]{thomasGrammarSakunSukur2014}} \\
\gll a	də́-rá=n	zúŋ	zəgəj	\\
\gl{pfv}	cook-\textsc{cent}=1\gl{sg}	one	time \\		
\glt `I cooked it one time.'

\ex\label{ex2:saku13}
\langinfo{Sakun verb \textit{zázá} `weave' with \textsc{cent} suffix}{}{\cite[372]{thomasGrammarSakunSukur2014}} \\
\gll  amá zázá=j atʃíji di kə́=j dzə-rá ⁿda tʃíji tsú	\\ 
but better=\gl{rel} \gl{dist}.\gl{dem} big \gl{prep}=\gl{rel} weave-\textsc{cent} person \gl{prox}.\gl{dem} also		\\
\glt `But, that one is bigger than the one that they are weaving here again.'
\z 

The suffix \textit{-kə́} is labeled as a ventive directional, but there are no examples given of this suffix with a ventive directional meaning with the possible exception of the ventive form of the verb \textit{ɬa} `take' to express caused accompanied motion (see below). With motion verbs, as in (\ref{ex2:saku1}), and non-motion verbs, as in (\ref{ex2:saku2}), there is no indication that this suffix has a directional or motion-related meaning.\footnote{Note that the glosses \textsc{vent} and \textsc{cent} are erroneously swapped in several places in \citet{thomasGrammarSakunSukur2014}. The glosses are corrected in the data presented here.} 

\ea\label{ex2:saku1}
\langinfo{Sakun verb \textit{təŋ} `move, find' with ``ventive'' suffix}{}{\cite[129]{thomasGrammarSakunSukur2014}} \\
\gll bəzaf ná, ɮa mə təŋtəŋ-kə́ kwá \\
useless \gl{top} only \gl{hab} move.\gl{pl}-\gl{vent} 2\gl{sg}		 \\
\glt `It was a problem, you are just wandering around aimlessly.'

\ex\label{ex2:saku2}
\langinfo{Sakun verb \textit{təŋ} `move, find' with ``ventive'' suffix}{}{\cite[317]{thomasGrammarSakunSukur2014}} \\
\gll a	ka	tʃír-kə́	səɓə́k	…	 \\
\gl{seq}	3\gl{f}	sprout-\gl{vent}	grass	\gl{seq}	 \\
\glt `…it (the groundnuts) grows some weeds.'
\z 

In addition to these, there are a few suffixes worth mentioning in passing. “The directional extension \textit{-ᵐta} `to the bush' only combines with the root \textit{dza} ‘to go’ and its perfective form \textit{ra}” \citep[122]{thomasGrammarSakunSukur2014}. Since this form only appears with one verb root, it can be considered a lexical form. The suffix \textit{-vá} \textsc{across} is included in Thomas’ list but not discussed or elsewhere attested in the dissertation. Note that it only differs from the outward suffix in tone marking. The suffix \textit{-tʃikə́} \textsc{spread} is also included in the list but not discussed either and only attested in one example where it is glossed \textsc{cover} \citep[232]{thomasGrammarSakunSukur2014}. Thomas also lists a suffix \textit{-j} \textsc{dir} noting “\textsc{dir} is an abbreviation for ‘direction’. This extension only occurs with \textit{dza} ‘go’ and is used when the direction of motion is unspecified” \citep[120]{thomasGrammarSakunSukur2014}. The verbal suffix \textit{-j} is only glossed as \textsc{dir} twice in the dissertation, and is normally labeled \textsc{prog}. I assume this is a tense-aspect marker and that the direction-related analysis was an early hypothesis that did not get fully corrected in the process of writing the final dissertation. The suffix \textit{-ʃì} is glossed \textsc{to} or \textsc{follow}. Its meaning is not discussed and is not apparent from the examples available.

Etymologically, some of the directional extensions are transparently derived from prepositions. The upward and downward suffixes share their form with prepositions of similar meaning, and the inward suffix \textit{-ɣə} \textsc{into} is similar to the preposition \textit{ɣi} ‘in’ \citep[273]{thomasGrammarSakunSukur2014}. 

\subparagraph{Caused accompanied motion (CAM)}

Since the glossing of verbs in \citet{thomasGrammarSakunSukur2014} is inconsistent, it is not clear whether there are any verbs that inherently express caused accompanied motion. The verb \textit{ɬa} is normally glossed `take'. The verb can have a directed caused accompanied motion (CAM) interpretation when combined with the ventive suffix, as in (\ref{ex2:saku16}), the ``centrifugal'' suffix, as in (\ref{ex2:saku17}), or the outward suffix, as in (\ref{ex2:saku18}).

\ea\label{ex2:saku16}
\langinfo{Sakun verb \textit{ɬá} `take' with ventive suffix}{}{\cite[259]{thomasGrammarSakunSukur2014}} \\
\gll  a ɬá-kə́ pəʃím gwamənátí xáⁿdʒiga kíŋ \\
\gl{pfv} carry-\gl{vent} clever government now many \\
\glt `The government has brought many plans (for farming).'

\ex\label{ex2:saku17}
\langinfo{Sakun verb \textit{ɬá} `take' with ``centrifugal'' suffix}{}{\cite[153]{thomasGrammarSakunSukur2014}} \\
\gll a	ɬá-rá	tʃá	ɬíɗi	a	fa	tə \\
\gl{pfv}	take-\textsc{cent}	3\gl{m}	chief	to	father	3\gl{m}.\gl{poss} \\
\glt `He took the kingship from his father.'

\ex\label{ex2:saku18}
\langinfo{Sakun verb \textit{ɬá} `take' with ``centrifugal'' suffix}{}{\cite[379]{thomasGrammarSakunSukur2014}} \\
\gll má {maʃín warakə́} ná, a kwá ɬá-tá-va vərʃín xaʔí \\
\textsc{hypothetical} early.morning \gl{top} \gl{sbjv} 2\gl{sg} take-3\gl{pl}.\gl{obj}-\gl{out} child.\gl{pl} here		\\
\glt `Early in the morning, you will take (it) out to the children.'
\z 

Directed CAM is also seen expressed by the outward suffix with the verb \textit{pəká} glossed `pick', `collect', `gather', `take', `put' and `bring', as in (\ref{ex2:saku19}), or with the ventive suffix and the verb \textit{ɓáts} which is glossed `gather',  `get', `take', `harvest' and `collect', as in (\ref{ex2:saku20}).

\ea\label{ex2:saku19}
\langinfo{Sakun verb \textit{pəká} `pick' with outward suffix}{}{\cite[87]{thomasGrammarSakunSukur2014}} \\
\gll da ɗwá-tə xá=j pəká-va ⁿda \\
thing eat-\gl{obj} \gl{pl}=\gl{rel} bring-\gl{out} person \\
\glt `the things to eat that they brought outside'

\ex\label{ex2:saku20}
\langinfo{Sakun verb \textit{ɓáts} `get' with ventive suffix}{}{\cite[263]{thomasGrammarSakunSukur2014}} \\
\gll a ka ɓáts-kə́ daf=j tsú káⁿgə də́ma=ju \\
\gl{sbjv} 3\gl{f} gather-\gl{vent} food=\gl{det} also \gl{ben} girl=\gl{det} \\
\glt `She will bring in that food also, for that girl.'
\z 

In addition, one example appears to show a multiverb construction in which the verb \textit{ɬa} `take' and a directional verb combine to express directed CAM. In (\ref{ex2:saku21}), the upward suffix on the directional verb would seem to indicate that the direction is both towards the deictic center (ventive) and upward, although the translation does not make this explicit. 

\ea\label{ex2:saku21}
\langinfo{Sakun multiverb construction}{}{\cite[206]{thomasGrammarSakunSukur2014}} \\
\gll  íⁿda mara ɬa sə́li ja-má xaʔí áj \\
\gl{neg} Mara carry shame come-\gl{up} here \gl{excl} \\
\glt `Don’t let Mara bring shame here!'
\z 

\subparagraph{Buy-sell verbs}

There are separate verbs for `buy' and `sell' in Sakun. The verb \textit{ɮamá} `buy' sometimes occurs with a ``centrifugal'' suffix, and in one case the \textit{-ma} morpheme is glossed as an upward suffix \citep[210]{thomasGrammarSakunSukur2014}. The verb \textit{pəra} is glossed `sell', and another verb \textit{dlə́má} is also glossed `sell'. 

\subsection{Teraic (2/5)}

Two Teraic languages are included here: Ga'anda and Tera. No relevant publications were found for Boga (boga1251), Hwana (hwan1240) or Jara (jara1274).

\paragraph*{Ga'anda (gaan1243)} \label{gaan}

Information on the grammar of Ga'anda is from a doctoral dissertation written within the Case Grammar framework \citep{newmanCaseGrammarGa1971}. Example sentence generally do not include interlinear glossing. Glosses are added below where possible.

\subparagraph{Directional verbs}~

\begin{table}[H]
	\caption{Ga'anda directional verbs} 
	\label{tab:gaanV}
	\begin{tabular}{lll} % add l for every additional column or remove as necessary
		\lsptoprule
		DIR & verb & gloss   \\ %table header
		\midrule
		ITIVE & \textit{ɗə} & `go' \\
		VENT & \textit{ɓa} & `come' \\
		INTO& \textit{yim} & `enter' \\
		\lspbottomrule
	\end{tabular}
\end{table}

\subparagraph{Directional extensions}~

\begin{table}[H]
	\caption{Ga'anda directional morphology}
	\label{tab:gaan}
	\begin{tabular}{llll} % add l for every additional column or remove as necessary
		\lsptoprule
		forms & source gloss & direction  & other functions  \\ %table header
		\midrule
		\textit{kaɗe}  &  particle  &    \textsc{itive/out/back} &  lexicalized \\
		\textit{xa}  &  particle  & \textsc{down}   &  \textsc{compl}  \\
%		\textit{in}  &  particle  &   NA &  \textit{subs.vent} \\
		\lspbottomrule
	\end{tabular}
\end{table} 

The ``particle'' \textit{kaɗə} is described has having two meanings: ``The first denotes that the action is done `away, out, back'... The second sense... indicates completeness of the action'' \citep[184--185]{newmanCaseGrammarGa1971}. The particle is not shown with typical translocative motion verbs, but has a directional function in a few examples such as in (\ref{ex2:gaan2}) and (\ref{ex2:gaan3}). 

\ea\label{ex2:gaan2}
\langinfo{Ga'anda verb \textit{ɬerta} `spill' with itive particle}{}{\cite[184]{newmanCaseGrammarGa1971}} \\
\gll na 'yem ɬerta kaɗə	\\
\gl{fut} water spill \gl{itive} \\
\glt `Water will spill out.'

\ex\label{ex2:gaan3}
\langinfo{Ga'anda verb \textit{məkani} `throw' with itive particle}{}{\cite[185]{newmanCaseGrammarGa1971}} \\
\gll  ə məkani xwermɗa kaɗə		\\
?? throw corn \gl{itive} \\
\glt `I threw out/away the corn.'
\z 

In at least one example, shown in (\ref{ex2:gaan1}), particle \textit{kaɗə} occurs with the ventive directional verb, tranlsated as `come back'.

\ea\label{ex2:gaan1}
\langinfo{Ga'anda verb \textit{ɓa} `come' with itive particle}{}{\cite[192]{newmanCaseGrammarGa1971}} \\
\gll  ə ɓa-incə kaɗə	\\
I come-?? \gl{itive}	\\
\glt `I came back yesterday.'
\z 

The particle \textit{xa} is also described as having two meanings: ``The first is associated with a downward direction of the action... The second meaning denotes that the action is both well done and completely done...'' \citep[186]{newmanCaseGrammarGa1971}. The only example found of \textit{xa} with a downward directional meaning is the transitive motion verb in (\ref{ex2:gaan4}). No motion-related translations are found in examples where \textit{xa} is used with non-motion verbs. 

\ea\label{ex2:gaan4}
\langinfo{Ga'anda verb \textit{ɗəfu} `put' with downward particle}{}{\cite[195]{newmanCaseGrammarGa1971}} \\
\gll  ɗəfu ɗelwertəɗi xa nec kə yina	\\
put book \gl{down} \gl{rel} ?? ?? 		\\
\glt `Put down the book which is mine.'
\z 

In addition, it is worth mentioning the particle \textit{in} which is not shown to have a directional function, but in some contexts may have a subsequent motion interpretation. \citet{newmanCaseGrammarGa1971} describes the meaning of \textit{in} as ``a non-directional `along' in transitive constructions...'' This can often be interpreted as meaning `in addition to' but it is not always clear how to interpret the English translations given. With some transitive non-motion verbs, the particle \textit{in} appears to have a benefactive or subsequent caused accompanied motion translation, as in (\ref{ex2:gaan5}) and (\ref{ex2:gaan6}). 

\ea\label{ex2:gaan5}
\langinfo{Ga'anda verb \textit{sər} `fry' with `along' particle}{}{\cite[184]{newmanCaseGrammarGa1971}} \\
\gll sər-áǹ-tǹ ɬiw in	\\
fry-??-?? meat along \\
\glt `Fry meat (and bring) to him!'

\ex\label{ex2:gaan6}
\langinfo{Ga'anda verb \textit{yark} `steal' with `along' particle}{}{\cite[184]{newmanCaseGrammarGa1971}} \\
\gll  ə yark-ìcə̀ in	\\
?? steal-?? along		\\
\glt `He stole (and brought) for me.'
\z

\subparagraph{Caused accompanied motion (CAM)}

\citet[26]{newmanCaseGrammarGa1971} writes that ``intransitive verbs like `come, go'... may take direct objects preceded by `with' to mean `bring, take'.'' However, in the one example found of a directional motion verb expressing CAM, there does not appear preposition before the object, as seen in (\ref{ex2:gaan7}).

\ea\label{ex2:gaan7}
\langinfo{Ga'anda verb \textit{ɓa} `come' with direct object}{}{\cite[189]{newmanCaseGrammarGa1971}} \\
\gll nanɗa ɓa-i-ta toxwatə xar in \\
they come-??-?? soup some along \\
\glt `They will bring me some soup.'
\z 

\subparagraph{Buy-sell verbs}

The verb \textit{xiyə} `buy' when followed by the `away' particle becomes \textit{xiyə kaɗə} `sell' \citep[44]{newmanCaseGrammarGa1971}. 

\paragraph*{Tera (tera1251)} \label{tera}

Information on Tera is from a Transformational Syntax analysis of the grammar \citep{newmanGrammarTeraTransformational1970}. This publication does not include interlinear glosses, so these are added below where possible with the help of a Tera lexicon \citep{muazuModernTeraEnglish2015}.

\subparagraph{Directional verbs}~

\begin{table}[H]
	\caption{Tera directional verbs} 
	\label{tab:teraV}
	\begin{tabular}{lll} % add l for every additional column or remove as necessary
		\lsptoprule
		DIR & verb & gloss   \\ %table header
		\midrule
		ITIVE & \textit{ɗí} & `go' \\
		VENT & \textit{ɓa} & `come' \\
		INTO/OUT & \textit{vi} & `enter, go out' \\
		\lspbottomrule
	\end{tabular}
\end{table}

%The expression \textit{vi nə́ dàm} glossed `enter to out' is translated as `go out' \citep{newmanGrammarTeraTransformational1970}. For evidence that \textit{vi} on its own can have outward motion meaning, the English lexicon of \citet[128]{muazuModernTeraEnglish2015} gloss `exit' as \textit{vi} without the adverbial phrase following. 

\subparagraph{Directional extensions}~

%\begin{table}[H]
%	\caption{Tera directional morphology}
%	\label{tab:tera}
%	\begin{tabular}{llll} % add l for every additional column or remove as necessary
%		\lsptoprule
%		forms & source gloss & direction  & other functions  \\ %table header
%		\midrule
%		\textit{ɓara}  &  particle  &  \textsc{dir.itive} (transitive) &   ??  \\
%		\lspbottomrule
%	\end{tabular}
%\end{table} 

According to \citet[18]{newmanGrammarTeraTransformational1970}, the ``particle'' \textit{ɓara} ``indicates action away from the speaker (like German \textit{ver-}),'' as in (\ref{ex2:tera1}) and (\ref{ex2:tera2}). However, \citet[314--315]{schuhChadicCornucopia2017} argues that this form is better considered an adverbial rather than a verbal extension, based on its syntactic distribution. 

%Note that these examples are of transitive motion verbs. No examples are given of  \textit{ɓara} with an intransitive motion verb.

\ea\label{ex2:tera1}
\langinfo{Tera verb \textit{mbukə} `throw' with itive particle}{}{\cite[19]{newmanGrammarTeraTransformational1970}} \\
\gll mbukə ɓara	\\
throw away		\\
\glt `to throw away'

\ex\label{ex2:tera2}
\langinfo{Tera verb \textit{vər} `give' with itive particle}{}{\cite[44]{newmanGrammarTeraTransformational1970}} \\
\gll Ali wà vər ye-rem ɓara	\\
Ali ?? give to-us away		\\
\glt `Ali gave it away to us.'
\z 

While \textit{ɓara} has the distribution of an adverb in the syntax, it is still used in idiomatic expressions, such as in the case of `buy' and `sell' (discussed below) and with apparent aspectual meaning, as in (\ref{ex2:tera3}). 

\ea\label{ex2:tera3}
\langinfo{Tera verb \textit{nji} `eat' with itive particle}{}{\cite[86]{newmanGrammarTeraTransformational1970}} \\
\gll Ali-E nji ɮu-a ɓara \\
Ali-?? ate meat-the away \\
\glt `Ali ate up the meat.'
\z 

\subparagraph{Caused accompanied motion (CAM)}

The verb \textit{ɓasi} `bring' expresses caused accompanied motion. It is not shown with any directional particles. It is possible for ventive CAM to be expressed by the ventive motion verb with a prepositional phrase, as in (\ref{ex2:tera4}). 

\ea\label{ex2:tera4}
\langinfo{Tera verb \textit{ɓa} `come' with comitative phrase}{}{\cite[113]{newmanGrammarTeraTransformational1970}} \\
\gll Ali n-a ɓa ndə woy-a fan ɓa \\
Ali ??-?? come with boy-\gl{def} ?? \gl{neg} \\
\glt `Ali didn't bring (come with) the boy here.'
\z 

\subparagraph{Buy-sell verbs}

The verb \textit{masa} `to buy' combined with the itive particle \textit{masa ɓara} means `to sell' \citep[19]{newmanGrammarTeraTransformational1970}.
